\documentclass[11pt,reqno]{amsart}

\usepackage[a4paper,margin=28mm]{geometry}
\usepackage{amsmath,amssymb,mathtools,mathrsfs}
\usepackage{enumitem,booktabs,longtable,array}
\usepackage{microtype}
\usepackage{xurl}
\usepackage[colorlinks=true,linkcolor=blue,citecolor=blue,urlcolor=blue]{hyperref}

\numberwithin{equation}{section}
\setlist[enumerate]{label=\textup{(\roman*)},leftmargin=2.2em}

\newtheorem{theorem}{Theorem}[section]
\newtheorem{proposition}[theorem]{Proposition}
\newtheorem{lemma}[theorem]{Lemma}
\newtheorem{corollary}[theorem]{Corollary}
\newtheorem{example}[theorem]{Example}
\newtheorem{remark}[theorem]{Remark}
\theoremstyle{definition}
\newtheorem{definition}[theorem]{Definition}

\newcommand{\N}{\mathbb N}
\newcommand{\Nzero}{\mathbb N_0}
\newcommand{\C}{\mathcal C}
\newcommand{\D}{\mathcal D}
\newcommand{\Acal}{\mathcal A}
\newcommand{\Lcal}{\mathcal L}
\newcommand{\Patt}{\mathsf P}
\newcommand{\hseq}{h_{\mathrm{seq}}}
\newcommand{\hpat}{h^{*}}
\newcommand{\hinv}{h_{\mathrm{inv}}}
\newcommand{\htop}{h_{\mathrm{top}}}
\newcommand{\Bor}{\mathrm{Bor}}
\newcommand{\clop}{\mathrm{clop}}
\newcommand{\lang}{\mathrm{lang}}
\newcommand{\supp}{\operatorname{supp}}
\newcommand{\doi}[1]{\href{https://doi.org/#1}{doi:\nolinkurl{#1}}}


\newcommand{\K}{\mathcal K}
\newcommand{\Ssafe}{\mathcal S}
\newcommand{\hstar}{h^*}
\newcommand{\one}{\mathbf 1}
\newcommand{\opnorm}[1]{\left\|#1\right\|_{\infty\to\infty}}
\newcommand{\esssup}{\operatorname*{ess\,sup}}
\newcommand{\R}{\mathbb R}
\newcommand{\Leb}{\operatorname{Leb}}
\newcommand{\cl}{\operatorname{cl}}
\newcommand{\intr}{\operatorname{int}}
\newcommand{\sgn}{\operatorname{sgn}}
\newcommand{\shift}{\sigma}
\newcommand{\ind}{\operatorname{ind}}
\newcommand{\const}{\mathrm{const}}
\newcommand{\ellA}{\ell_2^{A}}
\newcommand{\eps}{\varepsilon}
\allowdisplaybreaks[1]
\hypersetup{pdftitle={Sparse Sampling and Inverse Capacity in Safe Feedback Systems},pdfauthor={Ziqing Ding}}
\title[Sparse sampling and inverse capacity]{Sparse Sampling and Inverse Capacity in Safe Feedback Systems}
\author{Ziqing Ding}
\address{School of Mathematical Sciences, Nanjing Normal University,
Nanjing 210023, China}
\date{October 9, 2026}
\subjclass[2020]{93C55, 37B10, 37B40, 94A17}
\keywords{safe feedback, sequence entropy, maximal pattern entropy,
Borel partition, inverse capacity, symbolic dynamics, matrix products}
\begin{document}
\begin{abstract}
We study sparse observations of safe feedback names at a fixed control
period, allowing all legal Borel feedbacks and all finite Borel atom
refinements. Products of safe inverse-transfer operators give a
sufficient condition for every feedback to have positive-density
multivalued independence in its actual names. When this condition holds
at the minimum safety-cover cardinality $r$, the minimum maximal-pattern
rate is exactly $\log r/\tau$. For finite-cylinder safety domains, the
full operator products have exact matrix representations. We derive
conditions on the allowed cylinders that imply uniform contraction,
including a class with gain columns and no common positive nonstrict
left weight. A complementary result treats common successors: for a
finite closed cover of a mixing one-sided shift of finite type, every
prescribed sampling sequence with diverging gaps gives the exact minimum
$\log r$ over all Borel refinements. Countable-language attainment,
symbolic realizations, and affine examples describe the relation between
these two conclusions and invariance entropy. A cyclic-permutation
application and a weighted decomposition of actual even--odd pattern
pairs make the role and limitations of safe inverse coverage explicit. In particular, the
product criterion gives feedback-dependent independent positions; it
does not settle the corresponding common-sampling problem.
\end{abstract}
\maketitle
\tableofcontents

\section{Introduction}\label{sec:intro}

A safe feedback rule selects a control block from the current state and
repeats the same rule at successive block boundaries. Its instruction
names record which blocks are used. We ask how complicated these names
must remain when they are observed at selected times and the feedback
is chosen to minimize their complexity. The control period $\tau$ is
fixed, and a name observed at $n$ block times is normalized by $n\tau$.
The gaps between observations do not enter this normalization.

There are two different optimization problems. Let $h_A(s)$ be the
counting entropy of feedback $s$ along a sampling sequence $A$, and let
$h^*(s)$ be its maximal-pattern rate. We study
\[
 L_\tau=\sup_A\inf_s h_A(s),\qquad
 M_\tau=\inf_s h^*(s),\qquad L_\tau\le M_\tau.
\]
The infima range over all legal Borel state feedbacks. A lower bound for
$M_\tau$ may use coordinates chosen separately for each feedback. A
lower bound for $L_\tau$ needs one sampling sequence that works for the
whole class. Arbitrary Borel choices on overlapping safety domains can
change the future orbit, so the latter requirement is stronger.

Our first result concerns $M_\tau$. For each control block $u$, let
$K_u$ be the set of states from which the complete block is safe, and
let $L_u$ be its inverse-transfer operator relative to a full-support
probability. At each time retain a global list of $k-1$ labels and sum
their operators. Theorem~\ref{thm:capacity-main} shows that uniform
exponential decay of all products of these sums forces $k$-valued
independence on a positive fraction of the positions in every sufficiently
long feedback name set. The fraction is uniform over all Borel feedbacks.
Consequently $h^*(s)\ge\log k/\tau$. If $k=r$, where $r$ is the
minimum number of safety domains covering the state space, an ordered
subcover selector attains the matching upper bound:
$M_\tau=\log r/\tau$.

The proof first estimates the measures of actual initial-state fibers.
This yields an exponential lower bound for the number of label-list
boxes needed to cover the feedback names. A finite combinatorial lemma
of Kerr and Li \cite[Lemma~3.3]{KL2007}, also derived from the stronger
estimates of Huang and Ye \cite[Lemma~4.1 and Remark~4.2]{HuangYe2009},
then supplies independent positions. The control-specific step is the
connection between these boxes and products with all intermediate safety
restrictions retained. No invariant measure or continuity of the Borel
closed-loop map is required.

For branches $P_u\sigma$ on a full shift and finite-cylinder safety
domains, the product norms admit exact finite-dimensional
representations. Sections~\ref{sec:cylinder} and~\ref{sec:gainloss}
derive contraction from the allowed cylinders in three ways: a common
strict left weight, loss after two steps, and compensation of gain
columns by unavoidable loss at the next step. In the last class no
positive common nonstrict left weight exists in the original
first-coordinate representation. These are applications of established
positive-matrix ideas, including the loss mechanisms studied by
Shi et al.\ \cite{Shi2021} and the Lyapunov-template comparisons of
Debauche et al.\ \cite{Debauche2022}. The estimates here are proved
directly for every history of global label lists.

Our second result addresses $L_\tau$ when all legal blocks have the
same successor. Theorem~\ref{cc-thm:cover} states that for any finite
closed cover of a mixing one-sided shift of finite type, every prescribed
sampling sequence with diverging gaps has minimum counting entropy
$\log r$ over all Borel refinements. The proof uses the Baire property
of the atoms and concatenation of sufficiently separated cylinders.
The control consequence is $L_\tau=M_\tau=\log r/\tau$
(Theorem~\ref{cc-thm:control}). Romagnoli \cite{Romagnoli2007}
computed a cover--partition gap for an excluded-letter cover and clopen
partitions at consecutive times. Here the partition may be Borel and
the sparse sequence is fixed in advance. The consecutive-time question
for Borel refinements is not settled by this argument.

Common positive thresholds already occur in the measure-theoretic
theory: Huang, Maass, and Ye \cite[Lemma~3.3 and Theorem~3.5]{HuangMaassYe2004}
construct one sequence with a common positive threshold for refinements
of an appropriate measurable cover. Their result uses an invariant
measure; our closed-cover formula gives an exact counting value.
The broader sequence-entropy and independence background includes
Goodman \cite{Goodman1974}, Hulse \cite{Hulse1982}, Kamae and Zamboni
\cite{KamaeZamboniSystems2002}, and Kerr and Li \cite{KL2007,KL2009}.

The two principal arguments occupy
Sections~\ref{sec:capacity}--\ref{sec:common-successor}.
The remaining sections develop their scope. A feedback-dependent
example has $L_1=M_1<\log r$. A deletion argument gives simultaneous
attainment for countably many finite-alphabet languages. Symbolic safety
realizations include a Thue--Morse system with zero invariance entropy
and positive fixed-period pattern infimum, whereas optimization over
all periods recovers Kawan's invariance entropy \cite{Kawan2013}.
The affine examples and appendices give exact sampled values and
explain the consistency conditions needed to implement an open-loop
coding as a state feedback. The finite-control example in
Appendix~\ref{app:pointwise} distinguishes a common minimum from a
sequence attaining each feedback's own rate. The capacity and
probability examples in Appendix~\ref{app:capacity-boundaries} show
why neither transfer peaks nor Shannon bounds can replace support counts.
Corollary~\ref{cor:cyclic} gives a direct cyclic-permutation application
of the product criterion. Appendix~\ref{app:finiteblock} decomposes the
affine counting estimate using weighted actual pattern pairs, and
Example~\ref{ex:periodic-support} shows why countable exceptional states
cannot be discarded. Lemma~\ref{lem:finite-switch} and
Example~\ref{ex:one-wait} distinguish maximal-pattern stability from
stability at a prescribed sampling sequence.


\section{Control systems and invariant name languages}\label{sec:control}

Write $\N=\{1,2,\ldots\}$ and $\Nzero=\{0,1,2,\ldots\}$.
All logarithms are natural. Invariance entropy and consecutive-name entropy
are measured per physical time unit; sparse rates are defined below.
Let $X$ be a
metrizable space, let $U$ be a nonempty control set, and suppose that
$f_u:=f(\,\cdot\,,u):X\to X$ is continuous for every $u\in U$.  For a finite
control word $\omega=(\omega_0,\ldots,\omega_{N-1})\in U^N$, set
\begin{align*}
 \varphi(0,x,\omega)=x,
 \qquad
 \varphi(n+1,x,\omega)
 =f\bigl(\varphi(n,x,\omega),\omega_n\bigr)
 \quad(0\le n<N).
\end{align*}
A nonempty compact set $Q\subset X$ is \emph{controlled invariant} if every
$x\in Q$ admits an infinite control whose complete forward trajectory stays in
$Q$.

For $N\in\N$, let $r_{\rm inv}(N,Q)$ be the least cardinality of a set
$S\subset U^N$ such that for every $x\in Q$ some $\omega\in S$ satisfies
\begin{align*}
 \varphi(n,x,\omega)\in Q\qquad(0\le n\le N).
\end{align*}
If no finite such set exists, put $r_{\rm inv}(N,Q)=+\infty$.  The
discrete-time invariance entropy is
\begin{equation}\label{eq:hinv-definition}
 \hinv(Q)=\limsup_{N\to\infty}\frac1N\log r_{\rm inv}(N,Q).
\end{equation}
Throughout, the admissible control space is the full sequence space
$\mathcal U=U^{\Nzero}$; equivalently one may use Kawan's two-sided full
control space and retain its nonnegative coordinates.  Every finite word
therefore extends to an admissible infinite control, and finite-horizon
dependence identifies the preceding spanning number with Kawan's.  Thus
\eqref{eq:hinv-definition} is Kawan's definition with the logarithm converted
from base two to base $e$; see \cite[Definition~2.2, p.~44]{Kawan2013}.

Fix a sampling period $\tau\in\N$.  A \emph{finite invariant cover} of $Q$
with period $\tau$ is a triple $\C=(\Acal,\tau,\nu)$, where $\Acal$ is a finite
cover of $Q$ and $\nu(D)\in U^\tau$ is assigned to each $D\in\Acal$, such that
\begin{equation}\label{eq:invariant-cover}
 \varphi(s,x,\nu(D))\in Q
 \qquad(x\in D,\ 0\le s\le\tau).
\end{equation}
This is the discrete-time specialization of Kawan's Definition~2.8
\cite[p.~79]{Kawan2013}.  An \emph{invariant partition} is an invariant cover
whose underlying family is a finite Borel partition of $Q$ into nonempty
atoms.  It is \emph{clopen} if its atoms are clopen in $Q$.  We write
\begin{align*}
 \mathfrak C_{\Bor}(Q,\tau)&=\bigl\{\text{finite Borel invariant partitions of
 }Q\text{ of period }\tau\bigr\},\\
 \mathfrak C_{\clop}(Q,\tau)&=\bigl\{\C\in\mathfrak C_{\Bor}(Q,\tau):
 \C\text{ is clopen}\bigr\},\\
 \mathfrak C_{\Bor}(Q)&=\bigcup_{\tau\in\N}\mathfrak C_{\Bor}(Q,\tau).
\end{align*}
We use the extended-value convention $\inf\varnothing=+\infty$.

Let $\C=(\{D_i:i\in I_\C\},\tau,\nu)$ be an invariant partition, with
$I_\C$ a finite nonempty alphabet.  For
$a=(a_0,\ldots,a_{N-1})\in I_\C^N$, concatenate the assigned blocks in that
order.  The word $a$ is \emph{admissible} if some $x\in D_{a_0}$ follows this
concatenated control and satisfies
\begin{equation}\label{eq:admissible-word}
 \varphi(j\tau,x,\nu(D_{a_0})\cdots\nu(D_{a_{N-1}}))
 \in D_{a_j}\qquad(0\le j<N).
\end{equation}
The set of such words is $W_N(\C;Q)$.  These are the partition words of
Kawan's Definition~2.9 \cite[pp.~79--80]{Kawan2013}.

\begin{lemma}[Safety, suffixes, and right extension]\label{lem:safety}
Let $\C$ be an invariant partition.
\begin{enumerate}
\item A witness for $a\in W_N(\C;Q)$ remains in $Q$ at every physical time
from $0$ through $N\tau$ under the concatenated blocks.
\item Every prefix and every suffix of an admissible word is admissible.
\item Every $a\in W_N(\C;Q)$ has a right extension in
$W_{N+1}(\C;Q)$.  In particular, every finite admissible word has an
infinite right extension.
\item For $M,N\in\N$,
\begin{equation}\label{eq:word-submult}
 \#W_{M+N}(\C;Q)
 \le \#W_M(\C;Q)\,\#W_N(\C;Q).
\end{equation}
\end{enumerate}
\end{lemma}

\begin{proof}
On block $j$, condition \eqref{eq:invariant-cover}, applied to the state in
$D_{a_j}$, keeps the trajectory in $Q$ throughout the next $\tau$ steps.
Induction proves (i).  Prefixes are immediate.  A suffix beginning at block
$j$ is witnessed by the state at time $j\tau$, which proves (ii).

At time $N\tau$ the terminal state belongs to a unique atom $D_b$.
Appending $b$ and its assigned block gives an admissible word of length
$N+1$, proving (iii) without compactness.  Finally, an admissible word of
length $M+N$ is determined injectively by its length-$M$ prefix and
length-$N$ suffix.  Part (ii) places these in $W_M$ and $W_N$, respectively,
and gives \eqref{eq:word-submult}.
\end{proof}

Define the finite-prefix name system
\begin{equation}\label{eq:name-system}
 \Sigma_\C=\bigl\{a\in I_\C^{\Nzero}:a|_{[0,N-1]}\in W_N(\C;Q)
 \text{ for every }N\in\N\bigr\}.
\end{equation}
We give $I_\C$ the discrete topology and $\Sigma_\C$ the subspace topology
from $I_\C^{\Nzero}$.  Its complement is a union of cylinders determined by
forbidden finite prefixes, so $\Sigma_\C$ is closed.
Lemma~\ref{lem:safety}(iii) shows that every word in $W_N$ is the initial
word of some element of $\Sigma_\C$.  Lemma~\ref{lem:safety}(ii) shows
$\sigma(\Sigma_\C)\subset\Sigma_\C$; surjectivity need not hold.  For Borel
atoms, an element of $\Sigma_\C$ need not have one common state witness for
all its prefixes.  Our definitions use only finite projections and do not
assume such a witness.

There is a useful concrete description of $\Sigma_\C$. The feedback
block map and its consecutive atom names also appear in the numerical
framework of \cite[Section~2]{TomarKawanZamani2022}. Define
\begin{align*}
 T_\C(x)=\varphi(\tau,x,\nu(D_i))\quad\text{if }x\in D_i,
\end{align*}
and its atom itinerary $\iota_\C(x)$ by
$(\iota_\C(x))_j=i$ if $T_\C^j(x)\in D_i$.
Each branch of $T_\C$ is continuous, and $T_\C$ is Borel measurable;
it need not be continuous when the atoms are merely Borel. Every finite
admissible word is the prefix of the actual itinerary of its witness.
Conversely every actual itinerary has admissible prefixes. It follows that
\begin{equation}\label{eq:itinerary-closure}
 \Sigma_\C=\overline{\iota_\C(Q)}.
\end{equation}
Indeed, each prefix cylinder meeting $\Sigma_\C$ meets
$\iota_\C(Q)$, which proves density, and all actual itineraries belong
to the closed set $\Sigma_\C$.

\begin{remark}[Infinite witnesses]\label{rem:witness}
If the atoms are clopen, $T_\C$ and $\iota_\C$ are continuous.
Compactness of $Q$ makes $\iota_\C(Q)$ closed, so every element of
$\Sigma_\C$ has a common state witness. Equivalently, the sets of witnesses
for its finite prefixes are nested nonempty compact sets.

This conclusion can fail for a Borel partition even in a compact
continuous system with one control. Take $Q=X=[0,1]$,
$f(x)=\min\{2x,1\}$, $D_0=(0,1/2)$ and $D_1=Q\setminus D_0$.
Both atoms receive the unique control, so the partition is invariant.
For every $N$, the point $x_N=2^{-N-1}$ witnesses the word $0^N$:
for $0\le j<N$, $f^j(x_N)=2^{j-N-1}\in(0,1/2)$.
Thus $0^\infty\in\Sigma_\C$. No point witnesses this infinite name:
zero is outside $D_0$, and every positive point eventually leaves $D_0$.
\end{remark}

For a nonempty finite $F\subset\Nzero$, put
\begin{equation}\label{eq:control-pattern-count}
 \Patt_\C(F)=\#\{a|_F:a\in\Sigma_\C\}.
\end{equation}
Equivalently, if $m=\max F$, this is the number of projections of
$W_{m+1}(\C;Q)$ onto $F$.  Thus a counted pattern always has a finite witness
that remains in $Q$ also between the selected block times.

Let
\begin{align*}
 \mathscr S=\{A=(a_n)_{n\ge1}:0\le a_1<a_2<\cdots, a_n\in\Nzero\}.
\end{align*}
Define
\begin{align}
 p_\C^*(n)&=\max_{\substack{F\subset\Nzero\\\#F=n}}\Patt_\C(F),
 \label{eq:pstar-control}\\
 \hpat(\C;Q)&=\lim_{n\to\infty}\frac1{n\tau}\log p_\C^*(n),
 \label{eq:hpat-control}\\
 \hseq^A(\C;Q)&=\limsup_{n\to\infty}\frac1{n\tau}
 \log\Patt_\C(\{a_1,\ldots,a_n\}).
 \label{eq:hseq-control}
\end{align}
The existence of the limit in \eqref{eq:hpat-control} follows from
Lemma~\ref{lem:pattern-submult} below.  Although there are infinitely many
$n$-element subsets of $\Nzero$, the maximum in \eqref{eq:pstar-control}
exists: its possible values form a nonempty subset of
$\{1,\ldots,(\#I_\C)^n\}$.

The entries of $A$ are \emph{block indices}; the corresponding physical
observation times are $\tau A=(\tau a_n)_{n\ge1}$.  The normalization in
\eqref{eq:hseq-control} is $n\tau$, not the elapsed horizon $a_n\tau$.
For fixed $\tau$, $A\mapsto\tau A$ is a bijection between block-index
schedules and physical schedules contained in $\tau\Nzero$.  If sampling
periods vary, the same block-index sequence represents different physical
schedules.

The ordinary consecutive-name entropy is
\begin{equation}\label{eq:ordinary-partition-entropy}
 h(\C;Q)=\lim_{N\to\infty}\frac1{N\tau}\log\#W_N(\C;Q),
\end{equation}
whose limit exists by \eqref{eq:word-submult} and Fekete's lemma.


\section{Fixed-period feedbacks and the two optimizations}\label{sec:selectors}

Use the control setting of Section~\ref{sec:control}, now with a
finite control set $U$. Fix $\tau\ge1$ and write $W=U^\tau$.
For $w\in W$, its closed safety domain and endpoint map are
\begin{equation}\label{cc-eq:safe}
 K_w=\{x\in Q:\varphi(j,x,w)\in Q\ (0\le j\le\tau)\},
 \qquad g_w(x)=\varphi(\tau,x,w).
\end{equation}
The domains cover $Q$ by controlled invariance.

The class $\Ssafe_\tau$ consists of \emph{all} Borel selectors $s:Q\to W$
such that $x\in K_{s(x)}$. Its closed-loop block map is
$T_s(x)=g_{s(x)}(x)$. All Borel mixtures of legal selectors belong to
$\Ssafe_\tau$. No continuity or dependence on finitely many coordinates
is imposed. A selector is used without an additional clock, register,
or feedback memory.

For a finite window $F$, let
\[
 \Patt_s(F)=\left|\{(s(T_s^k x))_{k\in F}:x\in Q\}\right|.
\]
For $A=(a_i)_{i\ge1}\in\mathscr S$, define
\[
 h_A(s)=\limsup_{n\to\infty}\frac{1}{n\tau}
       \log\Patt_s(\{a_1,\ldots,a_n\}),\qquad
 h^*(s)=\lim_{n\to\infty}\frac{1}{n\tau}
       \log\max_{|F|=n}\Patt_s(F).
\]
The same submultiplicativity used in Section~\ref{sec:control} gives
the second limit. Normalization is per sampled control block, divided
by its fixed length; it is \emph{not} division by elapsed time $\tau a_n$.
Put
\begin{equation}\label{cc-eq:LM}
 L_\tau=\sup_A\inf_{s\in\Ssafe_\tau}h_A(s),\qquad
 M_\tau=\inf_{s\in\Ssafe_\tau}\hpat(s).
\end{equation}
Always $L_\tau\le M_\tau$. The pattern infimum is a minimum by the known integer-logarithm quantization; a self-contained argument and its precise fixed-period consequence are given in Appendix~\ref{app:minimum}.

The counts use actual closed-loop trajectories. If instead one takes
the closure of their one-sided name set, the finite patterns are
unchanged: a cylinder specifying finitely many coordinates is open, so
if it meets that closure, it meets the original name set. Taking a name
closure therefore introduces no extra finite-window witnesses.

\begin{lemma}\label{cc-lem:merge}
The infima in \eqref{cc-eq:LM}, and the inner infimum at any fixed $A$, are
unchanged if one allows every finite Borel safe partition with a
control block assigned to each atom, and counts atom names rather than
control-block names.
\end{lemma}
\begin{proof}
Merge all atoms assigned the same block. Their finite union is Borel
and remains inside that block's safety domain. The resulting selector
has exactly the same closed-loop state map. The map from original atom
names to block names is coordinatewise and onto the actual block-name
patterns, so every block-pattern count is no larger than the original
atom-pattern count. Conversely, the nonempty fibers of any legal
selector form a legal finite Borel safe partition whose atom and
block counts agree. Both directions of the infimum equality follow.
\end{proof}



\begin{example}[Merging observation labels and merging controls]
\label{ex:physical-merge}
Let $Q=\{0,1\}^{\Nzero}\cup\{2^\infty\}$ and let the controls
be $0,1,2$. A control equal to $x_0$ sends $x$ to $\sigma x$;
every other control sends it to an isolated cemetery point.
These maps are continuous, and safety forces the unique selector
$s(x)=x_0$. Every nonempty $n$-position window has $2^n+1$ actual control
patterns. Replacing the observed label $2$ by $0$ leaves exactly
$2^n$ patterns, so both observations have rate $\log2$.
The merged observed zero fiber, however, contains both $[0]$ and
$2^\infty$. The first set requires physical control $0$, whereas
the latter point requires physical control $2$. No one control is
safe on their union. Thus a symbolic factor may preserve the entropy
without defining a safe merged control. Lemma~\ref{cc-lem:merge}
applies because the atoms it merges were already assigned the
\emph{same} physical control block.
\end{example}


\section{Inverse capacities and maximal-pattern entropy}\label{sec:capacity}

\subsection{Branch measures, lists, and actual fibers}
We use the fixed-period rates of Section~\ref{sec:selectors}.
In particular $p_s^*(n)=\max_{|F|=n}\Patt_s(F)$ and
$h^*(s)=\lim_n(n\tau)^{-1}\log p_s^*(n)$.

Fix a period $\tau\ge1$ and a finite collection $W$ of admissible
control blocks, with $q=|W|$. For each $u\in W$ let $K_u\subset Q$
be its closed safety domain, including all intermediate-time safety
conditions, and $g_u:K_u\to Q$ its continuous endpoint map.
Assume these domains cover the nonempty compact metric space $Q$.
The feedback class and optimization quantities in this section are
relative to this specified block set $W$, as are its finite atom
refinements. For the system optimization of Section~\ref{sec:selectors},
we take $W=U^\tau$.
The legal class $\mathcal S_\tau$ consists of all Borel maps
$s:Q\to W$ with $x\in K_{s(x)}$. Define
$T_s(x)=g_{s(x)}(x)$. Reusing this same map at each block boundary
gives an infinite safe trajectory for every state. No continuity of
$s$ or of $T_s$ is imposed.

Let $\mu$ be a full-support Borel probability on $Q$ and suppose
\begin{equation}\label{eq:cap-density}
 \nu_u=(g_u)_*(\mu|_{K_u})\ll\mu,\qquad
 \rho_u=\frac{d\nu_u}{d\mu}\in L^\infty(\mu).
\end{equation}
The positive bounded operator $L_u:L^\infty(\mu)\to L^\infty(\mu)$
is specified by
\begin{equation}\label{eq:cap-transfer}
 \int hL_uf\,d\mu=\int_{K_u}h(g_ux)f(x)\,d\mu(x)
 \quad(h,f\text{ bounded Borel}).
\end{equation}
Indeed, the signed pushforward of $f\mu|_{K_u}$ is dominated by
$\|f\|_\infty\nu_u$; its density defines $L_uf$ and
$|L_uf|\le\|f\|_\infty\rho_u$. Changing $f$ on a $\mu$-null
set changes no pushforward measure, so this is well defined on
$L^\infty$ classes. In particular $L_u\one=\rho_u$.

Fix $2\le k\le q$ and put $l=k-1$. A list $S\subset W$ is a
\emph{global} set of $l$ labels, not a state-dependent choice.
Write $A_S=\sum_{u\in S}L_u$, and define
\begin{equation}\label{eq:cap-bn}
 b_0=1,\qquad
 b_n=\max_{|S_0|=\cdots=|S_{n-1}|=l}
 \|A_{S_{n-1}}\cdots A_{S_0}\one\|_\infty.
\end{equation}
Products are ordered in the forward pushforward order. Positivity gives
$\opnorm{P}=\|P\one\|_\infty$ for every positive product $P$.
Consequently $b_{n+m}\le b_nb_m$. If $b_h<1$ for some $h\ge1$,
choose $b_h<\theta^h<1$ and set
\begin{equation}\label{eq:cap-exp}
 H=\max_{0\le j<h}b_j\theta^{-j}.
 \qquad b_n\le H\theta^n\quad(n\ge0).
\end{equation}
This follows by writing $n=ah+j$ and using submultiplicativity.
The proof length $h$ is not a change of control period.

For a fixed feedback let
\[
 E_w(s)=\{x:s(T_s^jx)=w_j\ (0\le j<n)\},\qquad
 F_n(s)=\{w\in W^n:E_w(s)\ne\varnothing\}.
\]
The $E_w(s)$ are Borel, pairwise disjoint, and cover $Q$.
For $w=w_0\cdots w_{n-1}$ put
$L_w=L_{w_{n-1}}\cdots L_{w_0}$ and
$P_F=\sum_{w\in F}L_w$.

\begin{lemma}[Actual-fiber capacity]\label{lem:fiber-cap}
For every legal Borel feedback and every $n\ge1$,
\begin{equation}\label{eq:fiber-mass}
 1\le\int P_{F_n(s)}\one\,d\mu.
\end{equation}
If $F_n(s)$ is covered by $t$ boxes
$S^a_0\times\cdots\times S^a_{n-1}$ whose lists have size at
most $l$, then $1\le t b_n$.
\end{lemma}
\begin{proof}
Fix a nonempty fiber $E_w$ and let $\eta_j$ be the pushforward
of $\mu|_{E_w}$ after its first $j$ actual blocks. Set
$f_0=\one_{E_w}$. We claim inductively that
$\eta_j=f_j\mu$ and $f_{j+1}=L_{w_j}f_j$. Indeed,
$\eta_j$ is concentrated on $K_{w_j}$, and for every bounded Borel $h$,
\[
 \int h\,d\eta_{j+1}
 =\int_{K_{w_j}}h(g_{w_j}x)f_j(x)\,d\mu(x)
 =\int hL_{w_j}f_j\,d\mu.
\]
The boundedness of each $L_u$ permits the induction. None of these
restrictions removes mass, because all states in $E_w$ follow the
specified safe word. Thus $\int L_w\one_{E_w}\,d\mu=\mu(E_w)$.
Positivity gives $L_w\one_{E_w}\le L_w\one$; summing over the
disjoint fibers proves \eqref{eq:fiber-mass}. The set $F_n(s)$ includes
null but nonempty fibers as well.

Complete every nonempty list to size $l$. Distributivity and positivity
give
\[
 P_{F_n(s)}\one\le\sum_{a=1}^t
 A_{S^a_{n-1}}\cdots A_{S^a_0}\one.
\]
Repeated boxes and overlaps only enlarge this upper bound. Unsafe
words in a box are subject to the original safety gates in $L_u$.
Integrating and using \eqref{eq:cap-bn} yields $1\le tb_n$.
\end{proof}

\subsection{The finite combinatorial input and its application}
For $F\subset W^n$, let $C_l(F)$ be its minimum covering number by
$l$-list boxes. Given a $k$-element set $V\subset W$, let $C_V(F)$
be the minimum number of boxes of the form
\[
 \prod_{j=0}^{n-1}(W\setminus\{v_j\}),\qquad v_j\in V,
\]
needed to cover $F$. This is a cover, since $k\ge2$.

We use the following already established finite-set lemma. For every
$k\ge2$ and $b>0$ there are $c>0$ and $N$ such that if
$G\subset\{0,1,\ldots,k\}^n$, $n\ge N$, needs at least
$e^{bn}$ boxes omitting one nonzero letter per coordinate, then some
$J\subset\{0,\ldots,n-1\}$ with $|J|\ge cn$ has
\begin{equation}\label{eq:HY-cube}
 G|_J\supset\{1,\ldots,k\}^{J}.
\end{equation}
This is Kerr--Li's finite covering lemma
\cite[Lemma~3.3]{KL2007}. Huang and Ye recover the same conclusion
from their sharper estimates in \cite[Lemma~4.1 and Remark~4.2,
pp.~1701--1705]{HuangYe2009}. Its constants depend only on $b,k$;
its use here requires no dynamical assumptions on the closed loop.

For clarity, the entire reduction from list boxes to that lemma is
included. Choose a minimum omitted-letter cover for every $V$ and
intersect all these covers. The intersections cover $F$, with at most
$\prod_V C_V(F)$ members. In a nonempty intersection each coordinate
contains at most $k-1$ letters: if it contained a $k$-set $V_0$, its
factor from the $V_0$ cover would have omitted a letter of $V_0$.
Hence
\begin{equation}\label{eq:cap-intersect}
 C_l(F)\le\prod_{V\in\binom Wk}C_V(F).
\end{equation}
For actual names Lemma~\ref{lem:fiber-cap} gives
$C_l(F_n(s))\ge b_n^{-1}$. Under \eqref{eq:cap-exp}, for all
sufficiently large $n$ this lower bound is at least
$e^{\eta n/2}$, where $\eta=-\log\theta>0$. With
$J_0=\binom qk$, \eqref{eq:cap-intersect} implies that some $V$
satisfies
\[
 C_V(F_n(s))\ge e^{\eta n/(2J_0)}.
\]
Code the members of $V$ by $1,\ldots,k$ and all other letters by
$0$. Omitted-letter boxes pull back exactly to the boxes defining
$C_V$, so the coded set has precisely that covering number. Applying
\eqref{eq:HY-cube} and undoing this coding yields
$F_n(s)|_J\supset V^J$. No distinct members of $V$ were merged.

\begin{theorem}[Product inverse-capacity criterion]\label{thm:capacity-main}
Assume \eqref{eq:cap-density}. If $b_h<1$ for some $h\ge1$, then
there are $\alpha>0$ and $N$, depending only on the branches,
$\mu,k$, such that, for every legal Borel feedback $s$ and every
$n\ge N$, its actual length-$n$ names have a $k$-valued cube
projection on at least $\alpha n$ coordinates. Consequently
\begin{equation}\label{eq:cap-pattern}
 p_s^*(m)\ge k^m\quad(m\ge1),\qquad
 h^*(s)\ge\frac{\log k}{\tau}.
\end{equation}
The same lower bound holds for all finite Borel atom refinements.
If $k=r\ge2$ is the minimum number of safety domains covering $Q$,
then
\begin{equation}\label{eq:cap-M}
 M_\tau=\min_s h^*(s)=\frac{\log r}{\tau},
\end{equation}
both over unrefined selectors and over all finite Borel safe partitions.
\end{theorem}
\begin{proof}
The preceding argument supplies $\alpha,N$ independent of $s$.
Given $m$, choose $n$ so large that $\alpha n\ge m$ and retain
$m$ of the independent positions. All $k^m$ patterns are actual
feedback names, proving \eqref{eq:cap-pattern} with the prescribed
normalization by $m\tau$.

For $k=r$, choose a minimum subcover $K_{u_1},\ldots,K_{u_r}$ and
select $u_j$ on $K_{u_j}\setminus\bigcup_{i<j}K_{u_i}$. This is a
Borel state function, legal at every point. Each chosen block stays
safe at its intermediate times and ends in $Q$; applying the same
rule at the next state gives infinite safety. It uses $r$ labels and
has at most $r^m$ patterns on any $m$ positions. This proves
\eqref{eq:cap-M} and attainment. For a refined partition, projection
to its physical control labels is onto the actual coarse patterns
and preserves the dynamics. The lower bound therefore passes to
every refinement, while the unrefined selector remains an available
upper-bound witness in the enlarged class.
\end{proof}

The theorem asserts a uniform \emph{density}, not a common set of
independent positions. The sets $J$ and $V$ may depend on both $s$
and $n$. In particular it proves no equality $L_\tau=M_\tau$
without an additional common-sampling argument.

\subsection{Single-step scalar and joint-density criteria}
Define
\begin{equation}\label{eq:cap-Theta}
 \Theta_k=\left\|\max_{|S|=k-1}\sum_{u\in S}\rho_u\right\|_\infty.
\end{equation}
If $\Theta_k<1$, then $b_1=\Theta_k<1$ and
Theorem~\ref{thm:capacity-main} applies. If the branch capacities obey
$\mu(K_u\cap g_u^{-1}E)\le c_u\mu(E)$ for every Borel $E$,
then $\rho_u\le c_u$ almost everywhere, and the sum of the largest
$k-1$ of the $c_u$ bounds $\Theta_k$. Thus a scalar-capacity sum
smaller than one is a sufficient special case. The order of taking
the branch supremum and summing matters: joint densities may pass
the test while the individual optimal scalar bounds do not.

The joint-density case also has a direct quantitative proof worth
recording. For nonempty $F\subset W^n$, let $d_k(F)$ be the largest
$|J|$ for which there are sets $V_j\subset W$, $|V_j|=k$, such that
$F|_J\supset\prod_{j\in J}V_j$. The sets $V_j$ may differ with $j$.
Set $d_k(\varnothing)=-1$. Put $\theta=\Theta_k<1$, $C_0=\|\sum_u\rho_u\|_\infty$ and
$B=C_0\binom qk$. For nonnegative $f$ let
$\Phi f=\max_{|S|=k-1}\sum_{u\in S}L_uf$ and
$\mathcal Tf=\sum_uL_uf$. Their norms are bounded by $\theta$
and $C_0$, respectively.

\begin{proposition}[Weighted shattering estimate]\label{prop:weighted-sauer}
With $H(n,-1)=0$, for every $F\subset W^n$,
\begin{equation}\label{eq:cap-weighted}
 \opnorm{P_F}\le
 H(n,d_k(F)),\qquad
 H(n,d)=\sum_{j=0}^{\min(d,n)}\binom nj B^j\theta^{n-j}.
\end{equation}
For actual names this gives $d_k(F_n(s))>\alpha n$, uniformly in
every legal feedback, whenever $0<\alpha<1/2$ satisfies
\begin{equation}\label{eq:cap-alpha}
 \log\theta+\mathfrak h(\alpha)+\alpha\log(B/\theta)<0,
 \qquad\mathfrak h(t)=-t\log t-(1-t)\log(1-t).
\end{equation}
\end{proposition}
\begin{proof}
Take $H(n,-1)=0$ and treat the empty word at $n=0$ separately.
Split $F$ at its \emph{first} letter: $F_a=\{v:av\in F\}$,
$U_0=\bigcup_aF_a$, and $I_S=\bigcap_{a\in S}F_a$ for
$|S|=k$. If $d=d_k(F)$, then $d_k(U_0)\le d$ and
$d_k(I_S)\le d-1$: independence of $d$ tail positions in the latter
would also make the first position independent and give $d+1$.
For each $v\in U_0$ the set of its permitted first letters has
either at most $k-1$ members, or contains a $k$-set $S$. Applying
the positive tail product $L_v$ and summing gives
\[
 P_F f\le P_{U_0}\Phi f+\sum_{|S|=k}P_{I_S}\mathcal T f.
\]
This respects the order $L_{av}=L_vL_a$; no operators are
commuted. Induction and Pascal's identity give
$\theta H(n-1,d)+BH(n-1,d-1)=H(n,d)$. Positive domination
handles signed $f$ and proves \eqref{eq:cap-weighted}.

By actual-fiber mass, $1\le H(n,d_k(F_n(s)))$.
The cover property gives $C_0\ge1$, and positivity gives
$0<\theta\le C_0\le B$. For $d\le\alpha n$, therefore,
$B/\theta\ge1$ and
$\sum_{j\le\alpha n}\binom nj\le e^{n\mathfrak h(\alpha)}$
give
$H(n,d)\le\exp(n[\log\theta+\mathfrak h(\alpha)+
\alpha\log(B/\theta)])<1$, a contradiction.
The binomial bound follows by comparing each term to its probability
under independent Bernoulli$(\alpha)$ trials; those probabilities
sum to at most one. Condition~\eqref{eq:cap-alpha} is possible
because its last two terms tend to zero as $\alpha$ tends to zero.
\end{proof}

\subsection{A smooth joint-density example}\label{subsec:smooth}
On $Q=[0,1]$ take three full inverse branches with
\begin{align}
 l_u&=\frac{19u}{60},\qquad
 \rho_u(y)=\frac{11}{30}+\frac7{30}
            \cos(2\pi(y-u/3)),\label{eq:smooth-rho}\\
 \psi_u(y)&=l_u+\frac{11}{30}y+
 \frac7{60\pi}\big[\sin(2\pi(y-u/3))-\sin(-2\pi u/3)\big].
 \nonumber
\end{align}
Here $g_u=\psi_u^{-1}$, for $u=0,1,2$.
Their derivatives lie between $2/15$ and $3/5$, so the branches
are strictly increasing and expanding in the forward direction.
The closed safety intervals are
$[0,22/60]$, $[19/60,41/60]$, and $[38/60,1]$.
They cover $Q$ and each is needed: the middle interval alone covers
the open gap between $22/60$ and $38/60$, while the two endpoints
require the first and third intervals. Hence $r=3$.

With Lebesgue probability, the densities are exactly
$\psi'_u=\rho_u$. The three cosine terms sum to zero. The sum of
two is the negative of the remaining cosine, whose maximum is one;
therefore
\[
 \Theta_3=\frac{22}{30}+\frac7{30}=\frac{29}{30}<1.
\]
Each branch has optimal scalar bound $3/5$, since its continuous
density reaches that value. The largest two scalar bounds sum to
$6/5$. The joint criterion consequently gives $M_1=\log3$ although
that Lebesgue scalar criterion fails.

To place these interval branches in a continuous compact ambient
system, take $X=[-1,2]$. Extend each $g_u$ to the left of its
safety interval $[a_u,b_u]$ by $-\min\{a_u-x,1\}$ and to its
right by $1+\min\{x-b_u,1\}$. At the endpoints these formulas
match $0$ and $1$. Outside the closed safety interval the image
lies outside $Q$, so this extension changes neither safety domains
nor legal dynamics. The same construction applies to full monotone
interval branches elsewhere in the paper.

The fixed-reference separation proved here is not, by itself, a
separation from every full-support reference probability. The bounded
positive-density obstruction is recorded in
Appendix~\ref{app:measure}; its extension to arbitrary full-support
probabilities remains unresolved for this smooth example.

\subsection{Product contraction without a single-step reference test}
\label{subsec:product-example}
Let $Q=(\mathbb Z/3\mathbb Z)^{\Nzero}$, let $P_u$ add $u$
to each coordinate, and use $g_u=P_u\sigma$ at period one. Put
\[
 K_0=[02]\cup[12]\cup[22],\quad K_1=[00]\cup[20],\quad
 K_2=[01]\cup[10]\cup[11]\cup[12]\cup[21].
\]
The only overlap is $[12]$. The cylinders $[02],[00],[01]$
require controls $0,1,2$, respectively, so $r=3$.
On $[12]$, choosing $0$ or $2$ changes the successor's first
letter from $2$ to $1$. Every Borel mixture on this entire overlap
is allowed. Illegal controls go to an isolated fixed cemetery point;
the resulting maps on $Q\sqcup\{\dagger\}$ are continuous.

For fair Bernoulli probability, the operators leave first-coordinate
functions invariant and have matrices $M_u/3$, where
\[
 M_0=\begin{pmatrix}0&0&0\\0&0&0\\1&1&1\end{pmatrix},\quad
 M_1=\begin{pmatrix}0&0&0\\1&0&1\\0&0&0\end{pmatrix},\quad
 M_2=\begin{pmatrix}1&1&1\\0&1&0\\0&1&0\end{pmatrix}.
\]
This follows from splitting the source measure by its first letter,
and retaining exactly the allowed two-letter cylinders. For a
positive product the full $L^\infty$ norm equals its norm on
$\one$, whose entire orbit stays in this invariant space. Thus this
finite representation computes the full product norms, not an
approximation to arbitrary Borel functions.

Let $A=M_0+M_1$, $B=M_0+M_2$, $C=M_1+M_2$.
The vectors after two of these unnormalized factors are
\[
\begin{array}{ccc}
(0,3,5)&(0,7,8)&(0,4,7)\\
(5,2,7)&(8,1,9)&(7,3,10)\\
(5,5,2)&(8,8,1)&(7,7,3).
\end{array}
\]
They are coordinatewise dominated by one of the vectors in
\[
 V=\{(0,7,8),(7,3,10),(7,7,3),(8,1,9),(8,8,1)\}.
\]
For $v=(p,q',r')$,
$Av=(0,p+r',p+q'+r')$,
$Bv=(p+q'+r',q',p+2q'+r')$, and
$Cv=(p+q'+r',p+q'+r',q')$.
For the five vectors of $V$, their largest resulting coordinate is,
respectively, $22,23,24,19,25$. Also $BCB\one=(17,8,25)$.
Consequently $b_3=25/27<1$. Submultiplicativity, rather than an
extrapolation from finite calculations, controls all lengths.
Theorem~\ref{thm:capacity-main} gives $M_1=\log3$, and every
unrefined feedback has that maximal-pattern rate. An explicit safe
selector chooses $0$ if $x_1=2$, otherwise $1$ if $x_1=0$ and
$x_0\ne1$, and otherwise $2$.

On the three target first-letter cylinders the densities are
$(0,0,1)$, $(0,2/3,1/3)$, and $(1,0,1/3)$.
Hence $\Theta_3=4/3$. This failure persists under \emph{every}
full-support probability $\nu$. If one of the safe branch measures
has no bounded $\nu$-density, the single-step criterion already
fails its prerequisites. Otherwise $g_0(K_0)\subset[2]$ and
$g_1(K_1)\subset[1]$ give disjoint density supports. At most two
densities are nonzero at a point, so their largest two sum to their
total. Its integral is
\[
 \sum_{u=0}^2\nu(K_u)=1+\nu([12])>1.
\]
Thus $\Theta_3(\nu)>1$ for every such $\nu$, including singular
or atomic probabilities. This excludes both original single-step
tests, but not all possible changes of presentation or iteration.

There is also a distinction from iterating the pointwise envelope
$\Phi$. On nonnegative functions the disjoint supports give
$\Phi=\sum_uL_u$. Under fair measure its first-letter matrix is
$\frac13\left(\begin{smallmatrix}1&1&1\\1&1&1\\1&2&1\end{smallmatrix}\right)$,
with positive eigenvector
$(1,1,(\sqrt{13}-1)/2)$ and eigenvalue
$(3+\sqrt{13})/6>1$. Hence $\|\Phi^n\one\|_\infty$ grows,
while every history of \emph{global} two-label lists decays.
Maximization by state and composition cannot be interchanged.


\section{Finite-cylinder safety and left-weight estimates}\label{sec:cylinder}

\subsection{Exact finite representation without restricting feedbacks}
Let $D$ be a finite alphabet of size $d\ge2$, $Q=D^{\Nzero}$,
and $\mu$ its fair Bernoulli probability. The finite controls have
legal maps $g_u=P_u\sigma$, where $p_u$ permutes $D$ and $P_u$
acts in every coordinate. Suppose
\[
 K_u=\bigcup_{\xi\in R_u}[\xi],\qquad R_u\subset D^h,
 \quad h\ge2,
\]
and the $K_u$ cover $Q$. Shorter cylinders can be completed to the
same length $h$. The continuous cemetery extension used above
applies to these clopen domains. The control period is always one.

For every bounded Borel function the full operator is
\begin{equation}\label{eq:cyl-operator}
 L_uf(y)=\frac1d\sum_{a\in D}
 \one_{R_u}(a,p_u^{-1}y_0,\ldots,p_u^{-1}y_{h-2})
 f(aP_u^{-1}y).
\end{equation}
Indeed, conditioned on $x_0=a$ the tail of $x$ still has law $\mu$,
and $P_u$ preserves that law. Substitution in
\eqref{eq:cap-transfer} gives this formula with the exact safety gate.

Index source and target vectors by $D^{h-1}$. For target $\eta$
and source $\xi$, set
\begin{equation}\label{eq:cyl-matrix}
 T_u(\eta,\xi)=\frac1d\#\left\{a\in D:
 \begin{array}{l}
 (a,p_u^{-1}\eta_0,\ldots,p_u^{-1}\eta_{h-2})\in R_u,\\
 (a,p_u^{-1}\eta_0,\ldots,p_u^{-1}\eta_{h-3})=\xi
 \end{array}\right\}.
\end{equation}
When $h=2$ the second tuple is just $a$.
Functions of the first $h-1$ coordinates form an invariant subspace
of every $L_u$, represented by $T_u$. Starting at $\one$, every
list product stays in this space. Every such cylinder has positive
measure and a positive product has norm equal to its value on $\one$.
Therefore, for every length and every history,
\begin{equation}\label{eq:cyl-fullnorm}
 \opnorm{A_{S_{n-1}}\cdots A_{S_0}}
 =\|T_{S_{n-1}}\cdots T_{S_0}\one\|_\infty,
 \qquad T_S=\sum_{u\in S}T_u.
\end{equation}
This equality is on the full $L^\infty$ space. It does not claim
that an arbitrary Borel function is approximated by a fixed finite
space, or that a Borel feedback uses a finite-state rule. The finite
indices are proof tools, not new state coordinates or a clock.

For two-letter domains, rows are target letters and columns are
source letters, with
\begin{equation}\label{eq:cyl-two}
 T_u(j,i)=d^{-1}\one_{R_u}(i,p_u^{-1}j),\qquad
 d(w^\top T_u)_i=\sum_{z:(i,z)\in R_u}w(p_uz).
\end{equation}
The structural conditions below use this source-column identity.
A left budget $w^\top T_S\le\theta w^\top$ is different from
changing a reference density, which leads instead to a right
inequality $T_Sa\le\theta a$. They cannot be interchanged.

\subsection{A strict source-weight criterion}
Partition $D=A\sqcup B\sqcup C$ into nonempty sets of sizes
$a,b,c$, so $d=a+b+c$. Use three controls with
\begin{equation}\label{eq:weight-land}
 p_0(C)\subset D\setminus A,\qquad p_1(B)\subset A,
 \qquad p_2(A)\subset D\setminus A.
\end{equation}
Choose arbitrary relations $E_0\subset A\times C$ and
$E_1\subset A\times B$, with $E_0\cup E_1\ne\varnothing$,
and define
\begin{align}
 R_0&=((D\setminus A)\times C)\cup E_0,\nonumber\\
 R_1&=((D\setminus A)\times B)\cup E_1,\label{eq:weight-domains}\\
 R_2&=(A\times D)\cup(D\times A),\qquad K_u=[R_u].\nonumber
\end{align}
Put $e=\max_{i\in A}|\{z:(i,z)\in E_1\}|$. Require only
the following counts and landing conditions:
\begin{equation}\label{eq:weight-count}
 s_-:=\max\left\{\frac c{b+c},\frac e{b+c-e}\right\}
 <s_+:=\frac{\min\{a,c\}}b.
\end{equation}
These are direct conditions on the original allowed cylinders.

\begin{theorem}[Strict source budget]\label{thm:source-weight}
For every system satisfying \eqref{eq:weight-land}--\eqref{eq:weight-count},
$r=3$, all histories of global two-label lists have exponentially
decaying inverse capacity, and every legal unrefined Borel feedback
has $h^*(s)=M_1=\log3$. Every finite Borel atom refinement has
$h^*\ge\log3$. Moreover, no full-support reference probability
can directly satisfy the joint-density or scalar single-step
criterion for $k=3$ on these original branches.
\end{theorem}
\begin{proof}
Take $s=(s_-+s_+)/2$, $\gamma=1+s$, and let $w_i=\gamma$
on $A$, $w_i=1$ elsewhere. Write $H=\sum_iw_i=d+sa$.
Before division by $d$, the weighted contributions of controls
$0,1,2$ in a source column are
\[
\begin{array}{c|ccc}
 &0&1&2\\\hline
 i\in A&e_0(i)\le c&\gamma e_1(i)\le\gamma e&H\\
 i\notin A&c&\gamma b&a.
\end{array}
\]
Each entry follows from the original domains and
\eqref{eq:weight-land}. For example, on $A$ control $2$ permits
every tail, hence the sum of all weights irrespective of its
permutation. Off $A$, its permitted tails lie in $A$ and land
outside $A$, each with weight one.

Since $sb<a,c$, both $c$ and $\gamma e\le\gamma b$ are smaller
than $H$. Thus every two-label sum on $A$ is bounded by
$H+\max\{c,\gamma e\}$. Define
\[
 \theta=\max\left\{
 \frac{H+\max\{c,\gamma e\}}{d\gamma},
 \frac{c+\gamma b}{d},\frac{c+a}{d},\frac{\gamma b+a}{d}
 \right\}.
\]
We have $d\gamma-H=s(b+c)>c,\gamma e$ because $s>s_-$.
The other three terms are smaller than one by $sb<a,c$ and
$b>0$. Hence $0<\theta<1$ and
$w^\top T_S\le\theta w^\top$ for every global two-label list.
Repeatedly applying this single inequality gives
\[
 \|T_{S_{n-1}}\cdots T_{S_0}\one\|_\infty
 \le w^\top T_{S_{n-1}}\cdots T_{S_0}\one
 \le H\theta^n,
\]
since $\min_iw_i=1$. Equation~\eqref{eq:cyl-fullnorm} controls
the complete capacity for all lengths.

If $x_0\in A$, control $2$ is legal; otherwise its next letter
in $A,B,C$ requires, respectively, $2,1,0$. Thus the domains
cover $Q$. Taking $x_0\notin A$ and $x_1$ successively in
$C,B,A$ gives cylinders on which each of $0,1,2$ alone is legal,
so $r=3$. A clopen attaining selector chooses $2$ when $x_0\in A$
or $x_1\in A$, otherwise $1$ when $x_1\in B$, otherwise $0$.
It is safe, single-valued, and stays in $Q$ at each step; the same
rule at every successor gives all-time safety. Applying
Theorem~\ref{thm:capacity-main} proves all pattern conclusions,
including the lower bound for arbitrary refinements.

Let $O=[E_0\cup E_1]$. It is nonempty and clopen, and safety
multiplicity is exactly $1+\one_O$. For a full-support probability
$\nu$, either some branch fails the bounded-density prerequisite,
or all densities exist. In the latter case the images of controls
$0$ and $1$ lie, respectively, in $[D\setminus A]$ and $[A]$;
their supports are disjoint. The largest two densities consequently
sum to all three almost everywhere. Their integral is
$\sum_u\nu(K_u)=1+\nu(O)>1$. Thus $\Theta_3(\nu)>1$.
The scalar condition implies $\Theta_3<1$ and also fails.
\end{proof}

For instance
$1\le b\le a\le b+c$ and $c>b$ imply
\eqref{eq:weight-count} for every possible $e\le b$:
$e/(b+c-e)\le b/c<1$, $c/(b+c)<1$, while $s_+\ge1$.
Every allowed overlapping relation and permutation landing pattern
is covered. Constants can be uniform when these counts have fixed
bounds.

For a concrete example take
$D=\{0,1,2,3,4\}$, $A=\{0,1\}$, $B=\{2\}$,
$C=\{3,4\}$, $E_0=\{(0,3),(1,4)\}$,
$E_1=\{(0,2)\}$. Let $p_0$ be the identity,
$p_1=(0\ 2)$, and specify $p_2$ by
$0\mapsto2$, $1\mapsto3$, $2\mapsto0$, $3\mapsto4$,
$4\mapsto1$. Here $s=4/3$, $\gamma=7/3$, $H=23/3$,
$\theta=13/15$, and
$b_n\le(23/3)(13/15)^n$. The target densities at first letters
$0,1,2,3,4$ are, respectively,
\[
 (0,4/5,2/5),\quad(0,0,2/5),\quad(0,0,1),\quad
 (4/5,0,1),\quad(4/5,0,2/5).
\]
Thus $\Theta_3=9/5$. In $[14]$ the two legal successors have
first letters $4$ and $1$. No common successor is present. The
natural $A/B/C$ merging also fails to commute with the listed
permutations; this observation does not exclude every possible factor.

\subsection{Cross-layer dispersion without a strict left budget}
Let $U$ have size $q\ge2$. Partition $D=\bigsqcup_{t\in U}I_t$
into nonempty source classes, and let $m(i)=t$ on $I_t$.
For each $i$ partition the tail alphabet by
\[
 D=\bigsqcup_{u\ne m(i)}V_{i,u},\qquad F_i\subset D.
\]
Control $u\ne m(i)$ is the primary legal control on $V_{i,u}$,
while control $m(i)$ is also legal on $F_i$. Precisely,
\begin{equation}\label{eq:spread-relation}
 R_u=\{(i,z):u\ne m(i),\ z\in V_{i,u}\}
 \cup\{(i,z):u=m(i),\ z\in F_i\}.
\end{equation}
The permutations $p_u$ are arbitrary subject to the count below.
Assume
\begin{align}
 a_0&=\min_{i,u\ne m(i)}(|V_{i,u}|-|F_i|)>0,
 \label{eq:spread-size}\\
 a_1&=\min_{i,t}\left[d-
       \sum_{u\ne m(i)}|p_u(V_{i,u})\cap I_t|\right]>0.
 \label{eq:spread-land}
\end{align}
The second count includes multiplicities of coincident landings.
It says that the primary successor mass cannot be confined to any
one source class. Set $\eps=a_0/d$, $\eta=a_1/d$, and
$\zeta=1-\eps\eta\in(0,1)$.

\begin{theorem}[Dispersion criterion]\label{thm:spread}
Under \eqref{eq:spread-relation}--\eqref{eq:spread-land}, $r=q$ and
\begin{equation}\label{eq:spread-bound}
 b_n\le d\zeta^{\lfloor n/2\rfloor}
 \le\frac d{\sqrt\zeta}(\sqrt\zeta)^n\quad(n\ge0).
\end{equation}
Every legal unrefined Borel feedback has $h^*(s)=M_1=\log q$;
all finite Borel atom refinements have $h^*\ge\log q$.
There is no $w\in(0,\infty)^d$ and $\vartheta<1$ such that
$w^\top T_S\le\vartheta w^\top$ for every global
$q-1$ label list. A common nonstrict left budget does exist.
\end{theorem}
\begin{proof}
Write $S_t=U\setminus\{t\}$ and $T^{(t)}=T_{S_t}$.
Its column mass is exactly
\[
 c_i(t)=\sum_jT^{(t)}(j,i)=
 \begin{cases}
 1,&m(i)=t,\\
 1-|V_{i,t}|/d+|F_i|/d\le1-\eps,&m(i)\ne t.
 \end{cases}
\]
Thus every $T^{(t)}$ has induced $l^1$ norm at most one.
Let $Q_i(j)=d^{-1}\sum_{u\ne m(i)}\one_{p_u(V_{i,u})}(j)$.
It has total mass one, and by \eqref{eq:spread-land}
$\sum_{j\notin I_t}Q_i(j)\ge\eta$ for every $i,t$.

Fix two global lists $S_s,S_t$. The source-$i$ column mass of
$T^{(t)}T^{(s)}$ is $\sum_jc_j(t)T^{(s)}(j,i)$.
If $i\notin I_s$, it is at most $1-\eps\le1-\eps\eta$.
If $i\in I_s$, its first column is exactly $Q_i$, and at least
$\eta$ of that mass lands outside $I_t$, where the next list loses
at least $\eps$. Its mass is again at most $\zeta$.
Hence every pair has $l^1$ norm at most $\zeta$.
Grouping arbitrary histories into pairs proves their $l^1$ norm
is at most $\zeta^{\lfloor n/2\rfloor}$. Applying the product
to the all-one vector and using \eqref{eq:cyl-fullnorm} gives
\eqref{eq:spread-bound}. No restrictions on switching order or
on the feedback are used.

For every control $u$ choose $i$ with $m(i)\ne u$. Condition
\eqref{eq:spread-size} gives a $z\in V_{i,u}\setminus F_i$,
on whose cylinder only $u$ is legal. This proves $r=q$.
The unique primary control on each two-letter state is a clopen
selector, legal at every state and at each of its future successors.
The capacity theorem now proves the entropy statements.

If a strict common left budget existed, take a coordinate $i$ with
minimum $w_i$ and select the already assumed inequality with
$t=m(i)$. That column is $Q_i$, of mass one, and
\[
 (w^\top T^{(m(i))})_i\ge w_i>
 \vartheta w_i,
\]
a contradiction. This checks an inequality that was assumed for
every global list; it does not alter lists by source in an actual
product. Finally $\one^\top T^{(t)}\le\one^\top$ gives the
nonstrict budget, with strict contraction after two steps.
\end{proof}

An explicit family is obtained with $q\ge2$, $m\ge2$, and
$d=(q-1)m$. Partition tails into $q-1$ blocks of size $m$.
For each source, biject its $q-1$ primary labels with these blocks,
and let every $p_u$ preserve each tail block. Choose arbitrary
$F_i$ of size at most $m-1$, and any partition into $q$ nonempty
source classes whenever such a partition exists. The primary
landing mass is uniform, so $a_0\ge1$ and
$a_1=d-\max_t|I_t|>0$. This supplies arbitrary overlapping
choices and nontrivial permutations, within the same estimate.

For a small witness use $D=\{0,1,2,3\}$, $U=\{0,1,2\}$,
$I_0=\{0,1\}$, $I_1=\{2\}$, $I_2=\{3\}$, and
$F_i=\{0\}$ for every source. At sources $0,1$ the primary
labels $1,2$ use tails $\{0,1\},\{2,3\}$; at source $2$ the
labels $0,2$ use those tails; at source $3$ the labels $0,1$ do so.
Let $p_0$ be the identity, $p_1=(0\ 1)$, $p_2=(2\ 3)$.
The densities at target first letters are
\[
 (1,1/2,1/4),\quad(1/2,3/4,0),\quad
 (0,1/4,3/4),\quad(0,1/4,3/4).
\]
Thus $\Theta_3=3/2$, but $\eps=1/4$, $\eta=1/2$ give
$b_n\le4(7/8)^{\lfloor n/2\rfloor}$ and $M_1=\log3$.
On $[00]$ controls $0,1$ have different successors. The preceding
minimum-coordinate proof excludes all strict left budgets, not
just a numerically tested weight. This example does not inherit
the all-reference-probability exclusion of
Theorem~\ref{thm:source-weight}; its proved exclusion concerns
the specified matrix representation.


\section{Contraction with gain columns}
\label{sec:gainloss}

The preceding dispersion theorem retains a common nonstrict column
budget. We now describe a class in which every source has a gain
list, no positive common nonstrict left budget exists, and all global
list histories nevertheless contract. The condition is on original
safety sets and their cardinalities, not on a computed product norm.

\subsection{Original safety data and the count condition}
Let $q=|U|\ge3$ and partition $D=\bigsqcup_{t\in U}L_t$ into
nonempty source classes. Write $d=|D|$,
$\kappa(i)=t$ on $L_t$, and $M=\min_t|L_t|$.
For every source $i$ choose a spare label $m(i)\ne\kappa(i)$.
For each $u\in A_i:=U\setminus\{\kappa(i),m(i)\}$ choose
mutually disjoint tail sets $H_{i,u}$, all of size $a\ge1$.
Put
\[
 H_i=\bigsqcup_{u\in A_i}H_{i,u},\quad b=(q-2)a,\quad
 B_i=D\setminus H_i.
\]
Also choose an arbitrary spare tail set $F_i$ of size $f$.
It may overlap any of the primary tail sets. Define the original
allowed relation by
\begin{align}
 R_u={}&\{(i,z):u=\kappa(i),\ z\in B_i\}\nonumber\\
 &\cup\{(i,z):u\in A_i,\ z\in H_{i,u}\}\label{eq:gain-relation}\\
 &\cup\{(i,z):u=m(i),\ z\in F_i\},\qquad K_u=[R_u].\nonumber
\end{align}
Every two-letter state has one primary control and possibly a spare.
The permutations $p_u$ may be arbitrary. Use $g_u=P_u\sigma$
with the same continuous cemetery extension and period one.
Require
\begin{equation}\label{eq:gain-counts}
 M>b,\qquad a<f,\qquad b+f<d,
\end{equation}
and
\begin{equation}\label{eq:gain-margin}
 2d(f-a)+(f-a)^2<[d-(q-1)a][M-(q-2)a].
\end{equation}
These conditions contain only cardinalities of the original source
classes, small primary tails, and spare tails. The theorem applies
to every placement of these tails and every permutation, not only
to matrices previously known to be stable.

Set
\begin{align}
 \gamma&=(f-a)/d,& C&=1+\gamma,&
 \eps&=1-(b+f)/d,& \eta&=(M-b)/d,\nonumber\\
 \zeta&=\max\{C(1-\eps),C^2-(\eps+\gamma)\eta\},&
 \theta&=\sqrt\zeta,&H&=dC/\sqrt\zeta.
 \label{eq:gain-constants}
\end{align}

\begin{theorem}[Gain--loss safety criterion]\label{thm:gainloss}
For every system satisfying \eqref{eq:gain-relation}--\eqref{eq:gain-margin},
$r=q$, $0<\zeta<1$, and the complete inverse capacities satisfy
\begin{equation}\label{eq:gain-bn}
 b_n\le dC^{n\bmod2}\zeta^{\lfloor n/2\rfloor}
 \le H\theta^n\quad(n\ge0).
\end{equation}
Every source has an original $q-1$ label list whose column mass is
$C>1$. There is no $w\in(0,\infty)^d$ satisfying
$w^\top T_S\le w^\top$ for all original $q-1$ label lists.
Nevertheless every unrefined legal Borel feedback has
\begin{equation}\label{eq:gain-M}
 h^*(s)=M_1=\log q=\log r,
\end{equation}
and every finite Borel atom refinement has $h^*\ge\log q$.
\end{theorem}

\subsection{Complete estimate for every pair of global lists}
\begin{proof}
Write $S_t=U\setminus\{t\}$ and $T^{(t)}=T_{S_t}$.
The exact source-column masses, counted with edge multiplicities,
are
\begin{equation}\label{eq:gain-columns}
 c_i(t)=\sum_jT^{(t)}(j,i)=
 \begin{cases}
 (b+f)/d=1-\eps,&t=\kappa(i),\\
 1+(f-a)/d=C,&t\in A_i,\\
 1,&t=m(i).
 \end{cases}
\end{equation}
Indeed, omitting the large primary control leaves all small tails
and the spare; omitting a small primary tail removes $a$ edges
and adds $f$ spare edges to the primary total $d$; omitting the
spare leaves exactly the primary partition. Distinct controls landing
on the same target are still distinct counted edges.

Fix two whole-space lists $S_s,S_t$. The source-$i$ column mass
of their product is
\begin{equation}\label{eq:gain-identity}
 \sum_j(T^{(t)}T^{(s)})(j,i)
 =\sum_jc_j(t)T^{(s)}(j,i)
 =c_i(s)+\sum_j[c_j(t)-1]T^{(s)}(j,i).
\end{equation}
If $s=\kappa(i)$, then $c_i(s)=1-\eps$ and every next-column
mass is at most $C$. The product column is at most $C(1-\eps)$.

If $s\ne\kappa(i)$, the whole large primary tail is retained.
Its permutation image misses just $b$ letters. Hence, for every
next loss class $L_t$,
\begin{equation}\label{eq:gain-landing}
 \sum_{j\in L_t}T^{(s)}(j,i)
 \ge\frac{|p_{\kappa(i)}(B_i)\cap L_t|}{d}
 \ge\frac{|L_t|-b}{d}\ge\eta.
\end{equation}
Small and spare tails can only increase this mass. On $L_t$ the
next list has mass $1-\eps$, whereas elsewhere its mass is at
most $C$. Thus
\begin{align}
 \sum_jc_j(t)T^{(s)}(j,i)
 &\le Cc_i(s)-[C-(1-\eps)]
             \sum_{j\in L_t}T^{(s)}(j,i)\nonumber\\
 &\le C^2-(\eps+\gamma)\eta.
 \label{eq:gain-pair}
\end{align}
This includes both the current gain and the worst next gain, rather
than subtracting a loss while omitting the cost of spare labels.
The lists $s,t$ remain fixed over the state space throughout.

Dividing \eqref{eq:gain-margin} by $d^2$ yields
\[
 2\gamma+\gamma^2<(\eps+\gamma)\eta,
 \qquad\eps+\gamma=[d-(q-1)a]/d>0.
\]
Since $M\le d/q$ and $M>b$, $0<\eta<1$. The last inequality
therefore implies $\eps>\gamma+\gamma^2>\gamma$.
It follows that $C(1-\eps)<1$; it also gives
$C^2-(\eps+\gamma)\eta<1$ directly. Both quantities are positive:
the first by $b+f>0$, the second because
$\eps+\gamma<1$, $\eta<1$, and $C^2>1$.
Thus $0<\zeta<1$ and every list pair has induced $l^1$ norm
at most $\zeta$.

For every history, group consecutive factors into pairs. The possible
odd remaining factor has norm at most $C$, not one. Consequently
\[
 \|T^{(t_{n-1})}\cdots T^{(t_0)}\|_1
 \le C^{n\bmod2}\zeta^{\lfloor n/2\rfloor}.
\]
Equation~\eqref{eq:cyl-fullnorm} and $\|\one\|_1=d$ prove
\eqref{eq:gain-bn}, including $b_0=1\le d$.

To exclude \emph{all} common positive nonstrict left budgets,
suppose one existed and take a coordinate $i$ of smallest $w_i$.
Since $q\ge3$, choose $t\in A_i$. The already assumed
inequality for this global list would give
\[
 w_i\ge(w^\top T^{(t)})_i
 =\sum_jw_jT^{(t)}(j,i)\ge w_i c_i(t)=Cw_i>w_i,
\]
a contradiction. The contradiction holds for every positive vector $w$.

All states have a primary legal control. For each $u$ choose
$i\in L_u$. As $|B_i|=d-b>f$, a $z\in B_i\setminus F_i$
exists. Only $u$ is legal on $[iz]$. Thus every control is needed
in a subcover, giving $r=q$.
The actual primary selector is
\begin{equation}\label{eq:gain-selector}
 s_0(x)=\begin{cases}
 \kappa(x_0),&x_1\in B_{x_0},\\
 u,&x_1\in H_{x_0,u},\ u\in A_{x_0}.
 \end{cases}
\end{equation}
Its primary domains partition all states, so it is single-valued
and clopen. At every successor it again uses the same legal state
function; all infinite trajectories, including periodic and constant
states, stay safe. The capacity theorem gives the lower bound for
every Borel rule, not only this selector. With only $q$ physical
labels every unrefined rule has at most $q^m$ patterns on $m$
positions. This proves \eqref{eq:gain-M} and the refinement statement.
\end{proof}

The unavailable common budget is the \emph{original first-letter,
one-step, positive left} budget. Two-step products do have a strict
common budget. No claim excludes higher-dimensional lifts, right
budgets, arbitrary changes of reference measure, or path-dependent
Lyapunov functions. In particular no conclusion about $L_1$ follows
from \eqref{eq:gain-M}.

\subsection{A family for arbitrary control counts}
For any $q\ge3$ and integer $m\ge q+1$, take $d=qm$,
all $|L_t|=m$, $a=1$, and $f=2$. Each source may assign
its spare label arbitrarily, choose its $q-2$ distinct small-tail
letters arbitrarily, choose any two spare-tail letters, and use
arbitrary permutations. Conditions~\eqref{eq:gain-counts} hold.
The right side minus the left side of \eqref{eq:gain-margin} is
\[
 P(m)=qm^2-(q^2+q-1)m+q^2-3q+1.
\]
Here $P(q+1)=q^2-2q+2>0$ and
$P(m+1)-P(m)=2qm-q^2+1>0$ for $m\ge q+1$.
Thus the entire family satisfies the theorem, with constants
depending only on these counts. The condition is uniform over all placements of the allowed tails.
It does not assert that every member lacks all protected symbolic
subsystems; no such subsystem is a hypothesis of the proof.

\subsection{A three-control example with gain columns}
Let $D=\{0,\ldots,23\}$ and $U=\{0,1,2\}$.
Take $L_0=\{0,\ldots,7\}$, $L_1=\{8,\ldots,15\}$,
$L_2=\{16,\ldots,23\}$. For $i\in L_t$ the large label is
$t$, the small label is $t+1$, and the spare is $t+2$ modulo three.
The complete tails are specified in the following table; every
source in a row has that row's tails.
\begin{center}
\begin{tabular}{ccll}
\toprule
Source class&Large label&Small tail $H_t$&Spare tail $F_t$\\
\midrule
$L_0$&$0$&$\{0,1,2,3\}$&$\{4,5,6,7,8\}$\\
$L_1$&$1$&$\{8,9,10,11\}$&$\{0,1,2,3,4\}$\\
$L_2$&$2$&$\{8,9,10,11\}$&$\{0,1,2,3,4\}$\\
\bottomrule
\end{tabular}
\end{center}
The large tail is $D\setminus H_t$. Let $p_0$ be the identity,
$p_1=(0\ 1)$, $p_2=(16\ 17)$, fixing all other letters.
These permutations preserve each of their actually used tail sets.
Thus the exact matrix description, with subscripts modulo three, is
\[
 T_u(j,i)=\frac1{24}\big[
 \one_{L_u}(i)\one_{D\setminus H_u}(j)+
 \one_{L_{u-1}}(i)\one_{H_{u-1}}(j)+
 \one_{L_{u-2}}(i)\one_{F_{u-2}}(j)\big].
\]
It computes capacity without replacing the actual permutations by
common successor maps. Its target densities are
\begin{center}
\begin{tabular}{c|ccc|c}
\toprule
Target letter&$\rho_0$&$\rho_1$&$\rho_2$&Largest two sum\\
\midrule
$0,1,2,3$&$1/3$&$1$&$1/3$&$4/3$\\
$4$&$2/3$&$2/3$&$2/3$&$4/3$\\
$5,6,7$&$1/3$&$1/3$&$2/3$&$1$\\
$8$&$2/3$&$0$&$2/3$&$4/3$\\
$9,10,11$&$2/3$&$0$&$1/3$&$1$\\
$12,\ldots,23$&$1/3$&$1/3$&$1/3$&$2/3$\\
\bottomrule
\end{tabular}
\end{center}
In particular $\Theta_3=4/3$. Here $d=24$, $M=8$, $a=b=4$,
$f=5$, and the strict count margin is $64-49=15$.
The constants are
\[
 \eps=5/8,\quad\gamma=1/24,\quad C=25/24,\quad\eta=1/6,
\]
\[
 C(1-\eps)=25/64,\qquad
 C^2-(\eps+\gamma)\eta=187/192.
\]
The full-history estimate is therefore
\begin{equation}\label{eq:gain-example}
 b_n\le24(25/24)^{n\bmod2}(187/192)^{\lfloor n/2\rfloor}.
\end{equation}
Source columns have masses $3/8$, $25/24$, or $1$, depending
on which of their large, small, or spare labels is omitted.
The cylinders $[0,12],[8,12],[16,12]$ require controls
$0,1,2$. On $[16,0]$ controls $1,2$ are legal but have successor
first letters $1,0$. Consequently the full Borel class includes
choices genuinely changing the future trajectory.

There is no consecutive $000$ trajectory with uniquely legal
controls at every step. A unique $0$ at a source in $L_0$ has its
next first letter in $\{9,\ldots,23\}$; a unique $0$ in $L_2$
has its next first letter in $\{8,\ldots,11\}\subset L_1$.
In $L_1$, control $0$ is only spare and is never unique.
Thus from $L_0$ there can be at most one further unique $0$,
and from $L_2$ none. This excludes a direct consecutively forced
three-symbol full language. It does not exclude sparse bridges,
other factors, or every protected coding construction.

The gain--loss argument is elementary positive-matrix reasoning,
applied here through exact allowed-tail quotas and actual feedback
fibers. Shi et al. \cite[Definition~2.5 and Theorems~2.6, 2.8,
pp.~8--11 of arXiv:2004.01867v2]{Shi2021}
already allow losses to compensate growth in matrix products.
Their positive-diagonal and ordered-window assumptions do not
automatically hold here: for example $T^{(1)}(0,0)=0$ in this
witness. Our proof uses the mass quota
\eqref{eq:gain-landing} for every next global list and does not
depend on that theorem. The separation concerns a specified
common-left template; it is not a new general Lyapunov theory.


\section{Counting entropy and the closed-cover theorem}

Let $Q$ be a nonempty compact metric space and let $T:Q\to Q$ be a Borel
map. For a finite Borel partition $\alpha$ of $Q$ and a finite
$F\subset\Nzero$, put
\[
 \Patt_\alpha(F)=
 \left|\left\{(\alpha(T^k x))_{k\in F}:x\in Q\right\}\right|,
\]
where $\alpha(y)$ denotes the atom containing $y$ and empty atoms are
discarded. Thus $\Patt_\alpha(F)$ is the number of nonempty atoms in
$\bigvee_{k\in F}T^{-k}\alpha$. For a strictly increasing sequence
$A=(a_i)_{i\ge1}$ in $\Nzero$, define
\begin{align}
 h_A(T,\alpha)&=\limsup_{n\to\infty}\frac1n
       \log\Patt_\alpha(\{a_1,\ldots,a_n\}),\label{cc-eq:seq}\\
 p_\alpha^*(n)&=\max_{\substack{F\subset\Nzero\\ |F|=n}}
                    \Patt_\alpha(F),
 &\hpat(T,\alpha)&=\lim_{n\to\infty}\frac1n\log p_\alpha^*(n).
 \label{cc-eq:pat}
\end{align}
All logarithms are natural. These are combinatorial counting quantities;
for Borel atoms they are not Shannon entropies of an invariant measure.
The maximum in \eqref{cc-eq:pat} exists because the counts are integers
between $1$ and $|\alpha|^n$. Splitting a window into two parts gives
$p_\alpha^*(n+m)\le p_\alpha^*(n)p_\alpha^*(m)$. Subadditivity proves the
existence of the limit, and $h_A(T,\alpha)\le\hpat(T,\alpha)$.
No invertibility or continuity of the observation is used here.

Let $\K=(K_w)_{w\in W}$ be a finite closed cover of $Q$, and write
\[
 r=N(\K)=\min\{|I|:I\subset W,\ \bigcup_{w\in I}K_w=Q\}.
\]
The set $\mathcal B(\K)$ consists of all finite Borel partitions refining
$\K$: each atom must be contained in at least one $K_w$. Equivalently,
one may first consider all Borel maps $s:Q\to W$ with $x\in K_{s(x)}$,
and then allow arbitrary finite Borel refinements of their fibers.

Let $H$ be a primitive finite $0$--$1$ matrix, with vertex set $V$, and
let
\[
 Y_H=\{y\in V^{\Nzero}:H_{y_jy_{j+1}}=1\text{ for all }j\ge0\}.
\]
Primitivity means that all sufficiently large powers of $H$ have strictly
positive entries. Write $\sigma$ for the one-sided shift.

\begin{theorem}\label{cc-thm:cover}
Suppose $T:Q\to Q$ is topologically conjugate to $\sigma^p:Y_H\to Y_H$
for some primitive $H$ and integer $p\ge1$. Let $\K$ be any finite closed
cover of $Q$. For every strictly increasing $A=(a_i)$ satisfying
\begin{equation}\label{cc-eq:gaps}
 a_{i+1}-a_i\longrightarrow\infty,
\end{equation}
one has
\begin{equation}\label{cc-eq:coverformula}
 \min_{\alpha\in\mathcal B(\K)} h_A(T,\alpha)
 =\min_{\alpha\in\mathcal B(\K)}\hpat(T,\alpha)
 =\log r.
\end{equation}
An ordered difference partition of any minimum subcover attains both
minima, simultaneously for every $A$ satisfying \eqref{cc-eq:gaps}.
Consequently,
\begin{equation}\label{cc-eq:coverminimax}
 \max_A\inf_{\alpha\in\mathcal B(\K)}h_A(T,\alpha)=\log r,
\end{equation}
where the maximum ranges over all strictly increasing infinite sequences
and is attained, for example, by $a_i=(i-1)^2$.
\end{theorem}

The order of quantifiers is part of the assertion. Every $A$ satisfying
\eqref{cc-eq:gaps} works for every allowed partition, although the finite
initial segment discarded in the proof can depend on that partition.
The proof does not interchange an infimum over partitions with a
finite-window limit.

\subsection{A category and concatenation lemma}

We use meagerness relative to $Q$. A Borel set has the Baire property,
so a nonmeager Borel set is comeager in some nonempty open set. Indeed,
the sets differing from open sets by meager sets form a sigma-algebra:
for complements, the additional difference is contained in the
nowhere dense boundary of an open set. A compact metric space is a
Baire space. In particular, a nonempty open set cannot be covered by
countably many nowhere dense sets, as is also seen by choosing nested
nonempty open sets with closures avoiding those sets and diameters
tending to zero.

If $F:Q\to Q$ is continuous and open, preimages of meager sets are
meager. For a closed nowhere dense $N$, the closed set $F^{-1}N$ has
empty interior: otherwise the image of a nonempty open subset would be
a nonempty open subset of $N$. Passing to countable unions proves the
assertion.

\begin{lemma}\label{cc-lem:glue}
Under the dynamical assumptions of Theorem~\ref{cc-thm:cover}, $T$ is
continuous, open, and onto. Given finitely many nonempty cylinder sets
$C_1,\ldots,C_q$ in symbolic coordinates, there is an integer $b\ge1$
such that, for every $k\ge1$, every $t_1,\ldots,t_k\in\{1,\ldots,q\}$,
and all $0\le n_1<\cdots<n_k$ with $n_{j+1}-n_j\ge b$,
\[
 \bigcap_{j=1}^k T^{-n_j}C_{t_j}
\]
is nonempty and open.
\end{lemma}
\begin{proof}
Work on $Y_H$ through the conjugacy. The image under $\sigma$ of a
prefix cylinder of length at least two is the cylinder obtained by
removing its first symbol. The image of a one-symbol cylinder is the
union of the permitted successor cylinders. Hence $\sigma$ is open.
Every vertex has a predecessor, so every infinite path has a predecessor;
thus $\sigma$ is onto. These properties pass to powers and conjugates.

Choose within each $C_j$ a nonempty prefix cylinder, and extend their
defining words, if necessary, to a common length $\ell$. Choose $N$ with
$H^n>0$ for every $n\ge N$, and then choose $b$ with
$pb-\ell+1\ge N$. Place the word for $C_{t_j}$ at coordinate $pn_j$.
Between the end of one placed word and the start of the next there are
$p(n_{j+1}-n_j)-\ell+1$ edges to fill. Primitivity supplies a path of
that length between any two endpoint vertices. Predecessors fill the
finite segment before the first placed word, and successors provide an
infinite continuation after the last. This constructs an actual point
in the intersection. The intersection is open because it is a finite
intersection of preimages of open sets under continuous maps.
\end{proof}

\begin{lemma}\label{cc-lem:borel}
Under the same assumptions, let $\alpha$ be a finite Borel partition
with $q$ nonmeager atoms. For every $A$ satisfying \eqref{cc-eq:gaps} there
exists $m\ge1$ such that
\begin{equation}\label{cc-eq:categorycount}
 \Patt_\alpha(\{a_1,\ldots,a_n\})\ge q^{n-m+1}\qquad(n\ge m).
\end{equation}
In particular $h_A(T,\alpha)\ge\log q$.
\end{lemma}
\begin{proof}
List the nonmeager atoms as $D_1,\ldots,D_q$. The Baire property
provides a nonempty cylinder $C_j$ for which $C_j\setminus D_j$ is
meager. Let $b$ be supplied by Lemma~\ref{cc-lem:glue}, and choose $m$ so
that $a_{i+1}-a_i\ge b$ for all $i\ge m$.

Fix $n\ge m$ and any target labels $t_m,\ldots,t_n$. The set
\[
 G=\bigcap_{i=m}^n T^{-a_i}C_{t_i}
\]
is nonempty and open. Each power $T^{a_i}$ is continuous and open, so
\[
 \bigcup_{i=m}^n T^{-a_i}(C_{t_i}\setminus D_{t_i})
\]
is meager and cannot cover $G$. A point outside this union realizes
all the specified atoms at the sampled times. Distinct target label
strings give distinct patterns; restriction from the full window to
this tail is onto. This proves \eqref{cc-eq:categorycount}. Dividing its
logarithm by $n$ and taking a limit inferior proves the asserted bound.
\end{proof}

\subsection{Proof of the exact formula}

\begin{proof}[Proof of Theorem~\ref{cc-thm:cover}]
Let $\alpha\in\mathcal B(\K)$. Assign to each atom $D$ one index
$w(D)$ with $D\subset K_{w(D)}$. Let $J$ be the set of indices assigned
to the nonmeager atoms. If the closed sets $(K_w)_{w\in J}$ did not
cover $Q$, their complement would be a nonempty open set covered by
the finitely many meager atoms of $\alpha$, which is impossible.
Thus $|J|\ge r$, and the number $q$ of nonmeager atoms is at least $r$.
Lemma~\ref{cc-lem:borel} gives
\[
 h_A(T,\alpha)\ge\log q\ge\log r
\]
for every $A$ satisfying \eqref{cc-eq:gaps}. Closedness of the cover
members is used precisely in this passage from category to a full cover.

Choose a minimum subcover $K_{w_1},\ldots,K_{w_r}$, and set
\begin{equation}\label{cc-eq:ordered}
 E_j=K_{w_j}\setminus\bigcup_{k<j}K_{w_k},\qquad
 \alpha_0=\{E_1,\ldots,E_r\}.
\end{equation}
Every $E_j$ is Borel and is nonempty, since a minimum subcover has no
redundant member. The partition refines $\K$ and has $r$ atoms. Hence
$\Patt_{\alpha_0}(F)\le r^{|F|}$ for all finite $F$, giving
$h_A(T,\alpha_0)\le\hpat(T,\alpha_0)\le\log r$ for every $A$.
The established lower bound and $\hpat\ge h_A$ yield
\eqref{cc-eq:coverformula} and simultaneous attainment.
For arbitrary $A$, the partition $\alpha_0$ bounds the inner infimum
in \eqref{cc-eq:coverminimax} by $\log r$. Any $A$ satisfying
\eqref{cc-eq:gaps} gives the reverse bound and attains the maximum.
\end{proof}


\section{A cover--partition gap at prescribed sparse times}

For any finite cover $\K$ define its combinatorial sequence entropy by
\[
 h_A(T,\K)=\limsup_{n\to\infty}\frac1n
 \log N\left(\bigvee_{i=1}^n T^{-a_i}\K\right).
\]
For a partition this is \eqref{cc-eq:seq}. For overlapping sets, covering
each finite orbit window need not correspond to one fixed refining
partition used at every time.

\begin{corollary}\label{cc-cor:covergap}
Let $q\ge3$, $Q=\{0,\ldots,q-1\}^{\Nzero}$, $T=\sigma$, and
\[
 K_j=\{x:x_0\ne j\}\quad(0\le j<q),\qquad \K=(K_j)_{j=0}^{q-1}.
\]
For every increasing infinite $A$,
\begin{equation}\label{cc-eq:covervalue}
 h_A(\sigma,\K)=\log\frac q{q-1}.
\end{equation}
For every such $A$ with diverging successive gaps,
\begin{equation}\label{cc-eq:refvalue}
 \min_{\alpha\in\mathcal B(\K)}h_A(\sigma,\alpha)=\log2
 >h_A(\sigma,\K).
\end{equation}
\end{corollary}
\begin{proof}
The cover has no single member equal to $Q$, and any two distinct
members cover $Q$. Thus $r=2$, and \eqref{cc-eq:refvalue} follows from
Theorem~\ref{cc-thm:cover} once \eqref{cc-eq:covervalue} is proved.

For a window $F$ of size $n$, the full shift has all $q^n$ possible
input vectors. Each member of the joined cover specifies one excluded
letter at each coordinate, and covers $(q-1)^n$ of these vectors.
Consequently the minimum number $c(n)$ of such boxes is at least
$(q/(q-1))^n$. For completeness, the following elementary probabilistic
covering bound gives the matching exponent. Choose $R$ excluded-letter
vectors independently and uniformly from the $q^n$ possibilities. A
fixed input vector is left uncovered with probability at most
\[
 \left(1-\left(\frac{q-1}{q}\right)^n\right)^R
 \le \exp\left(-R\left(\frac{q-1}{q}\right)^n\right).
\]
Taking
$R=\left\lceil(q/(q-1))^n(n\log q+1)\right\rceil$ makes the sum of
these probabilities over all input vectors smaller than one. Some
choice therefore covers every vector. We have proved
\[
 \left(\frac q{q-1}\right)^n\le c(n)
 \le\left\lceil\left(\frac q{q-1}\right)^n(n\log q+1)\right\rceil.
\]
The count depends only on $n$, not on the positions of the sampled
coordinates, so its exponential rate proves \eqref{cc-eq:covervalue}.
\end{proof}

The excluded-letter cover and its covering exponent are already present
in Romagnoli \cite{Romagnoli2007}; the calculation above is included to
identify the sparse-time quantity without an external dependency.
The consequence of Theorem~\ref{cc-thm:cover} is the exact minimum over
\emph{all Borel} refinements for a sequence fixed in advance. It does not
assert that every Borel refinement has consecutive-time entropy at least
$\log2$, or that every sampling sequence gives that lower bound.


\section{Common-successor feedback systems}\label{sec:common-successor}

\begin{theorem}\label{cc-thm:control}
In the setting of Section~\ref{sec:selectors}, assume that there is
a map $T:Q\to Q$ such that
\begin{equation}\label{cc-eq:common}
 g_w(x)=T(x)\qquad\text{for every }w\in W\text{ and every }x\in K_w,
\end{equation}
and that $T$ is conjugate to $\sigma^p$ on a mixing one-sided shift of
finite type. With $r=N((K_w)_{w\in W})$, for every $A$ satisfying
\eqref{cc-eq:gaps} one has
\begin{equation}\label{cc-eq:controlformula}
 \min_{s\in\Ssafe_\tau}h_A(s)
 =\min_{s\in\Ssafe_\tau}\hpat(s)
 =L_\tau=M_\tau=\frac{\log r}{\tau}.
\end{equation}
The maximum defining $L_\tau$ is attained by every such $A$, and an
ordered minimum-subcover selector attains all the indicated minima.
The same minimum and minimax values hold when all finite Borel safe
partitions, including every finite Borel refinement, are allowed. Every
refinement satisfies the stated lower bound; no upper bound for each
individual refinement is asserted.
\end{theorem}
\begin{proof}
For every legal selector, \eqref{cc-eq:common} gives $T_s=T$ pointwise.
Its fibers form a finite Borel refinement of $(K_w)_{w\in W}$.
Theorem~\ref{cc-thm:cover}, with normalization by $\tau$, gives the lower
bound at every $A$ with diverging gaps. If $w_1,\ldots,w_r$ form a
minimum subcover, define $s_0(x)=w_j$ on the sets $E_j$ of
\eqref{cc-eq:ordered}. This is Borel and legal at every point. Each selected
block stays in $Q$ at its intermediate times by \eqref{cc-eq:safe}, and
the endpoint again lies in $Q$, so recursion gives an infinite safe
closed-loop orbit for every initial state. The selector uses only $r$
blocks and thus gives the upper bound for every window and every
sampling sequence. This proves attainment, \eqref{cc-eq:controlformula},
and the outer maximum. Lemma~\ref{cc-lem:merge} proves the last assertion.
\end{proof}

The standard safe-partition framework is described by Kawan
\cite[Chapter~2]{Kawan2013}. His minimization over control periods in
invariance entropy is a different optimization. Here $\tau$ is fixed,
intermediate safety is retained, and the exact value comes from the
additional common-successor and finite-type hypotheses.

\subsection{A positive overlapping example at period two}

Take
\[
 Q=\{0,1,2\}^{\Nzero},\quad X=Q\sqcup\{\dagger\},\quad
 U=\{0,1,2\},\quad\tau=2,
\]
where the isolated point $\dagger$ lies outside $Q$. Define
\[
 f_u(x)=
 \begin{cases}\sigma x,&x\in Q,\ x_0\ne u,\\
 \dagger,&x\in Q,\ x_0=u,\end{cases}
 \qquad f_u(\dagger)=\dagger.
\]
Each $f_u$ is continuous since the cases are clopen. For $w=(a,b)$,
\[
 K_{(a,b)}=\{x:x_0\ne a,\ x_1\ne b\},\qquad
 g_{(a,b)}|_{K_{(a,b)}}=\sigma^2.
\]
All nine control blocks are available to Borel feedbacks, and each
state permits four of them. A safety domain covers four of the nine
two-symbol input blocks, so two domains cannot cover $Q$. The three
diagonal blocks $(0,0),(1,1),(2,2)$ do cover $Q$, since some symbol
differs from both $x_0$ and $x_1$. Hence $r=3$, and
\[
 L_2=M_2=\frac{\log3}{2}.
\]
An attaining selector is
\[
 c(x)=\min(\{0,1,2\}\setminus\{x_0,x_1\}),\qquad
 s_0(x)=(c(x),c(x)).
\]
It is clopen and legal everywhere. Its actual block-boundary states
are $\sigma^{2k}x$. To prescribe any sequence of diagonal labels, use
successive input blocks $11$, $00$, and $01$ for the desired labels
$0$, $1$, and $2$, respectively. Therefore its actual name set is the
full shift on these three blocks. The lower bound, however, holds for
every legal Borel rule using any of the nine blocks.

This example does not preassign a nonconstant common control-letter
observation to every feedback. More precisely, suppose a map
$\rho:U^2\to V'$ and an observation $F:Q\to V'$ satisfy
$\rho(s(x))=F(x)$ for every legal Borel selector and every $x$.
Any two blocks $(a,b)$ and $(a',b')$ are colegal at a state beginning
with $ij$, where $i\notin\{a,a'\}$ and $j\notin\{b,b'\}$.
A legal baseline selector may be changed at this single state to use
either block, while remaining Borel and legal. Hence their $\rho$
values agree. All pairs of blocks are colegal somewhere, so $\rho$
is constant. This rules out this particular forced-letter explanation;
it makes no claim about every possible dynamical factor.


\section{Failure of the cover formula with different successors}

The common-successor assumption cannot be removed solely because each
branch is continuous and open. We give a system with individual
prefix cylinders as safety domains and compute its full-class value.

Let $Q=\{0,1\}^{\Nzero}$, let $\mathcal C=\{0,10,11\}$ be a set of
finite words, and take $U=\mathcal C\times\{0,1\}$ and $\tau=1$.
For $e\in\{0,1\}$ let $P_e$ add $e$ modulo two to every coordinate.
On $X=Q\sqcup\{\dagger\}$ define
\begin{equation}\label{cc-eq:boundarysystem}
 f_{(c,e)}(x)=
 \begin{cases}P_e\sigma x,&x\in[c],\\
 \dagger,&x\in Q\setminus[c],\end{cases}
 \qquad f_{(c,e)}(\dagger)=\dagger.
\end{equation}
The cylinders $[0],[10],[11]$ are disjoint and cover $Q$. Each state
has exactly two legal controls with different successors. The global
maps $f_{(c,e)}$ are continuous, and their legal restrictions are open.
The safety domains are $K_{(c,e)}=[c]$, and $r=3$.

\begin{proposition}\label{cc-prop:boundary}
For the system \eqref{cc-eq:boundarysystem}, over all legal Borel feedbacks
and arbitrary Borel mixtures,
\begin{equation}\label{cc-eq:boundaryvalue}
 \min_s h_{A_0}(s)=\min_s\hpat(s)=
 \max_A\inf_s h_A(s)=\log2<\log r,
 \qquad A_0=(0,1,2,\ldots).
\end{equation}
There is one clopen legal feedback $s_{\mathrm f}$ satisfying
$h_A(s_{\mathrm f})=\log2$ for every increasing infinite $A$.
The same full-class values hold if arbitrary finite Borel atom
refinements are allowed.
\end{proposition}
\begin{proof}
Let $c(x)$ be the unique member of $\mathcal C$ with $x\in[c(x)]$,
and set
\[
 s_{\mathrm f}(x)=(c(x),x_1),\qquad
 T_{\mathrm f}x=P_{x_1}\sigma x.
\]
This is a clopen selector, legal at every state. Its four nonempty
initial control fibers, for labels $(0,0),(0,1),(10,0),(11,1)$, are
$[00],[01],[10],[11]$, respectively. For all $n\ge1$, direct induction
gives
\begin{equation}\label{cc-eq:fliporbit}
 T_{\mathrm f}^n x=P_{x_n}\sigma^n x,
 \qquad s_{\mathrm f}(T_{\mathrm f}^n x)
        =(0,x_{n+1}\mathbin\oplus x_n),
\end{equation}
where $\oplus$ is addition modulo two. In the induction step the
chosen flip is $x_{n+1}\oplus x_n$; it cancels the previous flip and
leaves $P_{x_{n+1}}\sigma^{n+1}x$. The first coordinate is always zero
after time zero, which also proves all-time safety.

For every finite nonempty window $F$, with $n=|F|$, one has
\begin{equation}\label{cc-eq:flipcount}
 \Patt_{s_{\mathrm f}}(F)=
 \begin{cases}2^n,&0\notin F,\\2^{n+1},&0\in F.\end{cases}
\end{equation}
The four initial labels and two later labels give the upper bounds.
To prove the lower bounds, prescribe any of the four initial labels
if $0\in F$, thereby fixing $(x_0,x_1)$; if $0\notin F$, set
$x_0=x_1=0$. At positive sampled times prescribe the desired second
components $d_k$, and put $d_k=0$ at the other positive times.
Recursively set $x_{k+1}=x_k\oplus d_k$ for every $k\ge1$.
This constructs a full state in $Q$ whose actual trajectory under the
same selector has the required sampled controls by
\eqref{cc-eq:fliporbit}. Thus all counted patterns are realized by true
infinite closed-loop trajectories. Equation~\eqref{cc-eq:flipcount}
implies $h_A(s_{\mathrm f})=\hpat(s_{\mathrm f})=\log2$ for every $A$.

For the opposite bound fix any legal Borel selector $s$. If its first
$n$ actual controls from $x$ are $(c_k,e_k)$, let
$E_0=0$ and $E_k=e_0\oplus\cdots\oplus e_{k-1}$. Since shifts commute
with flips,
\[
 T_s^k x=P_{E_k}\sigma^k x.
\]
Let $b(c)$ be the first bit of the word $c$. Safety at time $k$ gives
\begin{equation}\label{cc-eq:decode}
 x_k=b(c_k)\oplus E_k\qquad(0\le k<n).
\end{equation}
Thus one actual length-$n$ control word determines at most one input
prefix of length $n$. Every binary input prefix has an extension in
$Q$, and the feedback is defined and safe on every such extension.
Distinct prefixes must yield distinct actual control words. Hence
\[
 \Patt_s(\{0,\ldots,n-1\})\ge2^n
\]
for every $n$ and every legal $s$. This proves $h_{A_0}(s)\ge\log2$
and $\hpat(s)\ge\log2$. Combining these inequalities with
$s_{\mathrm f}$ proves \eqref{cc-eq:boundaryvalue}, including attainment.
Lemma~\ref{cc-lem:merge} handles atom refinements.
\end{proof}

The obstruction has a precise orbit interpretation. In this example
$T_{\mathrm f}(Q)=[0]$: for the reverse inclusion, given $y\in[0]$,
take $x=0y$. Therefore neither $[10]$ nor $[11]$ is visited at a
positive time. The initial fibers with those labels are nonmeager,
but their labels cannot be freely concatenated in the future.
The map $T_{\mathrm f}$ is itself continuous and open, as follows by
restricting to the four clopen two-symbol cylinders. It is not onto
$Q$.

Even preservation of meager preimages does not repair this problem.
For any selector $s$ and any $n$, the value of $T_s^n x$ is either
$\sigma^n x$ or $P_1\sigma^n x$. Hence, for a meager set $N$,
\[
 (T_s^n)^{-1}N\subset
 (\sigma^n)^{-1}N\ \cup\ (P_1\sigma^n)^{-1}N
\]
is meager. What fails is the nonempty open intersection used in
Lemma~\ref{cc-lem:glue} for the \emph{selected} dynamics. Also, the
same $s_{\mathrm f}$ gives
\[
 \inf_s h_A(s)\le\log2=\log r-\log(3/2)
 \quad\text{for every }A.
\]
The failure of the cover formula is therefore a uniform numerical
gap over all samplings, rather than a failure at one finite window.
It does not refute $L_\tau=M_\tau$: equality holds in this very example.
Nor does the proof determine the inner infimum for every other fixed
sampling sequence.




\subsection{A cyclic-permutation application}
Let $q\ge3$, $G=\mathbb Z/q\mathbb Z$, and $Q=G^{\Nzero}$.
Write $P_u$ for coordinatewise addition of $u$. On
$X=Q\sqcup\{\dagger\}$ use the continuous controls
\[
 f_u(x)=P_u\sigma x\quad\text{if }x_0\ne u,
 \qquad f_u(x)=\dagger\quad\text{otherwise},
 \qquad f_u(\dagger)=\dagger.
\]
Their safety domains are the excluded-letter cylinders
$K_u=\{x:x_0\ne u\}$, so the minimum cover cardinality is two.

\begin{corollary}\label{cor:cyclic}
For this system, at period one and over all legal Borel feedbacks
and finite atom refinements,
\[
 M_1=\min_s h^*(s)=\log2,\qquad
 \log\frac q{q-1}\le L_1\le\log2,\qquad
 \hinv(Q)=\log\frac q{q-1}.
\]
There is a clopen feedback with $\Patt_s(F)=2^{|F|}$ for every
finite window $F$.
\end{corollary}
\begin{proof}
Use the uniform Bernoulli probability on $Q$. The safe inverse
operator is
\[
 (L_uf)(y)=\frac1q\sum_{a\ne u}f(aP_{-u}y),
 \qquad L_u\one=\frac{q-1}{q}\one.
\]
Products of any single-label lists therefore have norm
$((q-1)/q)^n$. Theorem~\ref{thm:capacity-main}, with $k=r=2$,
gives $M_1=\log2$. This is a direct application of the product
criterion, rather than an additional main theorem.

For a real control word $w=(w_0,\ldots,w_{n-1})$, put
$E_k=\sum_{j<k}w_j$ and $v_k=w_k-E_k$, with arithmetic in $G$.
The actual orbit satisfies $T_s^kx=P_{E_k}\sigma^kx$, so its
initial fiber lies in
\[
 K(w)=\{x:x_k\ne v_k\ (0\le k<n)\}.
\]
This is also the complete open-loop safety domain of $w$.
Each such domain contains exactly $(q-1)^n$ input prefixes.
Every input prefix extends to a state with an infinite safe
feedback trajectory; hence the actual length-$n$ word count is
at least $(q/(q-1))^n$. Consecutive sampling gives the stated
lower bound for $L_1$.

Define $s_0(x)=1$ if $x_0=0$, and $s_0(x)=0$ otherwise.
This is clopen and safe on all of $Q$. Given any infinite binary
sequence $d$, let $E_k=\sum_{j<k}d_j$ and let $b_k=0$ when
$d_k=1$, $b_k=1$ when $d_k=0$. The state with coordinates
$x_k=b_k-E_k$ has current first coordinate $b_k$ at time $k$,
by induction. Its real feedback name is therefore $d$.
This proves the finite-window equality and the upper bound for
$L_1$ with the same feedback for every sampling sequence.

The map $w\mapsto v$ is triangular and bijective: recover
$w_k=v_k+E_k$, $E_{k+1}=2E_k+v_k$, starting with $E_0=0$.
This requires no inverse of $2$, and includes even $q$.
Thus the open-loop spanning number is the least number $c_q(n)$
of boxes $\prod_{k<n}(G\setminus\{v_k\})$ covering $G^n$.
Writing $p=((q-1)/q)^n$, the counting lower bound and the
elementary random-cover upper bound give
\[
 p^{-1}\le c_q(n)\le
 \left\lceil p^{-1}(n\log q+1)\right\rceil.
\]
Indeed, $t=\lceil p^{-1}(n\log q+1)\rceil$ independent uniform
boxes leave some input uncovered with probability at most
$q^n(1-p)^t\le q^ne^{-pt}\le e^{-1}<1$.
A covering choice exists, and removal of repetitions can only
decrease its size. Taking normalized logarithms proves the
invariance-entropy formula. Lemma~\ref{cc-lem:merge} supplies
the refinement statement.
\end{proof}

The safe spanning number is a lower bound for the actual feedback
word count, not an upper bound. For example, when $q=3$, the
three safe words $00,12,21$ cover all two-symbol inputs: their
excluded vectors are $00,11,22$. The clopen rule
$s(x)=\min(G\setminus\{x_0,x_1\})$ repeatedly selects the
first control of a covering word and is safe everywhere.
Starting from $x=(1,1,0,0,\ldots)$, however, it selects $0$
and then $2$, giving the real word $02$ outside that cover.
The subsequent state is $2^\infty$ and all subsequent controls
are $0$, so the witness has an infinite safe continuation.
In general its real prefix fibers obey
\[
 H_s(uw)=s^{-1}(u)\cap(P_u\sigma)^{-1}H_s(w),
 \qquad H_s(w)\subset K(w).
\]
These compatibility equations are absent from a separately
chosen spanning family. The lower and upper bounds for $L_1$
remain distinct; neither their equality nor a strict minimax
gap is asserted for this cyclic system.


\section{Simultaneous attainment for pattern systems}\label{sec:comb}

A \emph{finite-alphabet pattern system} is a nonempty set
$\Omega\subset E^{\Nzero}$ with $E$ finite and nonempty.  No topology or
shift invariance is imposed.  For finite $F\subset\Nzero$, write
\begin{align*}
 \Patt_\Omega(F)=\#\{x|_F:x\in\Omega\},
 \qquad
 p_\Omega^*(n)=
 \max_{\substack{F\subset\Nzero\\\#F=n}}\Patt_\Omega(F).
\end{align*}
With a fixed scale $\tau>0$, define
\begin{align*}
 \hpat(\Omega;\tau)&=\lim_{n\to\infty}
       \frac1{n\tau}\log p_\Omega^*(n),\\
 \hseq^A(\Omega;\tau)&=\limsup_{n\to\infty}
       \frac1{n\tau}\log
       \Patt_\Omega(\{a_1,\ldots,a_n\}).
\end{align*}
The maximum exists: its possible values form a nonempty subset of the finite
integer set $\{1,\ldots,(\#E)^n\}$.

We use $\Patt_\Omega(\varnothing)=1$ when needed.
Taking the product-topology closure of $\Omega$ changes none of its
finite projections: a cylinder meeting $\overline\Omega$ already meets
$\Omega$. Thus dispensing with closedness is only a matter of
formulation. The absence of shift invariance, in contrast, matters to the
placement of the coordinate sets.

\begin{lemma}[Pattern submultiplicativity and deletion]\label{lem:pattern-submult}
Let $\Omega\subset E^{\Nzero}$ be a pattern system and $q=\#E$.
For $m,n\in\N$ and finite $B\subset F\subset\Nzero$,
\begin{align}
 p_\Omega^*(m+n)&\le p_\Omega^*(m)p_\Omega^*(n),
 \label{eq:pattern-submult}\\
 \Patt_\Omega(F)&\le q^{\#(F\setminus B)}\Patt_\Omega(B).
 \label{eq:deletion} 
\end{align}
If $B\subset G$, then $\Patt_\Omega(B)\le\Patt_\Omega(G)$.
Consequently,
\begin{equation}\label{eq:fekete-pstar}
 \hpat(\Omega;\tau)
 =\inf_{n\ge1}\frac1{n\tau}\log p_\Omega^*(n),
 \qquad
 \log p_\Omega^*(n)\ge n\tau\hpat(\Omega;\tau).
\end{equation}
\end{lemma}

\begin{proof}
Split any $(m+n)$-element set into disjoint subsets of sizes $m$ and $n$.
Restriction sends a pattern on the union injectively to the pair of its two
restrictions, proving \eqref{eq:pattern-submult}.  Restriction from $F$ to
$B$ has fibers of cardinality at most $q^{\#(F\setminus B)}$, which proves
\eqref{eq:deletion}.  If $B\subset G$, every pattern on $B$ extends to a
pattern on $G$ using the same element of $\Omega$, giving monotonicity.
Fekete's lemma applied to $\log p_\Omega^*(n)$ proves
\eqref{eq:fekete-pstar}.
\end{proof}

For comparison, let $\Omega=\{0^\infty,10^\infty\}$. This set is a
closed one-sided subshift, but
$\Patt_\Omega(\{0\})=2$ and $\Patt_\Omega(\{m\})=1$ for $m\ge1$.
One cannot simply translate a maximizing set. Deleting a finite initial
interval has a different advantage: the number of lost coordinates is
bounded independently of the size of the maximizing set. Inequality
\eqref{eq:deletion} turns this bound into a controlled information loss.

\begin{theorem}[Simultaneous attainment]\label{thm:A}
Fix $\tau>0$.  For every countable family
$(\Omega_j)_{j\ge1}$ of finite-alphabet pattern systems, there is one
$A\in\mathscr S$ such that
\begin{equation}\label{eq:simultaneous-attainment}
 \hseq^A(\Omega_j;\tau)=\hpat(\Omega_j;\tau)
 \qquad(j\ge1).
\end{equation}
Thus the maximal-pattern entropy rate of each system is attained along the same
sequence.  No closedness, compactness, shift invariance, shift surjectivity, or
backward extension is required.
\end{theorem}

\begin{proof}
Let $E_j$ be the alphabet of $\Omega_j$, put $q_j=\#E_j$ and
$h_j^*=\hpat(\Omega_j;\tau)$.  For every $A$ and $n$,
\begin{align*}
 \Patt_{\Omega_j}(\{a_1,\ldots,a_n\})\le p_{\Omega_j}^*(n),
\end{align*}
so \eqref{eq:fekete-pstar} gives
$\hseq^A(\Omega_j;\tau)\le h_j^*$.

Choose an enumeration of stages
\begin{align*}
 (j_k)_{k\ge1}=(1;\ 1,2;\ 1,2,3;\ 1,2,3,4;\ldots),
\end{align*}
where semicolons only separate finite groups. Every positive integer
occurs infinitely often. Choose $\eta_k>0$ with $\eta_k\to0$.  We recursively
construct nonempty finite sets $B_k$, each lying strictly to the right of all
previous ones.  Put
\begin{align*}
 G_{k-1}=\bigcup_{i<k}B_i,
 \qquad r_{k-1}=\#G_{k-1},
 \qquad M_{k-1}=\max G_{k-1},
\end{align*}
with $r_0=0$ and $M_0=-1$.

If $h_{j_k}^*>0$, choose $n_k$ so large that
\begin{equation}\label{eq:explicit-losses}
 n_k>M_{k-1}+1,
 \qquad
 \frac{(M_{k-1}+1)\log q_{j_k}}{n_k\tau}<\eta_k,
 \qquad
 \frac{r_{k-1}}{n_k}<\eta_k.
\end{equation}
The quantities on the numerators are fixed before $n_k$ is chosen,
so these requirements can be satisfied even if the previously selected
coordinates are extremely far apart. Notice that $M_{k-1}+1$ bounds the
number of \emph{possible} old positions, whereas $r_{k-1}$ counts the
positions actually selected. The former controls deletion; the latter
controls dilution by historical observations. No uniform bound on $q_j$
or on the coordinate gaps is required. Choose an
$n_k$-element $F_k$ attaining $p_{\Omega_{j_k}}^*(n_k)$ and set
\begin{align*}
 B_k=F_k\cap\{M_{k-1}+1,M_{k-1}+2,\ldots\}.
\end{align*}
At most $M_{k-1}+1$ coordinates are deleted, and the first condition in
\eqref{eq:explicit-losses} makes $B_k$ nonempty.  If $h_{j_k}^*=0$, set
$B_k=\{M_{k-1}+1\}$.

Let $A$ be the increasing enumeration of $\bigcup_{k\ge1}B_k$.  Fix $j$
with $h_j^*>0$ and restrict to the infinitely many stages with $j_k=j$.
Write $b_k=\#B_k$ and $R_k=r_{k-1}+b_k=\#G_k$; then $G_k$ consists of the
first $R_k$ entries of $A$.  By deletion, \eqref{eq:fekete-pstar}, and
monotonicity,
\begin{align}
 \log\Patt_{\Omega_j}(G_k)
 &\ge \log\Patt_{\Omega_j}(B_k)\notag\\
 &\ge \log p_{\Omega_j}^*(n_k)-(M_{k-1}+1)\log q_j\notag\\
 &\ge n_k\tau h_j^*-(M_{k-1}+1)\log q_j.
 \label{eq:endpoint-numerator}
\end{align}
Also $R_k\le r_{k-1}+n_k$.  At all sufficiently late such stages the last
line of \eqref{eq:endpoint-numerator} is positive.  Dividing by $R_k\tau$
and using \eqref{eq:explicit-losses} yields the explicit endpoint estimate
\begin{equation}\label{eq:endpoint-bound}
 \frac1{R_k\tau}\log\Patt_{\Omega_j}(G_k)
 >\frac{h_j^*-\eta_k}{1+\eta_k}.
\end{equation}
For a fixed $j$, the indices $k$ with $j_k=j$ tend to infinity, hence
$\eta_k\to0$ along those stages. Since every $B_i$ is nonempty,
$R_k\ge k\to\infty$. The limsup along these
endpoints is at least $h_j^*$, proving the lower
bound.  If $h_j^*=0$, the universal upper bound and nonnegativity give
equality for every $A$.  This proves \eqref{eq:simultaneous-attainment} for
all $j$.
\end{proof}

\begin{remark}[What is attained]\label{rem:attainment}
The theorem asserts equality of limsup rates. It does not assert a common
subsequence of endpoints on which every system is simultaneously close
to its maximal rate. The good endpoints used for different $j$ can be
different. Nor does it assert convergence of the normalized counts.
For example, take $\tau=1$, let $B\subset\Nzero$ and its complement both be infinite,
and let $\Omega_B$ consist of all binary sequences supported on $B$.
For any schedule, the two normalized logarithmic counts for $\Omega_B$
and $\Omega_{\Nzero\setminus B}$ add to $\log2$ at every $n$, although
both maximal-pattern rates are $\log2$. They cannot both converge to
$\log2$. The two limsups, however, can both equal $\log2$.

Exact finite-level maximization is also a stronger and generally false
requirement. Put
\begin{align*}
 \Omega=\{0x:x\in\{0,1\}^{\Nzero}\}\cup\{10^\infty,20^\infty\}
 \subset\{0,1,2\}^{\Nzero}.
\end{align*}
The unique maximizing singleton is $\{0\}$, with three patterns.
For a three-element set containing zero there are six patterns, whereas
any three positive coordinates realize eight patterns. No increasing
sequence maximizes the counts at both levels $n=1$ and $n=3$.
\end{remark}

\begin{corollary}[Fixed-language attainment]\label{cor:fixed}
For every finite-alphabet pattern system $\Omega$ and every $\tau>0$,
\begin{equation}\label{eq:fixed-attainment}
 \hpat(\Omega;\tau)=\max_{A\in\mathscr S}\hseq^A(\Omega;\tau).
\end{equation}
The same identity holds for every finite Borel invariant partition $\C$ after
putting $\Omega=\Sigma_\C$.
\end{corollary}


\section{A language-quotient minimax theorem}\label{sec:minimax}

Fix a compact controlled-invariant set $Q$ and one period $\tau\in\N$.
For two invariant partitions $\C$ and $\D$ of this period, write
$\C\sim_{\lang}\D$ if there is a bijection
$\pi:I_\C\to I_\D$ such that, for every $N\in\N$,
\begin{equation}\label{eq:language-equivalence}
 \pi^N(W_N(\C;Q))=W_N(\D;Q),
\end{equation}
where $\pi^N$ acts coordinatewise.  This is \emph{language equivalence}.
Identity, inverse, and composition of the alphabet bijections show that
$\sim_{\lang}$ is an equivalence relation.  If $F\subset\Nzero$ is finite and
$m=\max F$, then \eqref{eq:language-equivalence} at length $m+1$, followed by
projection onto $F$, gives a bijection between the two $F$-pattern sets.
Thus language-equivalent partitions have identical finite-coordinate pattern
counts and identical sequence and maximal-pattern complexities.

For a nonempty class $\mathfrak C\subset\mathfrak C_{\Bor}(Q,\tau)$, set
\begin{align}
 H_{\rm pat}^*(Q;\mathfrak C)
 &=\inf_{\C\in\mathfrak C}\hpat(\C;Q),
 \label{eq:system-pattern}\\
 H_{\rm seq}^A(Q;\mathfrak C)
 &=\inf_{\C\in\mathfrak C}\hseq^A(\C;Q).
 \label{eq:system-sequence}
\end{align}
Every individual rate is nonnegative and finite, and because the class is nonempty both
infima are finite; no uniform bound on the numbers of atoms is needed.

The elementary minimax inequality gives only
\begin{align*}
 \sup_A\inf_{\C\in\mathfrak C}\hseq^A(\C;Q)
 \le\inf_{\C\in\mathfrak C}\sup_A\hseq^A(\C;Q)
 =H_{\rm pat}^*(Q;\mathfrak C).
\end{align*}
Separate maximizing schedules for separate partitions do not give the
opposite inequality. For instance, a two-by-two array with zero diagonal
and unit off-diagonal has maximum row minimum zero and minimum column
maximum one. The next theorem avoids this logical obstruction by
attaining every individual rate at one schedule.

\begin{theorem}[Minimax over countably many language types]\label{thm:B}
Let $\varnothing\ne\mathfrak C\subset\mathfrak C_{\Bor}(Q,\tau)$ and suppose
that the quotient $\mathfrak C/{\sim_{\lang}}$ is finite or countable.  There
is one $A\in\mathscr S$ such that
\begin{equation}\label{eq:pointwise-class-attainment}
 \hseq^A(\C;Q)=\hpat(\C;Q)
 \qquad(\C\in\mathfrak C).
\end{equation}
Consequently,
\begin{equation}\label{eq:minimax}
 \max_{B\in\mathscr S}\inf_{\C\in\mathfrak C}
 \hseq^B(\C;Q)
 =\inf_{\C\in\mathfrak C}\hpat(\C;Q).
\end{equation}
The outer maximum is attained. The pattern infimum and the inner infimum at the common attaining schedule are minima, by Corollary~\ref{cor:pattern-minimum}.
\end{theorem}

\begin{proof}
Choose one representative $\C_j$ of each language-equivalence class.
Enumerate the representatives, repeating them periodically when the quotient
is finite, and apply Theorem~\ref{thm:A} to the resulting countable family
$\Sigma_{\C_j}$, using the common scale $\tau$.  The resulting $A$ gives equality for each
representative.  Language equivalence preserves all finite pattern counts, so
it gives \eqref{eq:pointwise-class-attainment} for every $\C\in\mathfrak C$.
Taking the infimum over $\C$ gives equality at this $A$.

For any $B\in\mathscr S$, Corollary~\ref{cor:fixed} gives
$\hseq^B(\C;Q)\le\hpat(\C;Q)$ for every $\C$.  Taking first the infimum over
$\C$ and then the supremum over $B$ proves the reverse minimax inequality and
hence \eqref{eq:minimax}.  If the pattern infimum is zero, the pointwise upper
bound and nonnegativity in fact give
$\inf_{\C\in\mathfrak C}\hseq^B(\C;Q)=0$ for every $B\in\mathscr S$, without
any countability assumption.
\end{proof}

The hypothesis concerns the symbolic behavior generated on $Q$, not the
names of controls used to present a strategy.  It identifies alphabet renumberings and strategies with the same
admissible language after such a renumbering. It does not merge distinct
atoms merely because their assigned controls agree.  It is therefore closer to an intrinsic
statement about the chosen partition class, although it remains relative to
that class and to the fixed period.

\begin{proposition}[Countability of clopen feedbacks]
\label{prop:digital}
Assume that $Q$ is compact metrizable and that the control alphabet $U$ is
countable.  Then
$\mathfrak C_{\clop}(Q,\tau)$ is countable.  Hence, if this class is nonempty,
Theorem~\ref{thm:B} applies to it.  In particular this holds when $U$ is a
finite digital alphabet.
\end{proposition}

\begin{proof}
A compact metrizable space has a countable base.  Every clopen subset is
compact and hence is a finite union of basic open sets contained in it.  Thus
$Q$ has only countably many clopen subsets and only countably many finite
clopen partitions.  Since $\tau$ is finite, $U^\tau$ is countable.  A finite
partition admits only countably many block assignments.  The invariant assignments form a
subcollection.
\end{proof}

Theorem~\ref{thm:B} fixes $\tau$.  It does not produce one common
physical-time schedule for partitions of different periods and proves no
minimax identity over $\mathfrak C_{\Bor}(Q)$.

\begin{example}[A purely combinatorial uncountable obstruction]
\label{ex:uncountable}
For every infinite $B\subset\Nzero$, let
\begin{align*}
 \Omega_B=\{x\in\{0,1\}^{\Nzero}:\supp(x)\subset B\}.
\end{align*}
At scale $\tau=1$,
\begin{align*}
 \Patt_{\Omega_B}(F)=2^{\#(F\cap B)},
 \qquad \hpat(\Omega_B;1)=\log2.
\end{align*}
For $A=(a_i)\in\mathscr S$,
\begin{align*}
 \hseq^A(\Omega_B;1)
 =\log2\cdot\limsup_{n\to\infty}
 \frac{\#(B\cap\{a_1,\ldots,a_n\})}{n}.
\end{align*}
If $\Nzero\setminus\{a_i:i\ge1\}$ is infinite, choose an infinite $B$ in that
complement.  Otherwise choose $B=\{a_{k^2}:k\ge1\}$.  In either case the
displayed upper density is zero.  Therefore
\begin{equation}\label{eq:abstract-obstruction}
 \sup_{A\in\mathscr S}
 \inf_{\substack{B\subset\Nzero\\B\text{ infinite}}}
 \hseq^A(\Omega_B;1)=0
 <\log2
 =\inf_{\substack{B\subset\Nzero\\B\text{ infinite}}}
 \hpat(\Omega_B;1).
\end{equation}
This proves only that minimax exchange is not a formal consequence of
pointwise attainment for arbitrary uncountable pattern families.  Every
invariant-partition name system satisfies
$\sigma(\Sigma_\C)\subset\Sigma_\C$ by suffix closure.  In contrast,
$\sigma(\Omega_B)\subset\Omega_B$ forces $k-1\in B$ whenever
$k\in B\setminus\{0\}$; an infinite such $B$ must equal $\Nzero$.  Hence the
proper sparse sets producing \eqref{eq:abstract-obstruction} cannot occur as
these autonomous control name systems.  The equation is not a deterministic
control counterexample and yields no failure statement for the unrestricted
Borel invariant-partition class.
\end{example}

Appendix~\ref{app:pointwise} supplies autonomous versions of the
obstruction. In particular, Proposition~\ref{prop:finite-pointwise}
uses four distinct controls and has $L_1=M_1=\log2$, although no
sequence attains every feedback's own maximal-pattern rate. It therefore
separates pointwise common attainment from a common full-class minimum.


\section{Realizing one-sided subshifts as safety systems}\label{sec:realization}

Let $E$ be a finite nonempty discrete alphabet and let
$\varnothing\ne Y\subset E^{\Nzero}$ be closed with $\sigma(Y)\subset Y$.
Write
\begin{align*}
 \Lcal_N(Y)=\{y_0\cdots y_{N-1}:y\in Y\},
 \qquad
 P_Y(F)=\#\{y|_F:y\in Y\}.
\end{align*}
The coordinate complexities are
\begin{align*}
 h_{\rm seq,coord}^A(Y)
 &=\limsup_{n\to\infty}\frac1n
   \log P_Y(\{a_1,\ldots,a_n\}),\\
 h_{\rm pat,coord}^*(Y)
 &=\lim_{n\to\infty}\frac1n
   \log\max_{\substack{F\subset\Nzero\\\#F=n}}P_Y(F).
\end{align*}

Adjoin an isolated cemetery point and give
$X=Y\sqcup\{\dagger\}$ the topological disjoint-union topology.  Put $U=E$
and define
\begin{equation}\label{eq:matching-control}
 f(y,u)=
 \begin{cases}
   \sigma y,&y_0=u,\\
   \dagger,&y_0\ne u,
 \end{cases}
 \qquad
 f(\dagger,u)=\dagger.
\end{equation}

\begin{theorem}[Symbolic-to-control realization]\label{thm:C}
For the matching-control system \eqref{eq:matching-control}:
\begin{enumerate}
\item $X$ is compact metrizable, $f:X\times U\to X$ is continuous, and $Y$
is a compact controlled-invariant set.
\item For every $N\in\N$,
\begin{equation}\label{eq:rinv-language}
 r_{\rm inv}(N,Y)=\#\Lcal_N(Y),
\end{equation}
and therefore
\begin{equation}\label{eq:hinv-htop}
 \hinv(Y)=\htop(Y,\sigma)
 =\lim_{N\to\infty}\frac1N\log\#\Lcal_N(Y).
\end{equation}
\item At period one identify $U^1$ with $U$.  If
$\C=(\{D_i:i\in I_\C\},1,\nu)\in\mathfrak C_{\Bor}(Y,1)$, set
$\lambda_\C(i)=\nu(D_i)\in E$.  Then every atom satisfies
\begin{equation}\label{eq:atom-cylinder}
 D_i\subset[\lambda_\C(i)]_Y
 :=\{y\in Y:y_0=\lambda_\C(i)\}.
\end{equation}
Moreover, the one-block label map
\begin{equation}\label{eq:label-factor}
 \Lambda_\C:\Sigma_\C\longrightarrow Y,
 \qquad (a_j)_{j\ge0}\longmapsto(\lambda_\C(a_j))_{j\ge0},
\end{equation}
is a continuous, onto, shift-commuting map.
\item For every $A\in\mathscr S$,
\begin{align}
 \inf_{\C\in\mathfrak C_{\Bor}(Y,1)}\hseq^A(\C;Y)
 &=\inf_{\C\in\mathfrak C_{\clop}(Y,1)}\hseq^A(\C;Y)
 =h_{\rm seq,coord}^A(Y),
 \label{eq:seq-realization}\\
 \inf_{\C\in\mathfrak C_{\Bor}(Y,1)}\hpat(\C;Y)
 &=\inf_{\C\in\mathfrak C_{\clop}(Y,1)}\hpat(\C;Y)
 =h_{\rm pat,coord}^*(Y),
 \label{eq:pat-realization}\\
 \inf_{\C\in\mathfrak C_{\Bor}(Y,1)}h(\C;Y)
 &=\inf_{\C\in\mathfrak C_{\clop}(Y,1)}h(\C;Y)
 =\htop(Y,\sigma).
 \label{eq:ordinary-realization}
\end{align}
Every Borel and clopen infimum in
\eqref{eq:seq-realization}--\eqref{eq:ordinary-realization} is attained by
the same first-coordinate cylinder partition.
\end{enumerate}
\end{theorem}

\begin{proof}
The space $X$ is the union of the two clopen compact sets $Y$ and
$\{\dagger\}$, hence is compact metrizable.  For fixed $u\in E$, the sets
$[u]_Y$, $Y\setminus[u]_Y$, and $\{\dagger\}$ are clopen; on them $f_u$ is,
respectively, the shift or a constant.  Thus $f_u$ is continuous.  Since $U$
is finite discrete, $f$ is jointly continuous.  Choosing $u=y_0$ sends $y$
to $\sigma y\in Y$, proving controlled invariance.

A direct induction gives
\begin{equation}\label{eq:safe-iff-match}
 \varphi(n,y,\omega)\in Y\ (0\le n\le N)
 \quad\Longleftrightarrow\quad
 \omega_j=y_j\ (0\le j<N).
\end{equation}
Indeed, after $j$ matches the state is $\sigma^j y$, whereas a mismatch enters
$\dagger$ permanently.  Hence a length-$N$ control word is safe from at
least one state of $Y$ exactly when it belongs to $\Lcal_N(Y)$, and then it
safely covers precisely its nonempty prefix cylinder.  Distinct language
words give disjoint cylinders. More explicitly:

\emph{Lower bound.} For each $w\in\Lcal_N(Y)$ choose $y^w\in Y$
with prefix $w$. If $S\subset U^N$ spans $Y$, some member of $S$ is safe
from $y^w$. By \eqref{eq:safe-iff-match} that member must equal $w$.
Thus $\Lcal_N(Y)\subset S$, and
$r_{\rm inv}(N,Y)\ge\#\Lcal_N(Y)$.

\emph{Upper bound.} Take $S=\Lcal_N(Y)$. For any $y\in Y$, its
prefix $y_0\cdots y_{N-1}$ belongs to $S$ and is safe through time $N$
by \eqref{eq:safe-iff-match}. Hence
$r_{\rm inv}(N,Y)\le\#\Lcal_N(Y)$.
This proves \eqref{eq:rinv-language}.  The language counts are
submultiplicative, so their normalized logarithms converge.

For completeness, this limit is the usual topological entropy even when
$\sigma|_Y$ is not onto.  The join of the first-coordinate cylinder cover
through times $0,\ldots,N-1$ has exactly $\#\Lcal_N(Y)$ nonempty atoms.
Conversely, given a finite open cover of compact $Y$, the partition into
cylinders of some fixed length $m$ refines that cover; its $N$-fold dynamical join is
then refined by the length-$(N+m-1)$ cylinders.  Taking exponential rates proves
\eqref{eq:hinv-htop}.

Let $\C=(\{D_i:i\in I_\C\},1,\nu)$ be a Borel invariant partition.  If
$D_i$ contained $y$ with $y_0\ne\lambda_\C(i)$, its assigned control would send $y$ to
$\dagger\notin Y$, contradicting safety.  This proves
\eqref{eq:atom-cylinder}.

Take $a=(a_j)_{j\ge0}\in\Sigma_\C$ and put $z_j=\lambda_\C(a_j)$.  For each
$N$, choose a witness $y^{(N)}$ for the first $N$ symbols of $a$.  Induction
using matching controls and \eqref{eq:atom-cylinder} gives
$y^{(N)}_j=\lambda_\C(a_j)=z_j$ for $0\le j<N$.  Hence
$y^{(N)}\to z$ in the product topology; since $Y$ is closed, $z\in Y$.  Thus
$\Lambda_\C$ is well defined, and it is continuous: the preimage of a cylinder in $E^{\Nzero}$ is a finite
union of cylinders in $I_\C^{\Nzero}$, restricted to $\Sigma_\C$.
It is shift commuting because
$(\Lambda_\C(\sigma a))_j=\lambda_\C(a_{j+1})
=(\sigma\Lambda_\C(a))_j$.  Given $y\in Y$, let $a_j$ be the unique index with
$\sigma^j y\in D_{a_j}$.  Its assigned label is $y_j$, and $y$ witnesses
every finite prefix of this atom itinerary.  Thus \eqref{eq:label-factor} is onto.  Here each prefix witness $y^{(N)}$ may be different. Its first $N$
\emph{control labels} nevertheless equal those of $z$, so convergence to
$z\in Y$ follows. This does not prove that $z$ follows the original atom
name $a$; its orbit may meet different atoms with the same labels.
Thus neither shift surjectivity nor a common atom-name witness is used.

For every finite $F$, the coordinate projection of $\Lambda_\C$ is onto, so
\begin{equation}\label{eq:partition-dominates-coordinate}
 \Patt_\C(F)\ge P_Y(F).
\end{equation}
It follows that every Borel, and hence every clopen, invariant partition has
sequence, maximal-pattern, and consecutive-name entropy at least the
corresponding coordinate quantity.

Let $E_Y=\{e\in E:[e]_Y\ne\varnothing\}$ and define
\begin{align*}
\C_0=(\{D_e=[e]_Y:e\in E_Y\},1,\nu_0),
 \qquad \nu_0(D_e)=e.
\end{align*}
This is a clopen invariant partition.  Under the symbol identification,
$W_N(\C_0;Y)=\Lcal_N(Y)$ and $\Sigma_{\C_0}=Y$; the latter follows from
closedness of $Y$.  Equality holds in \eqref{eq:partition-dominates-coordinate}
for every $F$.  Thus $\C_0$ attains all the lower bounds, proving
\eqref{eq:seq-realization}--\eqref{eq:ordinary-realization}.
\end{proof}


\section{A strict Thue--Morse separation}\label{sec:thue-morse}

Let $\vartheta(0)=01$, $\vartheta(1)=10$, and let
\begin{align*}
 \mathbf t=\lim_{r\to\infty}\vartheta^r(0)
 =0110100110010110\cdots.
\end{align*}
Equivalently, $t_n=s_2(n)\pmod2$, where $s_2(n)$ is the number of ones in
the binary expansion of $n$.  Let
\begin{equation}\label{eq:thue-morse-shift}
 Y=\overline{\{\sigma^m\mathbf t:m\ge0\}}
 \subset\{0,1\}^{\Nzero}.
\end{equation}
This is the one-sided Thue--Morse orbit closure.  The following direct computation supplies both rates used below.

\begin{proposition}[Exact symbolic rates]\label{prop:thue-morse}
For every $N\ge1$,
\begin{equation}\label{eq:TM-linear}
 \#\Lcal_N(Y)\le8N.
\end{equation}
Hence $\htop(Y)=0$.  On the single schedule
\begin{equation}\label{eq:TM-schedule}
 A_{\rm TM}=(a_i)_{i\ge1},
 \qquad a_i=4^{i-1}-1,
\end{equation}
one has, for every $k\ge1$,
\begin{equation}\label{eq:TM-full-pattern}
 P_Y(\{a_1,\ldots,a_k\})=2^k.
\end{equation}
The chosen coordinate set contains zero.  Since a cylinder is clopen, its
patterns in the orbit closure are exactly the patterns obtained along shifts
of $\mathbf t$; this calculation uses actual shifts of the fixed point.  Consequently
\begin{equation}\label{eq:TM-complexities}
 h_{\rm seq,coord}^{A_{\rm TM}}(Y)
 =h_{\rm pat,coord}^*(Y)=\log2.
\end{equation}
\end{proposition}

\begin{proof}
A finite word is in $\Lcal_N(Y)$ exactly when it occurs as a factor of
$\mathbf t$: one direction follows from the orbit, and the other from
convergence in a prefix cylinder.  Choose $\ell=\lceil\log_2N\rceil$.  The
identity $\mathbf t=\vartheta^\ell(\mathbf t)$ decomposes $\mathbf t$ into
aligned supertiles of length $2^\ell$.  Since $N\le2^\ell$, every length-$N$
factor is contained in $\vartheta^\ell(ab)$ for an adjacent two-letter factor
$ab$.  There are at most four choices for $ab$ and at most $2^\ell$ starting
offsets in the first supertile.  Therefore
$\#\Lcal_N(Y)\le4\,2^\ell\le8N$, proving zero topological entropy.

Fix $k\ge1$ and
$\varepsilon=(\varepsilon_0,\ldots,\varepsilon_{k-1})\in\{0,1\}^k$.
Put
\begin{align*}
 r=\#\{0\le i<k:\varepsilon_i=0\},
 \qquad
 q=\sum_{\substack{0\le i<k\\\varepsilon_i=0}}4^i
 +(r\bmod2)4^k.
\end{align*}
Only even binary positions occur in this sum.  The binary digit sum of $q$ is
$r+(r\bmod2)$, which is even, so $t_q=0$.  If $\varepsilon_i=1$, bit $2i$
of $q$ is zero, and adding $4^i=2^{2i}$ changes the digit-sum parity.  If
$\varepsilon_i=0$, bit $2i$ is one and bit $2i+1$ is zero; adding $2^{2i}$
replaces these two bits by zero and one without a further carry, so the parity
is unchanged.  In both cases
\begin{equation}\label{eq:TM-digit-identity}
 t_{q+4^i}=\varepsilon_i\qquad(0\le i<k).
\end{equation}
The extra bit at position $2k$ only corrects the initial parity of $q$.
For $i<k$, a carry stops at the empty odd position $2i+1$, which is
strictly below $2k$. The parity correction therefore interferes with
none of these tests. Each addition $q+4^i$ is performed separately;
we are not adding all the $4^i$'s successively to one integer.

Set $n=q+1$. Since $a_{i+1}=4^i-1$,
\begin{align*}
 (\sigma^n\mathbf t)_{a_{i+1}}
 =t_{q+4^i}=\varepsilon_i.
\end{align*}
Thus every binary $k$-pattern occurs on the first $k$ members of
$A_{\rm TM}$, proving \eqref{eq:TM-full-pattern}.  The binary alphabet gives
the reverse upper bound $p_Y^*(k)\le2^k$, and
\eqref{eq:TM-complexities} follows.
\end{proof}

\begin{corollary}[Zero invariance entropy, full binary sparse-time complexity]
\label{cor:separation}
Apply the matching-control construction \eqref{eq:matching-control} to the
Thue--Morse subshift \eqref{eq:thue-morse-shift}.  Then
\begin{equation}\label{eq:strict-control-separation}
 \hinv(Y)=0
 <\log2
 =\inf_{\C\in\mathfrak C_{\Bor}(Y,1)}\hpat(\C;Y)
 =\inf_{\C\in\mathfrak C_{\clop}(Y,1)}\hpat(\C;Y).
\end{equation}
Moreover, the explicit schedule $A_{\rm TM}$ satisfies
\begin{align}
 \max_{A\in\mathscr S}\inf_{\C\in\mathfrak C_{\Bor}(Y,1)}
 \hseq^A(\C;Y)
 &=\inf_{\C\in\mathfrak C_{\Bor}(Y,1)}\hseq^{A_{\rm TM}}(\C;Y)
 =\log2,
 \label{eq:TM-Borel-sequence}\\
 \max_{A\in\mathscr S}\inf_{\C\in\mathfrak C_{\clop}(Y,1)}
 \hseq^A(\C;Y)
 &=\inf_{\C\in\mathfrak C_{\clop}(Y,1)}\hseq^{A_{\rm TM}}(\C;Y)
 =\log2.
 \label{eq:TM-clopen-sequence}
\end{align}
\end{corollary}

\begin{proof}
Combine Theorem~\ref{thm:C} with Proposition~\ref{prop:thue-morse}.  For any
schedule the binary coordinate complexity is at most $\log2$, while
$A_{\rm TM}$ attains $\log2$.  Equations
\eqref{eq:seq-realization}--\eqref{eq:pat-realization} give the Borel and
clopen infima.
\end{proof}

The separation persists after minimization over all Borel invariant
partitions at period one. It therefore distinguishes the fixed-period
pattern infimum from invariance entropy.


\section{Comparison, all periods, and limitations}\label{sec:comparison}

\subsection{Classical entropy and sequence entropy}\label{sec:background}

For a finite open cover $\mathcal V$ of a compact space,
let $N(\mathcal V)$ be the minimum
cardinality of a subcover. The join of finitely many covers consists of
all intersections of one member from each cover, with empty intersections
removed. For a continuous map $T:Z\to Z$ and a finite open cover
$\mathcal V$, ordinary cover entropy is
\begin{align*}
 h_{\mathrm{top}}(T,\mathcal V)
 =\lim_{n\to\infty}\frac1n\log N\left(
 \bigvee_{i=0}^{n-1}T^{-i}\mathcal V\right).
\end{align*}
Submultiplicativity of cover numbers gives this limit. Taking the
supremum over finite open covers defines $h_{\mathrm{top}}(T)$.
For $A=(a_i)_{i\ge1}$ with $0\le a_1<a_2<\cdots$, set
\begin{align*}
 h_{\mathrm{top}}^A(T,\mathcal V)
 &=\limsup_{n\to\infty}\frac1n\log N\left(
     \bigvee_{i=1}^nT^{-a_i}\mathcal V\right),\\
 p_{Z,\mathcal V}^*(n)
 &=\max_{\substack{F\subset\mathbb N_0\\\#F=n}}
    N\left(\bigvee_{i\in F}T^{-i}\mathcal V\right),\\
 h_{\mathrm{top}}^*(T,\mathcal V)
 &=\lim_{n\to\infty}\frac1n\log p_{Z,\mathcal V}^*(n).
\end{align*}
The sequence in the last line is subadditive before division by $n$;
the arbitrary sequence in the first line need not have that property.
This is why a limit and a limsup appear in different definitions.

Huang and Ye's Theorem~2.1(1) states, under their convention that $T$ is a
homeomorphism of a compact metric space, that
\begin{equation}\label{eq:HY-background}
 h_{\mathrm{top}}^*(T,\mathcal V)
 =\max_A h_{\mathrm{top}}^A(T,\mathcal V).
\end{equation}
Its proof selects long nearly maximizing time sets and separates them
by translations. Surjectivity preserves the cover number under a common
translation. The previous sets contribute a vanishing fraction of the
selected observations. For a non-surjective map the translation step
need not preserve cover numbers; our proof uses deletion instead.

If $\mathcal V$ is a clopen partition, the join count is exactly the
number of names seen on the selected coordinates. An arbitrary open
cover may have overlaps, so its minimum subcover count must not be
replaced by the number of nonempty intersections. Similarly, the
measurable version of \eqref{eq:HY-background} uses Shannon entropy
$H_\mu(\alpha)=-\sum_{D\in\alpha}\mu(D)\log\mu(D)$, not the logarithm
of a support cardinality. Our language counts use the latter quantity.
Theorem~\ref{thm:A} is a counting statement, not a measurable
common-sequence theorem.

For a finite-alphabet subshift $Y$, the first-coordinate cylinders give
one particular clopen partition. Its coordinate sequence rate is a
cover-level quantity. Unlike ordinary topological entropy, sequence
entropy need not be computed by a generating partition: replacing
one-coordinate observations by fixed-length windows can change the
normalized sparse rate. We therefore keep the subscript
``$\mathrm{coord}$'' in the realization theorem and do not silently
identify this rate with the supremum over all open covers.

Kawan's Chapter~2 supplies the other interface. Definition~2.2 defines
invariance entropy from finite-horizon spanning controls.
Definitions~2.8 and 2.9 introduce invariant covers and their admissible
word entropies. Proposition~2.21 proves word-count submultiplicativity;
Proposition~2.22 bounds invariance entropy by invariant-cover entropy.
Lemma~2.3 disjointizes a cover without increasing its name entropy.
It does \emph{not} turn arbitrary nonmeasurable sets into Borel sets.
Borel measurability in Theorem~2.3 comes from the safe sets associated
with finitely many spanning controls and their successive differences.
The theorem takes an infimum over sampling periods as well as
partitions. We provide the discrete-time safe-set argument in
Section~\ref{sec:comparison}, where its application to pattern rates is
also made explicit.


\subsection{Comparison with Kawan's invariance entropy}
For every invariant partition, consecutive coordinates are one permissible
choice in the maximal-pattern count. Therefore
\begin{equation}\label{eq:ordinary-pattern-one-sided}
 h(\C;Q)\le\hpat(\C;Q)\le\frac{\log\#I_\C}{\tau}.
\end{equation}
Kawan's Theorem~2.3, with natural logarithms, gives
\begin{equation}\label{eq:kawan-characterization}
 \hinv(Q)=\inf_{\C\in\mathfrak C_{\Bor}(Q)}h(\C;Q).
\end{equation}
Its hypotheses are satisfied by the full-control discrete-time systems
used here: $X$ is metrizable, each controlled map is continuous, and
$Q$ is nonempty, compact and controlled invariant.
The period in this formula is allowed to vary.

For clarity, the relevant argument can be given directly. Spanning
families concatenate, so
$r_{\rm inv}(m+n,Q)\le r_{\rm inv}(m,Q)r_{\rm inv}(n,Q)$ when the
right-hand side is finite. Moreover $r_{\rm inv}(n,Q)$ is nondecreasing.
If any finite-horizon spanning number is finite, then $r_{\rm inv}(1,Q)$
is finite, all of them are finite by concatenation, and Fekete's lemma
gives a limit in \eqref{eq:hinv-definition}. If none is finite, that
limit equals $+\infty$.

Let $S_N=\{\omega^1,\ldots,\omega^m\}$ be a minimal finite
$N$-step spanning family. Its safe sets are
\begin{align*}
 K_i=Q\cap\bigcap_{s=1}^{N}
       \{x\in X:\varphi(s,x,\omega^i)\in Q\}.
\end{align*}
They are closed in $Q$, since finite control compositions are continuous
and $Q$ is closed. They cover $Q$. Set
$D_1=K_1$ and $D_i=K_i\setminus\bigcup_{j<i}K_j$, discard empty sets,
and assign $\omega^i$ to $D_i$. This gives a Borel invariant partition
$\C_N$ of period $N$ with at most $m$ atoms. This is the measurable
construction underlying Kawan's Theorem~2.3; Lemma~2.3 alone only
provides disjointization. The safe sets are already closed before
that lemma is applied.

For the other inequality in \eqref{eq:kawan-characterization}, all
admissible $k$-block words of any $\C$ form a spanning family after their
assigned controls are concatenated. Indeed, from each initial state the
feedback rule selects one such admissible word. Consequently
$r_{\rm inv}(k\tau,Q)\le\#W_k(\C;Q)$, and taking limits gives
$\hinv(Q)\le h(\C;Q)$. The reverse inequality follows from
$h(\C_N;Q)\le N^{-1}\log r_{\rm inv}(N,Q)$ and $N\to\infty$.
If no finite spanning family exists, no finite invariant partition exists,
and both sides are $+\infty$.

Thus, for every nonempty fixed-period subclass
$\mathfrak C\subset\mathfrak C_{\Bor}(Q,\tau)$, the general comparison is
\begin{equation}\label{eq:three-quantities}
 \hinv(Q)
 \le\inf_{\C\in\mathfrak C}h(\C;Q)
 \le\inf_{\C\in\mathfrak C}\hpat(\C;Q).
\end{equation}
The minimax theorem concerns the last quantity. The Thue--Morse
realization proves that it can strictly exceed $\hinv(Q)$.

\subsection{Why the unrestricted all-period infimum collapses}
There is a natural candidate for an all-period Borel quantity,
\begin{align*}
 H_{\rm pat,all}^{*,\Bor}(Q)
 :=\inf_{\tau\in\N}\inf_{\C\in\mathfrak C_{\Bor}(Q,\tau)}
       \hpat(\C;Q).
\end{align*}
The following proposition records the consequence of Kawan's spanning
construction for this all-period quantity.

\begin{proposition}[All-period collapse]\label{prop:all-period-collapse}
In the full-control discrete-time setting of Section~\ref{sec:control}, every
nonempty compact controlled-invariant set $Q$ satisfies
\begin{equation}\label{eq:all-period-collapse}
 H_{\rm pat,all}^{*,\Bor}(Q)=\hinv(Q).
\end{equation}
\end{proposition}

\begin{proof}
The lower bound follows from
\eqref{eq:three-quantities}. For the upper bound use the same
$\C_N$ just constructed. Since its alphabet has at most
$r_{\rm inv}(N,Q)$ symbols, every pattern on $k$ coordinates has at most
$r_{\rm inv}(N,Q)^k$ possibilities. Hence
\begin{align*}
 H_{\rm pat,all}^{*,\Bor}(Q)
 \le\hpat(\C_N;Q)
 \le\frac{\log r_{\rm inv}(N,Q)}{N}.
\end{align*}
Let $N\to\infty$. The infinite-entropy case was already covered by
nonexistence of finite invariant partitions. This proves
\eqref{eq:all-period-collapse}.
\end{proof}

For the matching-control Thue--Morse example, the collapse is particularly
visible. At period $N$, the length-$N$ cylinder partition with its matching
control blocks is clopen, has at most $8N$ atoms and has pattern rate at
most $N^{-1}\log(8N)$. Thus even the all-period clopen infimum is zero
in this example, while the period-one Borel infimum is $\log2$.
Changing the sampling period changes which information is packed into
one observed symbol and the cost assigned to that symbol.

\subsection{Countability, trajectories, and physical time}
The countability hypothesis in Theorem~\ref{thm:B} is used exactly when
one representative language from each class is listed and assigned
infinitely many construction stages. Compact metrizability gives
countably many clopen sets; it does not give countably many Borel sets
or Borel name languages. No approximation theorem allowing that
substitution is proved here. Example~\ref{ex:uncountable} rules out a
formal uncountable extension for arbitrary pattern sets, but its
non-shift-invariant examples do not settle the unrestricted autonomous
control problem. Theorem~\ref{thm:C} bypasses countability in its specific
realization by exhibiting a least language under the label projections.
The countable-language theorem does not settle the general Borel minimax
problem. The closed-cover theorem resolves a particular uncountable
class; the inverse-capacity theorem resolves maximal-pattern optimization,
but does not by itself resolve the corresponding common-sampling problem. Appendix~\ref{app:pointwise} separates common minimum attainment from pointwise attainment in genuine continuous control systems. Appendix~\ref{app:capacity-boundaries} explains why actual transfer-norm discounts and probability-entropy bounds do not automatically supply the missing counting estimate.

The patterns count admissible language projections. Every finite pattern
has a finite trajectory witness which stays safe between observations.
An infinite point in the completed Borel name system may have only
separate finite-prefix witnesses, as Remark~\ref{rem:witness} shows.
This distinction cannot be removed by compactness alone when the atoms
are not closed. The factor in Theorem~\ref{thm:C} is a factor of compact
\emph{name systems}; it is not a continuous factor defined by an arbitrary
Borel feedback map on the original state space.

For one fixed $\tau$, $A$ denotes block indices and $\tau A$ denotes
physical observation times. The normalization $n\tau$ charges one block
length per observation and disregards the lengths of the gaps. A common
block-index sequence for different periods would represent different
physical schedules. Theorem~\ref{thm:A} could be applied to countably many
languages with different positive normalizing constants by rescaling
each equality, but that would not solve the distinct physical-time
minimax problem. Formula~\eqref{eq:all-period-collapse} is an infimum
identity and makes no common-schedule assertion.

Finally, if one safe control block works for every point of $Q$, the
one-atom invariant partition makes all three fixed-period infima zero
for any class containing it. On a connected $Q$, every finite clopen
partition has one atom. The full clopen class is therefore either empty
or has zero pattern infimum, while Borel strategies can remain
nontrivial. These observations explain both the usefulness and the
limits of the clopen sufficient condition. The nontrivial separation
studied here records the richness of safe symbolic instructions at a
specified sampling period, despite a zero asymptotic count per elapsed
time.


\section{Affine examples, common sampling, and sharp single-step boundaries}
\label{sec:affine}

These examples separate fixed-period maximal-pattern optimization,
common sampling, and a fixed arithmetic sampling. Their proofs are
included because a matching upper bound must be supplied by a real
feedback. Volume estimates, integer digit coding, and protected inverse
branches are classical mechanisms, not new general tools.

\subsection{Integer affine branches and an absorbing feedback}
Let $m\ge2$, $D\ge m$ be integers, $R=D/(m-1)$,
$Q=[0,R]$, and controls $u=0,\ldots,D$ act legally by $mx-u$.
On $X=[-1,R+1]$ define the continuous ambient maps
\[
 f_u(x)=\max\{-1,\min\{R+1,mx-u\}\}.
\]
For $x\in Q$, this is in $Q$ exactly when $mx-u$ is in $Q$,
and then truncation has no effect. The safety intervals
$K_u=[u/m,(u+R)/m]$ cover $Q$. Fix any period $\tau\ge1$.

\begin{proposition}\label{prop:integer}
For this system, over all legal Borel block feedbacks and finite
Borel atom refinements,
\[
 \min_s h_{(0,1,\ldots)}(s)=L_\tau=M_\tau=\log m.
\]
There is an actual Borel selector $s_\tau$ with
$h_A(s_\tau)=h^*(s_\tau)=\log m$ for every infinite increasing
$A$, although its absorption time has no uniform bound.
\end{proposition}
\begin{proof}
For any actual length-$n$ block word, concatenate its controls into
$w_0\cdots w_{n\tau-1}$. Safety gives
\[
 x_{n\tau}=m^{n\tau}x-\sum_{j=0}^{n\tau-1}
                   w_jm^{n\tau-1-j}\in[0,R].
\]
Its initial-state fiber is contained in an interval of length
$R/m^{n\tau}$, whether or not that fiber itself is an interval.
These containing intervals cover $Q$, so there are at least
$m^{n\tau}$ actual consecutive block words. This gives the
full-class lower bound at the common sequence $(0,1,\ldots)$.

Define the one-step rule $v(x)=\min\{D,\lfloor mx\rfloor\}$
and $T(x)=mx-v(x)$. If $x<D/m$, the next state is in $[0,1)$;
if $x\ge D/m$, it is $mx-D\in[0,R]$.
The core $[0,1)$ is invariant and uses only
$\Gamma=\{0,\ldots,m-1\}$. Both $0$ and $R$ are fixed;
$D/m$ uses $D$ and goes to zero.
For $x<R$, while $D$ is used the state is
$R-m^j(R-x)$. Since $R-x>0$, eventually it is below $D/m$,
then one transitional control $v<D$ sends it into the core.
Only $R$ uses $D$ forever. The absorption times are unbounded
as $x$ tends to $R$.

From the current state define
$s_\tau(x)=(v(x),v(Tx),\ldots,v(T^{\tau-1}x))$.
It is Borel, each intermediate step is safe, and its actual endpoint
map is $T^\tau$. Put $\mathcal B=\Gamma^\tau$,
$b=m^\tau$, and $d_0=(D,\ldots,D)$.
All its real block names are contained in
\[
 \bigcup_{j\ge0}d_0^j\mathcal C_\tau\mathcal B^{\Nzero}
 \ \cup\ \{d_0^\infty\},\qquad
 |\mathcal C_\tau|\le D\frac{m^\tau-1}{m-1}=:C_\tau.
\]
The transitional block has some initial $D$'s, one control below
$D$, and then core controls; this proves the stated count.
For any window of $n$ block positions, an unobserved transition
leaves just one of $n+1$ cuts between constant initial labels and
core labels. A sampled transition has $n$ possible locations and
at most $C_\tau$ labels. Thus
\[
 \Patt_{s_\tau}(F)\le[n+1+nC_\tau]m^{\tau n},\qquad |F|=n.
\]
There is no factor depending on the last observation time or on
the unbounded absorption time.

For the matching lower count prescribe any $n$ core blocks at the
positions of $F$, fill all other one-step digits with zero, and let
$x=\sum_{j\ge0}e_jm^{-j-1}<1$. There are only finitely many
nonzero digits. Every tail is smaller than one, so
$\lfloor mT^jx\rfloor=e_j$ and these are real controls of the
same rule. Hence $\Patt_{s_\tau}(F)\ge m^{\tau n}$ for every
window. The all-window bounds and the full-class consecutive lower
bound prove the proposition, with minima actually attained.
Control-label projection treats arbitrary atom refinements.
\end{proof}

When $m\ge3$ and $D=m$, the minimum one-step safe subcover has
$r=m+1$, although $L_1=M_1=\log m$. Indeed, $m$ intervals of
length $R/m$ could cover an interval of length $R$ only by abutting
without positive-length overlap. Their left endpoints are multiples
of $1/m$, so this would force their consecutive control differences
to equal the noninteger $R=m/(m-1)$. All $m+1$ intervals are
therefore needed. There is a common state factor
$x\mapsto e^{2\pi ix}$, with induced map $z\mapsto z^m$;
it is preserved by every legal branch. The result cannot be presented
as eliminating common-factor explanations.

\subsection{Protected inverse branches and common sampling}
On $Q=[0,1]$, at period one, use
\[
 g_a(x)=\frac54x,\quad g_b(x)=\frac{20}3x-\frac{16}3,
 \quad g_c(x)=20x-19,\quad g_d(x)=100x.
\]
Their safety domains are
$[0,4/5]$, $[4/5,19/20]$, $[19/20,1]$, and
$[0,1/100]$, respectively. Continuous ambient extensions are
obtained by truncating these affine maps to $[-1,2]$, with no
change on legal paths. The minimum subcover is three and the largest
two inverse slopes sum to $4/5+3/20=19/20<1$; hence
$M_1=\log3$ by the scalar special case of the capacity theorem.

\begin{proposition}\label{prop:bridge}
For $A_2=(0,2,4,\ldots)$, every legal Borel feedback and every
finite Borel atom refinement has at least $3^n$ actual patterns
on its first $n$ observations. Consequently
\[
 \min_s h_{A_2}(s)=L_1=M_1=\log3.
\]
\end{proposition}
\begin{proof}
The full inverse branches are
$\psi_a(y)=4y/5$, $\psi_b(y)=4/5+3y/20$,
$\psi_c(y)=19/20+y/20$. Let $J=[1/2,999/1000]$ and
$\theta_u=\psi_u\psi_b$. They satisfy
\[
 \theta_a(y)=16/25+3y/25,\quad
 \theta_b(y)=23/25+9y/400,\quad
 \theta_c(y)=99/100+3y/400.
\]
Their complete images and the intermediate image are
\[
\begin{array}{c|c|c}
 \text{image}&\text{exact endpoints}&\text{only legal control}\\\hline
 \psi_b(J)&[7/8,18997/20000]&b\\
 \theta_a(J)&[7/10,18997/25000]&a\\
 \theta_b(J)&[149/160,376991/400000]&b\\
 \theta_c(J)&[159/160,398997/400000]&c.
\end{array}
\]
Each $\theta_u(J)$ lies in $J$. Every displayed closed interval
is strictly inside the appropriate uniquely legal region; in
particular $\theta_a(J)$ avoids the extra control $d$.

Given any $u_0\cdots u_{n-1}\in\{a,b,c\}^n$, set
$x_{2n}=1/2$ and recursively
$x_{2j+1}=\psi_b(x_{2j+2})$,
$x_{2j}=\theta_{u_j}(x_{2j+2})$.
The table shows that at every even state only $u_j$ is legal
and at every odd state only $b$ is legal. Thus \emph{every}
feedback must execute the actual period-one prefix
$u_0,b,\ldots,u_{n-1},b$ from the same constructed initial state.
The endpoint remains in $Q$ and the globally legal rule continues
safely forever. There is no added phase or period-two control.
All three-symbol sampled words are therefore actual patterns.
Refinement only increases their count.

The Borel selector using $a$ below $4/5$, $b$ from $4/5$ to
$19/20$, and $c$ at and above $19/20$, is legal at all endpoints
and uses only three labels. Its upper count is $3^n$ on every
window. This matches the lower bound, including attainment, and
$L_1\le M_1$ completes the equalities. The normalization is
$n\tau=n$, not $2n$.
\end{proof}

The closed contracting branches $\theta_a,\theta_b,\theta_c$
have disjoint first-level images in $J$. Their nested images give
a compact set homeomorphic to the three-symbol full shift, and
every $T_s^2$ on that set is the same shift under this coding.
To verify this, a word's image has diameter at most
$(3/25)^n\operatorname{diam}J$; nested images intersect in one
point, different first differing letters give disjoint images, and
the inverse identities used above give the shift relation.
The union of this set with its $\psi_b$ image is protected under
every legal feedback. These sets can overlap: the all-$b$ state
$16/17$ belongs to both and uniquely uses $b$. Hence the proof
does not introduce an implicit disjoint phase register. This is
precisely a protected-subsystem mechanism, even if the rest of the
state space has feedback-dependent successors.

\subsection{A critical boundary of the single-step density condition}
For $3/2<\lambda\le2$, on $Q=[0,1]$ use
\[
 g_a(x)=\lambda x,\quad
 g_b(x)=\lambda x-(\lambda-1)/2,\quad
 g_c(x)=3x-2.
\]
The safety intervals are
\[
 K_a=[0,1/\lambda],\quad
 K_b=[(\lambda-1)/(2\lambda),(\lambda+1)/(2\lambda)],
 \quad K_c=[2/3,1].
\]
They cover $Q$ and $r=3$. Both $a$ and $c$ are forced into
every subcover by endpoints $0,1$. Since $1/\lambda<2/3$,
they leave an open gap, so $b$ is also needed. The inverse
densities are constants $1/\lambda,1/\lambda,1/3$; therefore
$\Theta_2=1/\lambda<1$ and $\Theta_3=2/\lambda$.

\begin{proposition}\label{prop:critical}
For every $3/2<\lambda\le2$, this full Borel class has
$M_1=\log2<\log3$. Thus the condition $\Theta_3<1$ cannot be
uniformly replaced by $\Theta_3\le1$ or by
$\Theta_3<1+\eta$ for any $\eta>0$, while retaining the
conclusion $M_1=\log3$.
\end{proposition}
\begin{proof}
The $k=2$ single-step criterion gives every feedback
$h^*\ge\log2$. Put $t=(\lambda-1)/(2\lambda)$ and define
\[
 s_\lambda(x)=\begin{cases}
 a,&0\le x<t,\\ b,&t\le x<2/3,\\ c,&2/3\le x\le1.
 \end{cases}
\]
It is Borel and legal: $t<1/\lambda$ and
$(\lambda+1)/(2\lambda)\ge3/4>2/3$.
For $x<t$, $g_a(x)<(\lambda-1)/2\le1/2$.
For $t\le x\le1/2$,
$g_b(x)=\lambda x-(\lambda-1)/2\in[0,1/2]$.
Hence $J=[0,1/2]$ is invariant. Before entering $J$ only
$b,c$ occur; once in $J$ only $a,b$ occur. This remains true
for states which never enter and at all thresholds, including the
fixed points $1/2$ and $1$.

At any $n$ selected times a real name has some cut $0\le j\le n$
between the two groups, each of size two. Thus its pattern count
is at most $(n+1)2^n$ for \emph{every} window. This actual
feedback gives $h^*\le\log2$ and completes the equality.
At $\lambda=2$, $\Theta_3=1$; as $\lambda$ increases to $2$
from below, $\Theta_3$ decreases to $1$ from above. These values
give both claimed failures of uniform relaxation. The example
does not make $\Theta_3<1$ a necessary condition for every
system, and does not determine $L_1$ for this family.
\end{proof}

\subsection{A noninteger four-branch system}
\label{subsec:four-affine}
Let $Q=[0,6]$ and use, at period one,
\[
 g_{A0}(x)=\tfrac32x,\quad g_{A1}(x)=\tfrac32x-3,
 \quad g_{B0}(x)=\tfrac53x,\quad g_{B1}(x)=\tfrac53x-4.
\]
The original safety intervals are
$[0,4]$, $[2,6]$, $[0,18/5]$, $[12/5,6]$.
Truncation to the ambient interval $[-1,7]$ retains exactly these
safe transitions. All legal Borel feedbacks, arbitrary mixtures,
and finite atom refinements are allowed. Two $A$ domains cover
$Q$ and no single one does, so $r=2$. The largest scalar
inverse capacity is $2/3<1$. Thus $M_1=\log2$.
Every actual length-$N$ control prefix has its fiber inside an
interval of length at most $6(3/2)^{-N}$; these intervals cover
$Q$. Its number $D_N(s)$ is therefore at least $(3/2)^N$.
Appendix~\ref{app:finiteblock} retains the weights and actual
even--odd compatibility in a refinement of this volume argument.
This gives
\begin{equation}\label{eq:four-L}
 \log(3/2)\le L_1\le\log2.
\end{equation}

For $A_2=(0,2,\ldots)$ let
$\ell_2=\inf_s h_{A_2}(s)$, and let $\ell_2^A$ restrict to
all two-$A$ Borel feedbacks, without restricting their partitions.
The unrefined two-$A$ rules have the exact description
\[
 s_E(x)=A0\text{ on }[0,2)\cup E,
 \quad s_E(x)=A1\text{ elsewhere},\qquad E\subset[2,4]
 \text{ Borel}.
\]
Indeed the two overlap domains are exactly $[2,4]$.
For any feedback, an actual $2n$ prefix is determined by its even
and odd sampled patterns. Odd patterns occur among the even
patterns from states $T_sx\in Q$. Hence
$D_{2n}(s)\le N_n(s)^2$ and the volume lower bound gives
$N_n(s)\ge(3/2)^n$. Consequently $\ell_2\ge\log(3/2)$.

\begin{proposition}[Density-one samplings]\label{prop:density-one}
If $A=(a_i)$ satisfies $a_n/n\to1$, then
$\inf_s h_A(s)=\log(3/2)$ over the full class.
The greedy feedback $s_g=s_\varnothing$ is an actual minimizer
for every such $A$.
\end{proposition}
\begin{proof}
Set $\beta=3/2$. On the invariant core $[0,3)$, scale $x=3t$.
The greedy dynamics becomes $B(t)=\beta t-\lfloor\beta t\rfloor$
on $[0,1)$. Every nonempty length-$n$ digit cylinder is a
left-closed right-open interval whose $B^n$ image is $[0,l)$
for some $0<l\le1$. This is proved by splitting the current
image at $1/\beta$: it has a single child if $\beta l\le1$,
and otherwise a full-image child and a nonfull-image child.
Equality creates no extra singleton at the excluded right endpoint.

Let $F_n$ be the number of full-image cylinders, with $F_0=1$.
The sum of image lengths is $\beta^n$, since each source length
is $l/\beta^n$ and the sources partition $[0,1)$. Thus
$F_n\le\beta^n$. Each nonfull cylinder has a last full-image
ancestor, and each ancestor has at most one descendant at a given
later level without another full-image step. The total core count
$C_n$ consequently satisfies
\[
 C_n\le\sum_{j=0}^n F_j\le3\beta^n.
\]
For $x\in[3,6)$ the real greedy control remains $A1$ until
entering the core; while waiting its state is
$6-\beta^j(6-x)$, so entry is finite. State $6$ uses $A1$
forever. Therefore the full-space count satisfies
\[
 D_n(s_g)\le1+\sum_{j=1}^n C_j\le10\beta^n.
\]
This includes all thresholds and the exceptional endpoint.

Let $N=a_n+1$, so $N-n=o(n)$. Each sampled word has at most
$4^{N-n}$ completions to a consecutive word, giving
\[
 \Patt_s(\{a_1,\ldots,a_n\})\ge
 \frac{D_N(s)}{4^{N-n}}\ge\frac{\beta^N}{4^{N-n}}.
\]
For the greedy rule the same count is at most
$D_N(s_g)\le10\beta^N$. After division by $n$ the two
rates tend to $\log\beta$. Refinements preserve the lower
bound and the unrefined rule gives the matching upper bound.
\end{proof}

Let $\rho>1$ be the unique solution
\begin{equation}\label{eq:rho}
 \rho^6=\rho^5+\rho^4+\rho+1,
 \qquad\log\rho=0.55600938749\ldots.
\end{equation}
Appendix~\ref{app:greedy} proves the exact real-feedback value
$h_{A_2}(s_g)=\log\rho$, with full upper and lower counts.
Thus the currently established intervals are
\begin{equation}\label{eq:four-even}
 \log(3/2)\le\ell_2\le\ell_2^A\le\log\rho.
\end{equation}
Neither greedy optimality in the full class nor greedy optimality
in the two-$A$ class is established. These fixed-$A_2$ questions
are distinct from \eqref{eq:four-L}. A lower entropy abstract
open-loop language is constructed in Appendix~\ref{app:language};
its full Borel feedback implementation remains unresolved.


\section{Expansion, absence of continuous common factors, and phase obstructions}
\label{sec:extensions}

\subsection{Expansion extensions of the integer family}
The next statement retains a matching real feedback while allowing other
feedbacks to change the dynamics inside its absorbing core. It is a
geometric extension of Proposition~\ref{prop:integer}, rather than a
new digit-coding theorem.

\begin{proposition}\label{prop:extension}
Start with the integer family of Proposition~\ref{prop:integer}. Add
finitely many continuous branches $g_v:Q\to\mathbb R$ satisfying
\begin{equation}\label{eq:global-expansion}
 |g_v(x)-g_v(y)|\ge m|x-y|\qquad(x,y\in Q).
\end{equation}
A branch is legal precisely when its value is in $Q$.
For every fixed period $\tau$, over the entire enlarged Borel class,
\[
 \min_s h_{(0,1,\ldots)}(s)=L_\tau=M_\tau=\log m.
\]
If among the added branches is $g_e(x)=kmx$, with an integer
$k\ge2$ and $R\ge k$, there is no nonconstant continuous map
$\pi:Q\to Z$, to any compact metric space, and no map $H:Z\to Z$
such that $\pi(T_sx)=H(\pi(x))$ for all legal period-$\tau$
feedbacks and all states. Continuity of $H$ is not required.
\end{proposition}
\begin{proof}
Use the continuous ambient extension
$f_v=c_X\circ g_v\circ c_Q$ on $X=[-1,R+1]$, where $c$ denotes
interval truncation. A value outside $Q$ remains outside $Q$ after
truncation to $X$. The original branches still cover $Q$.
Two states in the same real length-$n$ block fiber follow the same
$n\tau$ legal controls, so \eqref{eq:global-expansion} gives
$m^{n\tau}|x-y|\le R$. Their fiber, however fragmented, has
diameter at most $Rm^{-n\tau}$. Its containing interval, together
with the other containing intervals, covers $Q$. There must be at
least $m^{n\tau}$ real consecutive words. The original actual
feedback $s_\tau$ is still available and retains every all-window
count proved in Proposition~\ref{prop:integer}. These matching
bounds prove the equalities and treat refinements by projection.

Suppose the asserted common continuous factor exists. Any legal block
at a given state can be chosen there by changing a globally safe
Borel selector at that singleton. Thus all legal blocks at that
state have the same projected endpoint. At $x=y/m^\tau$, for
$1\le y\le R$, the legal blocks $0^\tau$ and $0^{\tau-1}1$
end at $y,y-1$. All intermediate states are in $Q$. Hence
$\pi(y)=\pi(y-1)$ throughout $[1,R]$. The restriction to $[0,1]$
induces a continuous $F:\mathbb R/\mathbb Z\to Z$, and repeated
subtraction of integers gives $\pi(y)=F(y\bmod1)$ on $Q$.
At the same initial state the blocks $0^\tau$ and $0^{\tau-1}e$
end at $y,ky$, provided $0\le y\le R/k$. Since $R\ge k$,
this implies $F(t)=F(kt\bmod1)$ for all $t$ on the circle.
Iterating gives $F(j/k^n)=F(0)$ for every $0\le j<k^n$.
Their union is dense, so continuity makes $F$, and then $\pi$,
constant, a contradiction.
\end{proof}

For example, $m=2,D=3,Q=[0,3]$ admits both $g_e=4x$ and
$g_v=2x+x^2$ as additional branches. At $x=1/8$ the legal
successors under $0,e,v$ are $1/4,1/2,17/64$, all in the original
core. A rule choosing $e$ on $(0,1/4)$ and otherwise the original
greedy control is globally Borel and safe. These are genuine
feedback-dependent core dynamics, and the proposition excludes every
nonconstant continuous common factor. It does not exclude Borel
factors: for integer affine branches and $e$, rationality of the
state is a nonconstant Borel common factor.

Condition~\eqref{eq:global-expansion} cannot be replaced by expansion
on each monotonicity interval. In the $m=2,D=3$ system add
$g_b(x)=2x$ on $[0,3/2]$ and $g_b(x)=6-2x$ on $[3/2,3]$.
Both slopes have absolute value two, but this single branch maps
all of $Q$ into $Q$. The constant real block $b^\tau$ has one
instruction name; both optimized rates are zero. Folding hides
the spatial branching from the control-label alphabet.

\subsection{An alternating open-loop language with no Borel realization}
Let $R>0$, $1<a,b<2$, and $\log a/\log b\notin\mathbb Q$.
On $[0,R]$ take four legal affine branches
\[
 g_{Ai}(x)=ax-i(a-1)R,\qquad
 g_{Bi}(x)=bx-i(b-1)R\quad(i=0,1).
\]
Each type's two full inverse intervals cover $Q$. Let
$\Lambda_{\rm alt}$ be the two-phase language whose $A,B$ types
alternate strictly and whose digits are otherwise arbitrary. Let
\[
 \widehat\Lambda=\Lambda_{\rm alt}\ \cup\
 \{w c^\infty:w\text{ is a finite alternating-type word},\ c\in U\}.
\]
The empty prefix is allowed. This compact shift-invariant language
includes all constant endpoint tails and their finite alternating
prehistories. The inverse contractions
$\psi_{Ci}(y)=(y+i(\lambda_C-1)R)/\lambda_C$, with
$\lambda_A=a,\lambda_B=b$, define a continuous coding
$\pi(d)=\lim_n\psi_{d_0}\cdots\psi_{d_{n-1}}(0)$.
Every state has a coding in $\Lambda_{\rm alt}$: in each prescribed
type choose a safe digit and repeat. The contraction error is at
most $R\min(a,b)^{-n}$. This proves full open-loop coverage, not
feedback consistency.

\begin{proposition}\label{prop:alternation}
No globally safe Borel period-one feedback has all its real names
in $\widehat\Lambda$. This remains true with arbitrary Borel
mixtures and finite atom refinements. Nevertheless for every
$n$-position window $F$,
\[
 2^{n+1}\le\Patt_{\widehat\Lambda}(F)\le8(n+1)2^n.
\]
\end{proposition}
\begin{proof}
The alternating part has two distinct type patterns and $2^n$
independent digits. For a constant-tail member there are at most
$n+1$ cuts on $F$, two prefix phases, $2^j$ prefix digits, and
four constant tails. Summing by cuts bounds the count by
$8(n+1)2^n$ (and also covers the infinite alternating case).

For nonrealizability, let $C$ be the union of all finite affine
preimages of $\{0,R\}$. This set is countable, and its complement is preserved by every
legal branch: a successor in $C$ would put the original state in $C$. A constant safe tail under any one
branch forces its tail state to be the corresponding endpoint,
by expansion about $0$ or $R$. Thus any hypothetical feedback
has strictly alternating types forever on $Q\setminus C$.
Choose
\[
 0<\eta<\min\{R/a,R/b,(a-1)R/a,(b-1)R/b\}.
\]
On $(0,\eta)$ only zero-digit controls are legal. For
$0<x<\eta/(ab)$ outside $C$, the true first two steps are
$x,ax,abx$ if the first type is $A$, and $x,bx,abx$ otherwise.
The third type is the first type. Put $p=\log a,q=\log b,h=p+q$
and encode the feedback type at $\eta e^t$, $t<0$, by a Borel
$c(t)\in\{0,1\}$, with $A$ represented by one. Outside a
countable exceptional set,
\[
 c(t+h)=c(t),\quad
 c(t)=1\Rightarrow c(t+p)=0,\quad
 c(t)=0\Rightarrow c(t+q)=1\qquad(t<-h).
\]
Translating each negative interval of length $h$ to $[-h,0)$
gives a measurable circle function, with all countably many
translated exceptions still null. Apply the last two implications
at $t=r-h$, $-h\le r<0$. For its set $E$ of type one, circle
rotations satisfy $R_pE\subset E^c$ and $R_qE^c\subset E$
modulo null sets. Since $q=-p$ modulo $h$, these imply
$R_pE=E^c$, measure $1/2$, and $R_{2p}E=E$ modulo null sets.
The number $2p/h$ is irrational. For
$G=2\one_E-1$, invariance under that rotation and Fourier
coefficients give
$\widehat G(j)=e^{2\pi i j(2p/h)}\widehat G(j)$.
Every nonzero coefficient vanishes; the mean vanishes as well.
Completeness of the trigonometric system gives $G=0$ almost
everywhere, contradicting $|G|=1$. This proves impossibility.
Refinements project to the same physical control rule and cannot
avoid the obstruction.
\end{proof}

The noninteger system in Section~\ref{subsec:four-affine} has
$a=3/2,b=5/3,R=6$. Its log ratio is irrational by unique prime
factorization: $(3/2)^j=(5/3)^k$ would imply
$3^{j+k}=2^j5^k$. None of its nonempty finite affine composites
has an integer slope. The proposition is therefore a substantive
implementation boundary for this system; it is not a proof of a
strict optimized entropy gap.

\section{Further questions}\label{app:scope}
The main remaining question is whether hypotheses that force
$M_\tau=\log r/\tau$ through product capacities also force
$L_\tau=M_\tau$, or whether there is a genuine strict numeric gap.
The positive-density independent positions furnished here may vary
with the feedback, so they do not answer this question. For the
four-branch affine system, both fixed-even-sampling infima and greedy
optimality remain open within the stated intervals. The ten-block
language's full Borel implementation and the full-support scalarization
of the smooth example are also unresolved.


\appendix

\section{Exact even-sampling entropy of the real greedy feedback}
\label{app:greedy}

We prove $h_{A_2}(s_g)=\log\rho$ for the system of
Section~\ref{subsec:four-affine}. An upper language bound is followed by
inverse paths that realize a family with the same exponential rate.
On $C_0=[0,3)$ the actual map $R=T_{s_g}^2$ is
\begin{equation}\label{eq:greedy-R}
 R(x)=\begin{cases}
 9x/4,&0\le x<4/3,\\
 9x/4-3,&4/3\le x<2,\\
 9x/4-9/2,&2\le x<3.
 \end{cases}
\end{equation}
The middle state really uses $A0$ or $A1$ according to the original
threshold, so these are actual two-step trajectories of a period-one
rule. Let $0,1$ denote its even labels $A0,A1$.

If a core even name begins $11$, its second state is in
$[2,9/4)$, its third in $[0,9/16)$, and the next two images
are in $[0,81/64)$ and $[0,729/256)$. The first of these
is below $4/3$. If the last image is at least $2$, its next
image is smaller than $1953/1024<2$.
Consequently every core even name avoids
\begin{equation}\label{eq:greedy-forbidden}
 111,\qquad1101,\qquad110011.
\end{equation}
All interval upper endpoints in these calculations are excluded by
the core or its actual branch convention; $4/3$ and $2$ both
map to zero under $R$.

Every word avoiding \eqref{eq:greedy-forbidden} has a unique
left-to-right parsing into the prefix-code blocks
$a=0$, $b=10$, $c=1100$, with no adjacent $cc$, followed by
one of $\varnothing,1,11,110$. The first two exclusions bound
one-runs and force two zeros after a full double-one run; the third
prevents two complete $c$ blocks. Let $H_m$ count complete block
strings of bit length $m$ with no $cc$, with $H_m=0$ for
$m<0$. If $X$ and $Z$ count strings ending in $a/b$ and in
$c$, respectively, their formal generating series satisfy
\[
 X=(z+z^2)H,\qquad Z=z^4(1+X),\qquad H=1+X+Z.
\]
Thus
\begin{equation}\label{eq:greedy-H}
 H(z)=\frac{1+z^4}{1-z-z^2-z^5-z^6},\qquad
 H_m=H_{m-1}+H_{m-2}+H_{m-5}+H_{m-6}\quad(m\ge5),
\end{equation}
with $H_0,\ldots,H_4=1,1,2,3,6$ and $H_5=10$.
The equation for $\rho$ in \eqref{eq:rho} is equivalent to
$\rho^{-1}+\rho^{-2}+\rho^{-5}+\rho^{-6}=1$, whose left side
is strictly decreasing. Its unique root obeys $5/3<\rho<7/4$.
The initial counts are at most $\rho^m$, and the recurrence
proves $H_m\le\rho^m$ at all lengths. The incomplete-tail
count is at most
$H_m+H_{m-1}+H_{m-2}+H_{m-3}\le D\rho^m$, where
$D=1+\rho^{-1}+\rho^{-2}+\rho^{-3}$.

For a real full-space even word, either all labels are $1$, or its
first $0$ occurs at some position $j$. At that moment the state
is below $2$ and lies in the core; its remaining word is governed
by \eqref{eq:greedy-forbidden}. This includes arbitrarily long
waiting outside the core and the fixed endpoint $6$.
Therefore
\begin{equation}\label{eq:greedy-upper}
 N_n(s_g)\le\sum_{m=0}^n D\rho^m
 \le\frac{D\rho}{\rho-1}\rho^n.
\end{equation}
This is a real feedback upper bound for every length, not an equality
of its language with the larger forbidden-block language.

\subsection*{Inverse paths furnishing the matching lower bound}
Put $\eta=4/9$, $L=162/65$, $d_0=729/256$, and define
\[
 P_0=[0,16/27]\cup[72/65,4/3),\qquad P_1=[L,3),
\]
\[
 C=P_0\cup P_1,\qquad
 D_0=P_0\cup[2,550/243]\cup[L,d_0).
\]
These are subsets of the original interval used only to prove
realizability, not feedback states or phase variables. The following
inverse paths realize the indicated even blocks:
\begin{center}
\begin{tabular}{cclc}
\toprule
Even block&Target $y$&Source state&Real one-step controls\\
\midrule
$a=0$&$y\in C$&$\eta y$&$00$\\
$b=10$&$y\in C$&$2+\eta^2y$&$1000$\\
$b=10$&$y\in P_0$&$70/27+\eta^2y$&$1001$\\
$c=1100$&$y\in D_0$&$26/9+\eta^4y$&$10100000$\\
\bottomrule
\end{tabular}
\end{center}
The left $b$ rule has its source in $D_0$ for every $y\in C$,
and in $C\cap D_0$ when $y\in P_1$. For a target in $P_0$,
the right $b$ rule supplies a source in $C$.

Here are all intermediate checks. For $a$, the source
$\eta y$ is in $P_0$, the intermediate state is $2y/3<2$,
and the endpoint is $y$; both controls are zero.
For left $b$, the states after its source are
$8y/27$, $\eta y$, $2y/3$, $y$. If $y\in P_0$, the
source is in $[2,550/243)$; if $y\in P_1$, it is in
$[L,70/27)$. Each source is above $2$ and outside
$I=[12/5,L)$, while the three subsequent control states are
below $2$. This gives $1000$ with the stated source memberships.

For right $b$ the states are
$8/9+8y/27$, $4/3+\eta y$, $2+2y/3$, $y$.
The source belongs to $[70/27,694/243)\subset P_1$.
The first two subsequent control states are below $2$.
If $y\in[0,16/27]$, the last control state is at most
$194/81<12/5$; if $y\in[72/65,4/3)$, it is at least
$178/65>L$ and below $26/9$. Thus it is above $2$ and
outside $I$, proving $1001$ at all included endpoints.

For $c$, its source is in $[26/9,3)\subset P_1$, since
$26/9+\eta^4d_0=3$ and $d_0$ is excluded from $D_0$.
Its successive states are
\[
\begin{array}{llll}
4/3+(3/2)\eta^4y<3/2,&2+\eta^3y<9/4,&
(3/2)\eta^3y<3/8,&\eta^2y<9/16,\\
(3/2)\eta^2y<27/32,&\eta y<81/64,&
(3/2)\eta y<243/128<2,&y.
\end{array}
\]
Only its source and its second subsequent state are above $2$;
both avoid $I$. This proves $10100000$ with every intermediate
state safe. We have obtained true inverse relations
\[
 C\xrightarrow{a,b}C,\qquad C\xrightarrow cD_0,
 \qquad D_0\xrightarrow{a,b}C,
\]
where each indicated target set is covered in its entirety by the
specified inverse rule.

For any complete block string without $cc$, start with proof label
$C$ and update to $D_0$ after $c$, and to $C$ after $a,b$.
End at $0\in C\cap D_0$ and build the path backwards using
the table: for $b$ with source label $D_0$ use its left rule;
with source label $C$ use the left rule for a target in $P_1$
and the right rule for a target in $P_0$. The checked memberships
justify every preceding target requirement. Forward execution by
the same greedy state rule now realizes the prescribed even word,
ends at zero, and continues safely forever. Different block strings
give different words because the blocks are a prefix code. Thus
$N_m(s_g)\ge H_m$.

Let $c_* =\min_{0\le j\le5}H_j/\rho^j>0$. The full recurrence
for $m\ge6$ proves $H_m\ge c_*\rho^m$. This and
\eqref{eq:greedy-upper} yield
\begin{equation}\label{eq:greedy-exact}
 h_{A_2}(s_g)=\log\rho.
\end{equation}
All inverse paths above avoid $I$. Therefore the same lower
construction applies when an arbitrary Borel subset
$E\subset[12/5,162/65)$ is changed to control $B1$, with all
other states greedy. Such a feedback is globally legal and Borel,
but has $h_{A_2}\ge\log\rho$; this entire short-reset template
cannot lower the greedy value. This is not a lower bound for the
full four-control class or for arbitrary two-$A$ Borel rules.

\section{Arithmetic sampling and finite-block optimization}
\label{app:finiteblock}
For a fixed globally safe Borel rule and arithmetic sampling
$(0,p,2p,\ldots)$, a length-$n+m$ sampled pattern is determined
by its first $n$ entries and the length-$m$ pattern starting from
$T_s^{pn}x\in Q$. Consequently $N_{n+m}(s)\le N_n(s)N_m(s)$.
The logarithmic rate is a limit and equals
$\inf_n n^{-1}\log N_n(s)$ (with division by $\tau$ when
the control period is $\tau$). This uses forward safety, not
surjectivity or unrestricted concatenation.

For the full four-control class at $A_2$, let $K_n$ be the least
actual count $N_n(s)$ among unrefined legal Borel feedbacks.
It exists as a least member of a nonempty subset of
$\{1,\ldots,4^n\}$. Merging refinements cannot increase counts,
and unrefined rules are themselves allowed, so
\begin{equation}\label{eq:finiteblock}
 \ell_2=\inf_s\inf_n\frac{\log N_n(s)}n
 =\inf_n\frac{\log K_n}n.
\end{equation}
Only two infima are exchanged. No common feedback is asserted to
attain every $K_n$, no minimizer of the asymptotic infimum is
deduced, and no supremum over sampling sequences is interchanged.
This formula turns a \emph{proved real-feedback} finite-block
upper bound into an asymptotic upper bound. An open-loop word
table or a finite enumeration without global implementation is
insufficient.


\subsection{Weighted actual pattern pairs}
The volume argument in Section~\ref{subsec:four-affine} has a
more informative form that retains the compatibility of even
and odd patterns. Throughout this subsection the feedback is
an unrefined legal Borel rule for the original four-branch system.
Set $a=3/2$, $b=5/3$, and $r_1=a/b=9/10$. Let
\[
 V_n=\{(s(T_s^{2j}x))_{0\le j<n}:x\in Q\},\qquad N_n=|V_n|,
\]
and let $I_n$ consist of the pairs $(\alpha,\gamma)$ of even
and odd patterns from the \emph{same} initial state. Since
$T_sx\in Q$, one has $I_n\subset V_n\times V_n$; surjectivity
of $T_s$ is unnecessary. Let $B(\alpha)$ count the $B0,B1$
labels of $\alpha$, and define
\[
 W_n=\sum_{\alpha\in V_n}r_1^{B(\alpha)},\qquad
 Z_n=\sum_{(\alpha,\gamma)\in I_n}
            r_1^{B(\alpha)+B(\gamma)},\qquad
 \kappa_n=Z_n/W_n^2.
\]
In particular, $0<\kappa_n\le1$ and $W_n\le N_n$.

\begin{proposition}\label{prop:weighted-pairs}
With the preceding definitions, put
\[
 \theta_n=-\frac1n\log\frac{W_n}{N_n},\qquad
 d_n=-\frac1{2n}\log\kappa_n.
\]
Then $\theta_n,d_n\ge0$ and
\begin{equation}\label{eq:weighted-pairs}
 \frac{\log N_n}{n}\ge\log a+\theta_n+d_n.
\end{equation}
More precisely, if $\Lambda_w$ is the slope of the complete
affine word $w$ and
\[
 \Omega_m=\sum_{w\text{ actual},\ |w|=m}\Lambda_w^{-1},
 \qquad e_n=\frac{\log\Omega_{2n}}{2n},
\]
then $\Omega_m\ge1$, $e_n\ge0$, and
\begin{equation}\label{eq:weighted-pairs-exact}
 \frac{\log N_n}{n}-\log a=\theta_n+d_n+e_n.
\end{equation}
\end{proposition}
\begin{proof}
For each actual complete word $w$, its Borel initial fiber
$E_w$ has diameter at most $6/\Lambda_w$, because two states
following the same word separate by the factor $\Lambda_w$
and their final states lie in $[0,6]$. All these fibers partition
$Q$. Consequently
\[
 1\le\sum_{w\text{ actual},\ |w|=2n}\Lambda_w^{-1}
       =a^{-2n}Z_n.
\]
The equality uses the bijection between complete words and their
actual interlaced pairs. Nonempty zero-measure fibers are included.
Substitution of $Z_n=\kappa_nW_n^2$ proves
\eqref{eq:weighted-pairs}; the same substitution with
$\Omega_{2n}=a^{-2n}Z_n$ proves
\eqref{eq:weighted-pairs-exact}.

Each original inverse branch carries $Q$ into its safety domain.
Thus the complete inverse interval $J_w=G_w^{-1}(Q)$ is contained
in $Q$ and has length $6/\Lambda_w$; every intermediate open-loop
state is safe. The real fiber need only satisfy $E_w\subset J_w$.
The intervals $J_w$ for actual words cover $Q$, and, writing
$\nu_m=\sum_w\one_{J_w}$, one obtains
\[
 6(\Omega_m-1)=\int_Q(\nu_m-1)\,dx
              =\sum_w\bigl(|J_w|-|E_w|\bigr)\ge0.
\]
Here $|E_w|$ denotes Lebesgue measure; no interval or convexity
assumption is made on a feedback fiber. This also proves the
claimed nonnegativity at every length.
\end{proof}

The three terms distinguish fast-branch weighting, absent actual
pattern pairs, and excess complete inverse coverage. They are
support-counting quantities, rather than empirical frequencies
or Shannon entropies. The estimates transfer to arbitrary atom
refinements by projecting to their physical control labels;
the equalities themselves refer to those merged labels.
Splitting an actual trajectory at time $2n$ also gives
\[
 N_{n+m}\le N_nN_m,\qquad W_{n+m}\le W_nW_m,
 \qquad Z_{n+m}\le Z_nZ_m.
\]
The weights multiply under this split, and the suffix is realized
from a state in $Q$. The subadditivity argument at the beginning
of this appendix therefore gives finite limits
\[
 h=\lim_n\frac{\log N_n}{n},\quad
 p=\lim_n\frac{\log W_n}{n},\quad
 t=\lim_n\frac{\log Z_n}{2n},\qquad
 \log a\le t\le p\le h,
\]
and the exact limiting decomposition
$h-\log a=(h-p)+(p-t)+(t-\log a)$.

These formulas do not supply a uniformly positive correction.
Even the uniformly strict one-step inequality
$\Omega_1\ge6/5$ does not iterate: controls at $0$ and $6$
are different and each inverse slope is at least $3/5$, which
proves that inequality, but the greedy feedback has
$\Omega_1=4/3$ and
\[
 1\le\Omega_m(s_g)=a^{-m}D_m(s_g)\le10.
\]
The last bound is the all-initial-state word estimate in
Section~\ref{subsec:four-affine}. Hence $e_n(s_g)\to0$.
The complete inverse intervals overcount the true fibers at one
step without producing an exponential overcount at later steps.

\begin{example}[Full support from a countable modification]
\label{ex:periodic-support}
There is a globally safe two-$A$ Borel feedback for which
$N_n=2^n$, $I_n=V_n\times V_n$, and $\theta_n=d_n=0$ at
every length. It agrees with the greedy feedback on the full
infinite trajectories of all but countably many initial states.
\end{example}
\begin{proof}
For a purely periodic binary sequence $d$, set
\[
 \pi(d)=3\sum_{j\ge0}d_ja^{-(j+1)}.
\]
This address is injective on all purely periodic sequences.
Indeed, if two different sequences share a common period $L$
and the same address, subtracting their periodic-sum formulas
gives a nonzero integer polynomial $P$ with leading coefficient
$\pm1$ and $P(3/2)=0$. If its degree is $m$, multiplication by
$2^m$ leaves an odd leading term $\pm3^m$ and only even
remaining terms, a contradiction. Degree zero is immediate.

Let $P_i$ be the addresses of the periodic sequences with first
digit $i$, and $P=P_0\cup P_1$. These sets are countable Borel
sets, and $P_0\cap P_1=\varnothing$. Define $s_{\mathrm{per}}=Ai$
on $P_i$ and $s_{\mathrm{per}}=s_g$ on $Q\setminus P$. The identity
\[
 a\pi(d)-3d_0=\pi(\sigma d)\in P\subset Q
\]
proves safety on $P$, including $0$ and $6$; safety elsewhere
is that of $s_g$. This one state rule is therefore safe on all
infinite trajectories. Injectivity and induction show that its
actual name from $\pi(d)$ is precisely $d$.

Any prescribed binary pattern on a nonempty finite window extends to a
purely periodic sequence: choose a period larger than its largest
coordinate and fill the other digits arbitrarily. The empty window has
one pattern. Thus
$\Patt_{s_{\mathrm{per}}}(F)=2^{|F|}$ for every finite $F$.
Applying this to interlaced even and odd patterns gives
$V_n=\{A0,A1\}^n$, $I_n=V_n^2$, $W_n=2^n$, $Z_n=4^n$,
and the asserted zero corrections.

Finally $R=\bigcup_{j\ge0}T_g^{-j}P$ is countable: each point
has at most two one-step greedy preimages. Outside $R$ the
greedy trajectory never enters $P$, so induction gives identical
states and controls under the two rules at every time. Nevertheless,
their consecutive-time counting entropies are respectively
$\log(3/2)$, by Proposition~\ref{prop:density-one}, and $\log2$.
Discarding countably many initial states would therefore change
the counting question. No realization of all infinite nonperiodic
binary names is needed or claimed.
\end{proof}

This example excludes an unconditional positive lower bound for
$\theta_n+d_n$ over the whole feedback class. Its own pattern
count is already large, so it does not exclude an improved estimate
restricted to feedbacks with small $N_n$, and it supplies no
low-entropy feedback or strict full-class minimax gap.


\section{Itinerary closures and feedback consistency}
\label{app:language}

We construct a closed interval-covering language with smaller even
entropy than the greedy feedback. The construction also identifies the
additional consistency condition needed for a feedback realization.
For an arbitrary two-$A$ rule $s_E$ in
Section~\ref{subsec:four-affine}, write its real binary name as
$B_E(x)_j=\one_{\{s_E(T_E^jx)=A1\}}$ and put
\[
 Y_E=\cl B_E(Q),\qquad
 \pi(d)=2\sum_{j\ge0}d_j(2/3)^j,
 \qquad P_2(d)=(d_0,d_2,d_4,\ldots).
\]
Every coordinate of $B_E$ is Borel, and the original state rule
gives $\sigma B_E(x)=B_E(T_Ex)$. Thus $Y_E$ is nonempty,
compact, and $\sigma Y_E\subset Y_E$, without continuity of
the closed loop. The true affine recurrence gives
\[
 x=2\sum_{j=0}^{n-1}B_E(x)_j(2/3)^j+(2/3)^nT_E^nx.
\]
The remainder is at most $6(2/3)^n$, so $\pi(B_E(x))=x$.
Continuity of $\pi$ and the range $[0,6]$ imply
$\pi(Y_E)=Q$.

A cylinder fixing the first $n$ even digits is open. If it meets
the closure $Y_E$, it meets $B_E(Q)$; therefore closure adds
no finite even patterns. Since $\sigma P_2=P_2\sigma^2$,
$P_2(Y_E)$ is a one-sided subshift with exactly these real
prefixes. Its prefix entropy is its topological entropy, by the
same finite-cylinder refinement argument used in
Theorem~\ref{thm:C}. Consequently
\begin{equation}\label{eq:closure-even}
 h_{\rm top}(P_2(Y_E))=h_{A_2}(s_E).
\end{equation}
No factor $1/2$ is inserted, because the rate is per sampled
symbol at control period one.

\subsection{A closed shift-invariant interval-covering language}
Let
\[
 W_{10}=\{(a,b,a,c,d,e,f,g,h,g):a,\ldots,h\in\{0,1\}\},
\]
let $Y_0$ be arbitrary infinite concatenations of these ten-digit
blocks, and set $Y=\bigcup_{i=0}^9\sigma^iY_0$.
The equality conditions $d_{10j}=d_{10j+2}$ and
$d_{10j+7}=d_{10j+9}$ make $Y_0$ closed.
Its compact shift images have a finite closed union; adding or
removing one full block gives $\sigma^{10}Y_0=Y_0$ and
$\sigma Y=Y$.

\begin{proposition}\label{prop:abstract-low}
The language $Y$ satisfies
\[
 \pi(Y)=Q,\qquad
 h_{\rm top}(P_2(Y))=\frac45\log2<\log\rho.
\]
It therefore disproves the assertion that closedness,
$\sigma Y\subset Y$, and full interval projection alone force
the greedy even entropy. It does not provide an implemented
Borel feedback or improve $\ell_2^A$.
\end{proposition}
\begin{proof}
Put $r_0=2/3$, $\vartheta=r_0^{10}=1024/59049$,
$\Delta=6\vartheta=2048/19683$.
The block projection $t_w=2\sum_{j=0}^9w_jr_0^j$ is a
subset sum of the eight weights
\[
 26/9,\quad4/3,\quad16/27,\quad32/81,\quad64/243,
 \quad128/729,\quad3328/19683,\quad512/6561.
\]
In increasing order their numerators over denominator $19683$ are
\[
 1536,3328,3456,5184,7776,11664,26244,56862.
\]
The corresponding values of
$19683(\Delta+\sum_{j<i}\kappa_j-\kappa_i)$ are
\[
 512,256,3456,5184,7776,11664,8748,4374,
\]
all positive.
Thus each new weight is at most the earlier total plus $\Delta$.
Inductively, the intervals $[t,t+\Delta]$ over all subset sums
fill the single interval $[0,\Delta+\sum_i\kappa_i]=[0,6]$:
adding a weight unites the preceding interval with its overlapping
translate. Hence
\[
 Q=\bigcup_{w\in W_{10}}(t_w+\vartheta Q).
\]
Given $x_0\in Q$, choose a containing interval and write
$x_0=t_{w_0}+\vartheta x_1$, $x_1\in Q$; repeat forever.
Then
$x_0=\sum_{j<n}\vartheta^jt_{w_j}+\vartheta^nx_n$,
with remainder at most $6\vartheta^n$. The infinite block
concatenation belongs to $Y_0$ and projects to $x_0$.
This proves genuine infinite coverage, including all endpoints,
not just finite-depth coverage. The reverse inclusion is automatic
for binary sequences.

The even and odd block projections of $Y_0$ are
$(a,a,d,f,h)$ and $(b,c,e,g,g)$: each has exactly four free
bits among five output positions. For a shift of a block
concatenation, an even prefix consists of complete five-position
outputs and at most two partial boundary blocks. There is a
constant $C$, independent of the length, with prefix count at
most $10\cdot2^{4n/5+C}$. On lengths $n=5m$, $P_2(Y_0)$
has $2^{4m}$ distinct prefixes. Since $P_2(Y)$ is a subshift,
its entropy limit is exactly $(4/5)\log2$.
The strict comparison with $\rho$ can be made without decimals:
integer arithmetic gives
$\rho>1743/1000>2^{4/5}$ by substituting $1743/1000$ in
the decreasing reciprocal equation for $\rho$, and checking
$1743^5>16\cdot1000^5$.
\end{proof}

The coding constructed in this proof is open loop. Its successive
tails are safe states, since
$\pi(\sigma d)=(3/2)\pi(d)-3d_0$, but different codings
of the same state may require incompatible first controls.
The exact unresolved implementation question is whether there is
a Borel $B:Q\to Y$ with, for \emph{every} $x$,
\begin{equation}\label{eq:implementation}
 \pi(B(x))=x,\qquad
 B((3/2)x-3B(x)_0)=\sigma B(x).
\end{equation}
A Borel right inverse of $\pi$ alone does not meet this requirement.
If both equations were proved, the real rule would have
$h_{A_2}\le(4/5)\log2$ and would disprove greedy optimality.
Neither that implementation nor its global impossibility is asserted.


\begin{example}[A safe periodic seed need not extend]
\label{ex:periodic-seed}
The word $w=1010001101$ belongs to $W_{10}$. Let
$\omega=w^\infty$ and $p_j=\pi(\sigma^j\omega)$, with indices
modulo ten. Direct evaluation gives
\[
 (p_0,\ldots,p_9)=\frac1{5275}
 (17358,10212,15318,7152,10728,16092,24138,20382,14748,22122).
\]
The points are distinct and lie in $Q$. On them prescribe the
control $A_{w_j}$, prescribe $A0$ at $0$, $A1$ at $6$, and use
the greedy rule elsewhere. This is a globally safe Borel rule:
$g_{A_{w_j}}(p_j)=p_{j+1}$ follows from the projection identity,
and the other states use already safe controls. Its names on
this finite invariant seed all lie in $Y$.

There is, however, no implementation of $Y$ extending these
prescribed names. The point
$x_*=2p_1/3=6808/5275<2$ has only $A0$ legal and its next
state is $p_1$. Consistency would force its name to be
$z=0\sigma\omega$. The periodic sequence
$e=\sigma\omega=(0100011011)^\infty$ satisfies
$e_r=e_{r+2}$ exactly at residues $0,2,6,9$ modulo ten.
The phase test below therefore puts $e$ in exactly phases $1,8$.
If $z\in Y$ its phase must consequently be $0$ or $7$, because
deleting its first digit adds one to the phase. Both phases require
$z_0=z_2$, whereas these digits are $0$ and $1$.
Thus $z\notin Y$. This excludes extension of this particular
seed, not every Borel implementation of $Y$.
\end{example}


\subsection{The exact forced-predecessor compatibility kernel}
For $a\in\{0,1\}$ write $\kappa_a(d)=ad$ and define
\[
 Z_0=Y,\qquad Z_{n+1}=\{d\in Z_n:\sigma d\in Z_n,
 \ \pi(d)<3\Rightarrow0d\in Z_n,
 \ \pi(d)>3\Rightarrow1d\in Z_n\}.
\]
These are closed and nested. For example the excluded condition
$\pi(d)<3$ together with $0d\notin Z_n$ is open; the other
strict-threshold exclusion is likewise open.
Every implementation in \eqref{eq:implementation} has
$B(Q)\subset Z_n$ for every $n$. Indeed, when $x<3$,
its original predecessor $\psi_0(x)=2x/3<2$ has only control
$0$ legal, so $B(\psi_0x)=0B(x)$. When $x>3$,
$\psi_1(x)=2+2x/3>4$ has only $1$ legal and
$B(\psi_1x)=1B(x)$. This, forward consistency, and induction
give the claim. At $x=3$, its predecessors $2,4$ are in
overlap safety domains and supply neither compulsory condition.


The first deletion layer still covers the entire state interval:
\[
 \pi(Z_1)=Q.
\]
For $x<3$, choose $z\in Y$ with $\pi(z)=2x/3<2$, using the
full projection of $Y_0$ proved above. Its first digit must be
zero, since a first digit one gives projection at least two.
Writing $z=0d$ gives $d\in Y$, $\pi(d)=x$, $\sigma d\in Y$,
and $0d=z\in Y$, so $d\in Z_1$. For $x>3$ choose instead
$\pi(z)=2+2x/3>4$. Its first digit must be one, since a first
digit zero gives projection at most four; $z=1d$ gives the
required coding in $Z_1$. At $x=3$ any $Y$ coding suffices,
because neither forced-prefix condition applies. This finite-layer
coverage leaves the intersection of all layers to be tested.


The limit kernel has the following exact form:
\begin{equation}\label{eq:kernel-H}
 H=\left\{d\in Y:\ \forall j\ge0,
 \begin{array}{l}
 \pi(\sigma^jd)<3\Rightarrow00\sigma^jd\in Y,\\
 \pi(\sigma^jd)>3\Rightarrow11\sigma^jd\in Y
 \end{array}\right\}=\bigcap_nZ_n.
\end{equation}
Here each tested member may use its own phase in the union $Y$.
For the equality, we need the prefix implication
$aa d\in Y\Rightarrow a^m d\in Y$ for every $m\ge2$.
A sequence $z$ belongs to $\sigma^iY_0$, $0\le i<10$, exactly when
\[
 z_r=z_{r+2}\quad\text{whenever }r\ge0
       \text{ and }r+i\equiv0\text{ or }7\pmod {10}.
\]
The missing coordinates before zero can be supplied independently,
since the prescribed copying pairs in each ten-block are disjoint.
Suppose $z=aa d$ satisfies these conditions in phase $i$.
The sequence $az$ satisfies them in phase $i-1$ modulo ten:
every pair starting at $r\ge1$ is an old pair translated by one,
and a possible pair starting at zero compares the equal letters
$(az)_0=(az)_2=a$. Iteration proves the prefix implication.

The set $H$ is closed, is forward shift-invariant, and is closed
under $d\mapsto0d$ when $\pi(d)<3$ and under $d\mapsto1d$
when $\pi(d)>3$. For example $00d\in Y$ gives $0d\in Y$;
the new zero-tail test requires $000d\in Y$, supplied by the
prefix fact, while the later tests are those of $d$. Thus
$H\subset Z_n$ inductively. Conversely, the nested intersection
is closed under shift and these forced prefixes: membership of
$d$ in every $Z_{n+1}$ places the required image in every $Z_n$.
For $\pi(d)<3$, $\pi(0d)=2\pi(d)/3<3$, so two such
prefixes belong to the intersection, hence to $Y$; the same
argument holds above three. Applying it to every shifted tail
proves the reverse inclusion in \eqref{eq:kernel-H}.

For each $x$, the fibers $Z_n\cap\pi^{-1}\{x\}$ are nested
compact sets. Therefore
\begin{equation}\label{eq:kernel-compact}
 x\in\pi(H)\quad\Longleftrightarrow\quad
 Z_n\cap\pi^{-1}\{x\}\ne\varnothing\text{ for every }n.
\end{equation}
This is an exact compactness interface, not a proof that the fibers
are nonempty.

There is also the exact central reduction
\begin{equation}\label{eq:kernel-central}
 [2,3]\subset\pi(H)\quad\Longleftrightarrow\quad\pi(H)=[0,6].
\end{equation}
Complementation $d\mapsto1-d$ preserves $Y$ and $H$ and
sends the projection to $6-\pi(d)$. If $0<x<2$, take the
least $k$ with $y=(3/2)^kx\ge2$; then $2\le y<3$.
Starting with an $H$ coding of $y$, repeated forced zero
prefixes give a coding of $x$. The endpoints zero and six are
the constant names, and $2,3,4$ are supplied by the central
interval and complementation. This proves the nontrivial direction;
the other is immediate. No condition is added at projection three.

Neither \eqref{eq:kernel-compact} nor
\eqref{eq:kernel-central} decides the remaining coverage of
$[2,3]$. Even full coverage of $H$ would not by itself give a
single-step-consistent Borel section. A missing fiber would exclude
this language's global implementation but would not prove greedy
optimality or exclude other lower-entropy feedbacks. Accordingly
the bounds \eqref{eq:four-L} and \eqref{eq:four-even} are unchanged.

\section{A bounded-density obstruction to simultaneous scalar capacities}
\label{app:measure}
The smooth branches in Section~\ref{subsec:smooth} cannot all be
scalarized by an equivalent probability with a density bounded both
above and away from zero. The proof treats every such Borel density,
without inserting it at a null periodic point.

\begin{proposition}\label{prop:bounded-density}
There are no Borel $a$ and positive $c_0,c_1,c_2$ with
$0<\operatorname*{ess\,inf}a\le\operatorname*{ess\,sup}a<\infty$,
\[
 \rho_u(y)a(\psi_u(y))\le c_u a(y)\quad\text{almost everywhere},
 \qquad\max_{u\ne v}(c_u+c_v)<1.
\]
Here all three branches use the same density.
\end{proposition}
\begin{proof}
Choose $0<m\le M<\infty$ bounding $a$ almost everywhere, and
remove the null exceptions for these bounds and the three inequalities.
Also remove their preimages under every finite composition of the
$\psi_u$. These are null because every composition has a Lipschitz
inverse on its image: its derivative is bounded below by
$(2/15)^n$. The remaining Borel set $G$ has full measure and is
preserved by every inverse branch. For $y_0\in G$ and
$y_{j+1}=\psi_{u_j}(y_j)$, multiplication gives
\begin{equation}\label{eq:density-iterate}
 \left(\prod_{j<n}\rho_{u_j}(y_j)\right)a(y_n)
 \le\left(\prod_{j<n}c_{u_j}\right)a(y_0).
\end{equation}
Since $\psi_0(0)=0$, $\psi_0(y)\le3y/5$, and
$\rho_0(0)=3/5$, every $\eta>0$ has a small invariant interval
on which $\rho_0\ge3/5-\eta$. Choose a point of $G$ in its
interior, iterate, and take $n$th roots in
$m(3/5-\eta)^n\le Mc_0^n$. Letting $n$ grow and then
$\eta$ decrease proves $c_0\ge3/5$.

For a second constraint use the positive-length closed intervals
\begin{align*}
 I_0&=[9/10,23/25],& I_1&=[3/10,33/100],\\
 I_2&=[9/20,47/100],& I_3&=[53/100,14/25],\\
 I_4&=[19/25,79/100],& I_5&=I_0.
\end{align*}
The inverse word is $(0,1,1,2,2)$. The following bounds are
rational containing intervals for the \emph{complete} images:
\begin{center}
\begin{tabular}{ccc}
\toprule
Input and branch&Containing interval&Next interval\\
\midrule
$\psi_0(I_0)$&$[3081/10000,3195/10000]$&$\intr I_1$\\
$\psi_1(I_1)$&$[4510/10000,4692/10000]$&$\intr I_2$\\
$\psi_1(I_2)$&$[5385/10000,5494/10000]$&$\intr I_3$\\
$\psi_2(I_3)$&$[7672/10000,7836/10000]$&$\intr I_4$\\
$\psi_2(I_4)$&$[9002/10000,9170/10000]$&$\intr I_0$\\
\bottomrule
\end{tabular}
\end{center}
For an explicit verification of every endpoint bound, use
$314159/100000<\pi<314160/100000$ and
$173205/100000<\sqrt3<173206/100000$.
The latter bounds follow by squaring; the former follow from
Machin's identity with the alternating arctangent sums through
degrees $11$ and $3$ at $1/5$ and $1/239$.
At the two control-zero endpoints use $z=2\pi(y-1)$, and at
the others use $z=2\pi(y-u/3)$. In every case $|z|<1$.
Substitute the rational interval for $z$ in
\[
 S_7(z)=z-z^3/6+z^5/120-z^7/5040,\qquad
 |\sin z-S_7(z)|\le1/362880.
\]
The constant sine terms are $0,-\sqrt3/2,\sqrt3/2$, and
$7/(60\pi)$ is bounded by the reciprocal rational interval.
Interval addition and multiplication at each endpoint lie strictly
inside the displayed containing bounds. Monotonicity of $\psi_u$
then proves the full-image inclusions, including endpoints.

On each input interval the distance to its corresponding cosine
center ($1,1/3,2/3$) is at most $41/300$. Thus
\[
 \rho_{u_j}(y)\ge\frac35-\frac7{60}(41\pi/150)^2
 >\frac{1214099}{2362500}>\frac12.
\]
The composite $\psi_2\psi_2\psi_1\psi_1\psi_0$ preserves
$I_0$. Starting at any $y\in G\cap I_0$, repeat its five-step
word $n$ times in \eqref{eq:density-iterate}. This yields
$m(1/32)^n\le M(c_0c_1^2c_2^2)^n$, and hence
$c_0c_1^2c_2^2\ge1/32$. But the pairwise sums and
$3/5\le c_0<1$ give $c_1,c_2<1-c_0$. The function
$t(1-t)^4$ decreases on $[3/5,1)$, so
\[
 c_0c_1^2c_2^2<\frac35(2/5)^4=\frac{48}{3125}<\frac1{32},
\]
a contradiction. All iterations occur in the same full-measure
set; no regularity of $a$ was assumed.
\end{proof}

For $d\mu_a=a\,dy/\int a$, the density of the branch measure
relative to $\mu_a$ is exactly
\[
 \rho_u^a(y)=\rho_u(y)a(\psi_u(y))/a(y),
\]
by the change of variables
$x=\psi_u(y)$. The proposition excludes this bounded positive
density class. It does not exclude arbitrary full-support singular
probabilities or densities without uniform bounds; that question
remains open for this smooth example. It does not alter its already
proved value $M_1=\log3$.

There is an exact reduction of that remaining question, without a
density assumption. If a full-support probability $\mu$ satisfied
$\mu(\psi_u(E))\le c_u\mu(E)$ for all Borel sets and the
same target pair sums, then $q_0=\max c_u<1$ would force it to
be nonatomic. Indeed, an atom of mass $m>0$ can be followed along
any successive legal controls, since the safety intervals cover
$Q$. Applying the set inequality to the successor singleton raises
its mass to at least $m/q_0$ at each step, eventually exceeding
one. Repetition of states does not avoid that contradiction.

Its distribution function $h(x)=\mu([0,x])$ would therefore be a
continuous strictly increasing homeomorphism of $[0,1]$, sending
$\mu$ to Lebesgue probability. Applying the inequalities to
$(x,y]$ shows that every $h\psi_uh^{-1}$ is $c_u$-Lipschitz.
Conversely, any one increasing homeomorphism with these simultaneous
Lipschitz bounds defines a full-support nonatomic measure
$\mu=(h^{-1})_*\Leb$ with the required inequalities on every
Borel set. To see the last assertion, cover a Borel set in the
$h$ coordinate by an open set of arbitrarily close measure.
Decompose that open set into countably many intervals; each image
has length at most $c_u$ times its length. Subadditivity and
outer regularity give the full Borel estimate, not merely an
interval estimate. This equivalence leaves existence of $h$
undecided and is not an extension of the preceding exclusion.

\section{Conditional real returns and closed viability boundaries}
\label{app:returns}

\subsection{A two-branch return condition}
The affine system below is the four-branch system of
Section~\ref{subsec:four-affine}. A sufficient certificate is stated
with its true fibers; extracting it for every Borel feedback is an
unresolved step.

\begin{proposition}\label{prop:return-certificate}
Suppose one legal feedback has distinct real control words $w^0,w^1$
of the same length $k$, and a nonempty open interval $I\subset(0,6)$,
such that its real prefix fiber $E_{w^i}$ is comeager in the complete
inverse interval $J_i=G_{w^i}^{-1}(I)$, and $J_i\subset I$.
For the single predetermined sequence $A=(j!)_{j\ge1}$ this
feedback has $h_A(s)\ge\log2$.
\end{proposition}
\begin{proof}
Each inverse $\psi_i:I\to J_i\subset I$ is an affine
homeomorphism, and its bad subset $J_i\setminus E_{w^i}$ is
meager. For a prescribed finite sequence of the two words, compose
their complete inverse branches backwards from $I$. The resulting
initial interval is nonempty and open; the inverse images of every
bad set along the finite affine maps are meager. A point outside
their finite union executes precisely the prescribed sequence by
the same feedback. It remains safe thereafter by global legality.

Choose an offset $0\le t<k$ where the two words differ. For
all $j\ge k$, $j!$ is divisible by $k$. Prescribe independently
the word in block $j!/k$ of a sufficiently long concatenation,
fill all other blocks arbitrarily, and start observation at the
actual state $T_s^t x$. Its control at time $j!$ is the offset-$t$
letter of that chosen block. Thus the first $n$ factorial samples
have at least $2^{n-O(k)}$ patterns. The ignored head is fixed
and the normalized limsup gives the claim. The word length $k$
is a proof device; the feedback period is still one.
\end{proof}

For any expanding affine word $G_w(x)=\lambda_w x-t_w$ and
$I=(\alpha,\beta)$, elementary endpoint inequalities give
\[
 G_w^{-1}(I)\subset I
 \quad\Longleftrightarrow\quad
 t_w/(\lambda_w-1)\in[\alpha,\beta].
\]
The interval must be the complete inverse of $I$; a locally feasible
subinterval is insufficient. The proposition would give
$L_1=M_1=\log2$ if its certificate were extracted for every
legal Borel rule, even with $k$ depending on the rule. Such an
extraction is not proved here. Compactness of the state interval
does not make the true return graph finite, and openness of each
individual affine branch does not imply Borel closed-loop stitching.


A common covered target need not admit the returns in
Proposition~\ref{prop:return-certificate}. To see the distinction
on real trajectories, use the globally safe greedy rule $s_g$,
whose interval $[0,3]$ is forward invariant. The target
$I=(2,5/2)$ has complete one-step inverse intervals
\[
 J_0=(4/3,5/3),\qquad J_1=(10/3,11/3),
\]
and the rule really uses $A0$ throughout $J_0$ and $A1$
throughout $J_1$. Starting in $I$, however, no trajectory can
ever reach $J_1$, since it remains in $[0,3]$. Thus this
covered target has no return to its second source at any length.
The same rule does have a return certificate on the different
target $I_0=(0,3)$: the words $A0,A0,A0$ and $A1,A0,A0$
have complete inverse intervals $(0,8/9)$ and $(2,26/9)$,
both inside $I_0$. Their intermediate states lie successively
in $(0,4/3)$, $(0,2)$, and $(0,3)$, where the indicated
controls are actually chosen. All points on these inverse
intervals follow the stated words and continue safely thereafter.
The failed return for the first target therefore concerns the
choice of certificate, rather than this feedback's sampled entropy
or the existence of every possible return certificate.


\subsection{Nested viable sets}
For a closed $K\subset Q=[0,6]$, let
\[
 \Gamma(K)=\{x\in Q:\exists\text{ a legal }u,\ g_u(x)\in K\},
 \quad c(K)=\sup\{\operatorname{diam}C:C\text{ a component of }K\}.
\]
For closed $C$, define $C^{(0)}=C$,
$C^{(n+1)}=C^{(n)}\cap\Gamma(C^{(n)})$, and
$V(C)=\bigcap_n C^{(n)}$. All these sets are compact.
Finite controls and nested compactness show $V(C)\subset\Gamma(V(C))$:
the nested nonempty sets of legal controls carrying a point into
$C^{(n)}$ have a common control. Conversely every closed viable
$K\subset C$ remains in every iterate and hence in $V(C)$.

Call a closed set $K$ viable if $K\subset\Gamma(K)$.
For $\delta>0$, say that $s$ has \emph{complete-branch growth at scale
$\delta$} if every nonempty relatively open $V\subset Q$ contains a
set
\[
 J=G_w^{-1}(B_Q(y,\delta))\subset V
\]
for some real control word $w$ and $y\in Q$, such that its actual
prefix fiber $E_w(s)$ is comeager in $J$. Here $G_w$ is the complete
affine composition, $B_Q(y,\delta)=Q\cap(y-\delta,y+\delta)$,
and the empty word is allowed with $G_w$ the identity. The inverse
set is required in its entirety, including any relative endpoint.

\begin{proposition}\label{prop:nested-selector}
Every decreasing sequence of closed viable sets $K_j$ has one
globally safe Borel selector preserving all $K_j$. If also
$\intr_QK_j\ne\varnothing$ and $c(K_j)\to0$, this selector
fails complete-branch growth at every positive fixed scale.
No sequence satisfying these additional interior and component-scale
conditions can have a finite intersection.
\end{proposition}
\begin{proof}
On $K_j\setminus K_{j+1}$ choose the first original control
carrying the state into $K_j$. On $K_\infty=\bigcap_jK_j$
choose the first control carrying it into all $K_j$; a control
exists because the finite nonempty admissible-control sets are
decreasing. Outside $K_1$ use any ordered safe selector.
The selection domains are Borel, the choices are pointwise
consistent, and the resulting rule preserves each $K_j$ for
every state, hence along infinite real trajectories.

Fix $\delta>0$, choose $j$ with
$c(K_j)<\min\{\delta,6\}$, and an open interval
$V\subset\intr_QK_j$. A complete actual inverse branch starting
in $V$ and ending onto a relative open ball $B_Q(y,\delta)$
would send a comeager subset into $K_j$. Its affine homeomorphism
makes the image dense in that ball. Closedness forces the entire
ball into $K_j$, but this connected ball has diameter at least
$\min\{\delta,6\}$, contradicting the component bound.
Length zero gives the same contradiction. This is precisely the
complete-branch growth obstruction, not an entropy counterexample.

For the last assertion, impose the additional conditions just stated.
Nested nonempty compactness makes $F=\bigcap K_j$ nonempty.
Suppose that $F$ is finite. Choose
a closed $\eta$-neighborhood of $F$ small enough that its
nearest center is unique and every actual safe step staying in it
has $g_u(p_j)=p_{j+1}$ for the corresponding centers. Such an
$\eta$ exists: for the finitely many nonzero distances
$|g_u(p)-q|$, choose $\eta$ smaller than their minimum divided
by $1+\max_u\lambda_u$, and also smaller than one third of each positive distance between
distinct centers. For a singleton $F$ there is no center-separation
restriction; if the first collection of distances is empty, its
restriction is omitted. Along a trajectory
remaining in this neighborhood,
$|x_{n}-p_n|\ge(3/2)^n|x_0-p_0|$, whereas the left side is
at most $\eta$. Therefore $x_0\in F$, and the full viability
kernel of the neighborhood is contained in $F$. Nested compactness
places some $K_j$ inside this neighborhood. Its viability then
places it in $F$, forcing eventual equality with $F$. It cannot
have nonempty interior, a contradiction.
\end{proof}

There is a further exact restriction: a nondegenerate closed interval
$[a,b]\subset[0,6]$ is viable if and only if $b-a\ge3$.
The allowable $A0$ source interval inside it is
$[a,2b/3]$, and its $A1$ source interval is
$[2+2a/3,b]$, with empty intervals discarded. The corresponding
$B0,B1$ source pieces are contained in these, because
$3b/5\le2b/3$ and $(3a+12)/5\ge2+2a/3$ for $a\le6$.
Thus they cover the interval exactly when
$2b/3\ge2+2a/3$, equivalently $b-a\ge3$.
Nonconvex viable sets need not preserve their separate components,
so this statement is not extended to their convex hulls.

An infinite nested viable chain with nonempty interior and vanishing
component scale has not been constructed or excluded. The conditional
growth obstruction above neither excludes all real return certificates
nor changes $L_1$. These unresolved steps are kept separate from
the complete maximum-pattern theorems.

\section{Attainment of the fixed-period infimum}\label{app:minimum}

Integer-logarithm quantization of maximal-pattern entropy is a known
combinatorial fact. For one-sided subshifts it is stated in
Huang and Ye \cite[Theorem~5.5]{HuangYe2009}. We record the finite-set
argument needed here, since it also settles attainment of the feedback
infimum. This argument does not address the common-sampling problem.

\begin{proposition}\label{prop:quantization}
Let $\varnothing\ne\Omega\subset E^{\Nzero}$, where $|E|=q<\infty$,
and fix $\tau>0$. There is an integer $m\in\{1,\ldots,q\}$ such that
$\hpat(\Omega;\tau)=\log m/\tau$.
\end{proposition}
\begin{proof}
For $H\subset E^n$, let $d_k(H)$ be the largest size of a coordinate
set $J$ for which $H|_J$ contains $\prod_{j\in J}V_j$, with
$V_j\subset E$ and $|V_j|=k$. Put $d_k(\varnothing)=-1$.
The elementary finite bound is
\[
 |H|\le B_k(n,d):=
 \sum_{i=0}^{\min(d,n)}\binom ni(k-1)^{n-i}\binom qk^i
 \quad\text{if }d_k(H)\le d,\qquad 2\le k\le q.
\]
Set $B_k(n,-1)=0$. To prove the bound, write $H_e$ for the
$(n-1)$-letter prefixes ending in $e$, $V=\bigcup_eH_e$, and
$V_S=\bigcap_{e\in S}H_e$ for $|S|=k$. Projection gives
$d_k(V)\le d$, whereas $d_k(V_S)\le d-1$: a shattered set in
$V_S$ can be extended by the last coordinate. A prefix appearing
in $t$ of the sets $H_e$ contributes
$t\le(k-1)+\binom tk$. Hence
\[
 |H|\le(k-1)|V|+\sum_{|S|=k}|V_S|
 \le(k-1)B_k(n-1,d)+\binom qk B_k(n-1,d-1)
 =B_k(n,d).
\]
The empty class and the nonempty class at $n=0$ start the induction.

Let $D_k=\sup_Fd_k(\Omega|_F)$, over finite $F\subset\Nzero$,
and let $m$ be the largest $k$ with $D_k=\infty$.
Since $\Omega$ is nonempty, $D_1=\infty$.
Restricting shattered sets gives $p_\Omega^*(n)\ge m^n$ for every
$n$. If $m=q$, the alphabet bound gives equality. Otherwise
$D_{m+1}=d<\infty$, so the preceding bound gives, for a constant $C$,
\[
 m^n\le p_\Omega^*(n)
 \le\sum_{i=0}^{\min(d,n)}\binom ni m^{n-i}\binom q{m+1}^i
 \le C n^d m^n\qquad(n\ge1).
\]
Taking logarithms and dividing by $n\tau$ proves the assertion.
\end{proof}

\begin{corollary}\label{cor:pattern-minimum}
For every nonempty class of finite Borel invariant partitions at
one fixed period $\tau$, the infimum of $\hpat(\C;Q)$ is a minimum.
At the common attaining schedule of Theorem~\ref{thm:B}, the inner
sequence-entropy infimum is also a minimum. In particular, for a
finite control alphabet the quantity $M_\tau$ in
\eqref{cc-eq:LM} is attained without a countability assumption.
\end{corollary}
\begin{proof}
Apply Proposition~\ref{prop:quantization} to each $\Sigma_\C$.
The corresponding integers form a nonempty subset of $\N$, which
has a least member. At the schedule of Theorem~\ref{thm:B}, all
individual sequence rates equal their pattern rates.
For the full finite-control class, Lemma~\ref{cc-lem:merge}
identifies the infimum with that over legal selectors, whose
nonempty fibers have at most $|U|^\tau$ labels.
\end{proof}

No attainment assertion for an arbitrary fixed sequence $A$ follows
from this corollary. The integer quantization applies to maximal-pattern
entropy, and fixed-sequence entropy need not have that form.

\section{The obstruction to pointwise common attainment}
\label{app:pointwise}

Theorem~\ref{thm:B} attains the rate of every member of a countable
language family. The next constructions show why this stronger
conclusion must be distinguished from attaining the minimum rate
over all legal feedbacks. We first construct a compact family of
genuine subshifts, then use it in continuous control systems.

\subsection{A continuous family of sparse hereditary subshifts}

Write $\Theta=\{0,1\}^{\N}$, with its fair Bernoulli probability
$\nu$. For a finite binary word $v$, including the empty word, set
\[
 j(v)=2^{|v|}+\operatorname{val}_2(v),\qquad
 b_v=4^{j(v)},\qquad
 B_\theta=\{b_{\theta|_\ell}:\ell\ge0\}.
\]
In $K=\{0,1\}^{\Nzero}$ define
\[
 Y_\theta=
 \operatorname{cl}\bigcup_{t\in\mathbb Z}
 \{y:\supp(y)\subset(B_\theta+t)\cap\Nzero\}.
\]

\begin{lemma}\label{lem:sparse-bundle}
Each $Y_\theta$ is a nonempty hereditary subshift with
$\sigma Y_\theta=Y_\theta$ and zero ordinary entropy.
The map $\theta\mapsto Y_\theta$ is Hausdorff continuous and injective.
Every $Y_\theta$ has coordinate maximal-pattern entropy $\log2$.
For each fixed sampling sequence $A=(a_i)$,
\[
 \frac1n\log\Patt_{Y_\theta}(\{a_1,\ldots,a_n\})
 \longrightarrow0
 \quad\text{for }\nu\text{-almost every }\theta.
\]
The exceptional set may depend on $A$.
\end{lemma}
\begin{proof}
The positive differences of $D=\{4^j:j\ge1\}$ have unique
representations. Indeed, the $4$-adic valuation of
$4^j-4^i$ with $j>i$ is $i$, which determines $i$ and then $j$.
A finite pattern appears in $Y_\theta$ exactly when its set of
one-positions is contained in a translate of $B_\theta$;
cylinders are clopen and prescribed zero-positions impose no
further restriction. If there are at least two one-positions,
their difference fixes the candidate nodes and the translation.
The remaining one-positions either give incompatible nodes, or
the appearance parameters form the cylinder determined by the
deepest node. Patterns with at most one one appear for every
parameter. Thus each finite prefix language is locally constant
in $\theta$, proving Hausdorff continuity. The pattern $11$ on
$\{0,b_v-b_\varnothing\}$ detects whether $\theta$ extends a
nonempty word $v$, proving injectivity.

Setting coordinates to zero preserves the defining supports and
then their closure. Shifting replaces $t$ by $t-1$, while
prepending zero replaces it by $t+1$; compactness supplies
predecessors of limit points. This proves heredity and surjectivity.
An interval of length $N$ meets any translate of $D$ in at most
$L_N=2+\lceil\log_4(N+1)\rceil$ points. Consequently
\[
 |\mathcal L_N(Y_\theta)|
 \le (L_N+1)N^{L_N}\qquad(N\ge2),
\]
which gives zero ordinary entropy. On any $n$ positions in
$B_\theta$ all binary patterns occur, so $p_{Y_\theta}^*(n)=2^n$.

Fix $F_n=\{a_1,\ldots,a_n\}$ and put
$M_\theta(F_n)=\sup_t|F_n\cap(B_\theta+t)|$.
A translation having at least two intersections is determined
by a pair in $F_n$ and a difference in $D$, so there are at most
$\binom n2$ candidates, independent of $\theta$. Each candidate
involves at most $n$ nodes. A branch containing at least $k$
of these nodes contains one at depth at least $k-1$. The
parameter-cylinder probabilities therefore give
\[
 \nu\{M_\theta(F_n)\ge k\}\le n^3 2^{-(k-1)}
 \qquad(k\ge2).
\]
The events are finite unions of parameter cylinders. With
$k_n=\lceil6\log_2(n+1)\rceil+2$, the probabilities are summable.
Borel--Cantelli gives $M_\theta(F_n)<k_n$ eventually almost surely.
Every occurring pattern then has at most $k_n$ one-positions, and
\[
 \Patt_{Y_\theta}(F_n)\le(k_n+1)n^{k_n}\qquad(n\ge2).
\]
Its normalized logarithm tends to zero.
\end{proof}

\subsection{Compactness of a designated feedback family is insufficient}

\begin{proposition}\label{prop:compact-family}
There is a compact continuous control system with a compact
designated family of clopen period-one feedbacks $\{\C_\theta\}$
whose actual name languages are $Y_\theta$ and for which
\[
 \sup_A\inf_\theta\hseq^A(\C_\theta;Q)
 =0<\log2=\min_\theta\hpat(\C_\theta;Q).
\]
The full Borel and clopen classes in the same system have both
optimization values zero. The control alphabet in this construction
is uncountable.
\end{proposition}
\begin{proof}
Define $r_\theta:K\to Y_\theta$ by retaining the longest allowed
prefix of $x$ and replacing the rest by zeros; if every prefix
is allowed, set $r_\theta x=x$. Heredity makes this a retraction.
Both one-letter prefixes are allowed, so $(r_\theta x)_0=x_0$.
Every fixed output prefix depends on finitely many prefix tests,
whose parameter sets are clopen by Lemma~\ref{lem:sparse-bundle}.
Thus $(\theta,x)\mapsto r_\theta x$ is jointly continuous.

Take $X=K\sqcup\{\dagger\}$, $Q=K$ and
$U=\Theta\times\{0,1\}$, with $\dagger$ isolated, and define
\[
 f(x,(\theta,s))=
 \begin{cases}\sigma r_\theta x,&x_0=s,\\
 \dagger,&x_0\ne s,\end{cases}
 \qquad f(\dagger,u)=\dagger.
\]
The matching regions are clopen, so $f$ is jointly continuous.
Use the two fixed clopen atoms $\{x:x_0=s\}$ and assign
$(\theta,s)$ to them. Their actual names are exactly $Y_\theta$:
the first retraction preserves the initial symbol and thereafter
the state shifts inside $Y_\theta$. The assignments form a
compact family in $U^2$. Lemma~\ref{lem:sparse-bundle} gives the
displayed strict gap.

For the full class, let $\theta^0,\theta^1$ be the constant
parameters and use $(\theta^0,0)$ on the zero atom and
$(\theta^1,1)$ on the one atom. The positive difference sets of
$B_{\theta^0}$ and $B_{\theta^1}$ are disjoint. Hence retracting
a point of one language into the other leaves at most one one.
The languages have no adjacent ones. An initial zero trajectory
shifts in $Y_{\theta^0}$ until its first one, after which it
becomes zero. An initial one trajectory enters $Y_{\theta^1}$
with next symbol zero; the next retraction leaves at most one
further one. All actual and completed names have at most two ones.
Thus $p^*(n)\le1+n+\binom n2$ for this legal clopen feedback.
Nonnegativity gives zero for both full-class infima and for
every fixed-sequence inner infimum.
\end{proof}

\subsection{Four controls and a positive full-class minimum}

\begin{proposition}\label{prop:finite-pointwise}
There is a compact system with four distinct continuous control
maps and a controlled-invariant compact set $Q$ such that
\[
 \hinv(Q)=0,\qquad L_1=M_1=\log2,
\]
yet no sampling sequence attains $\hpat(\C;Q)$ for every legal
period-one Borel invariant partition $\C$.
This already fails for a family of four-label selectors with
Hausdorff continuous, compact name languages and rate $\log4$.
\end{proposition}
\begin{proof}
Let $Z$ be the Thue--Morse subshift of
Section~\ref{sec:thue-morse}, and let
\[
 \mathcal E=\{(\theta,y):y\in Y_\theta\},\qquad
 Q=Z\times\mathcal E\times\{0,1\},\qquad
 X=Q\sqcup\{\dagger\},\qquad U=\{0,1\}^2.
\]
Hausdorff continuity makes $\mathcal E$ compact.
For $x=(z,\theta,y,r)$ define
\[
 f(x,(a,b))=
 \begin{cases}
  (\sigma z,\theta,\sigma y,b),&a=z_0,\\
  \dagger,&a\ne z_0,
 \end{cases}
 \qquad f(\dagger,u)=\dagger.
\]
These maps are continuous on clopen matching regions. At each
state exactly the two controls with $a=z_0$ are safe. Their
successors differ in the register coordinate $b$, so the four
control maps are distinct and $Q$ is controlled invariant.
The register is a state coordinate of this system; the feedback
has no additional memory or time variable.

For every legal Borel partition, the first control label of
each atom must equal $z_0$. The first state coordinate always
shifts, independently of every Borel mixture and atom refinement.
Label projection therefore gives
$\Patt_\C(F)\ge\Patt_Z(F)$ for each finite $F$.
The two-atom clopen feedback $(z_0,0)$ attains equality for
every $F$. It follows that
\[
 \inf_{\C\in\mathfrak C_{\Bor}(Q,1)}\hseq^A(\C;Q)
 =h_{\rm seq,coord}^A(Z)\quad\text{for every }A,
 \qquad M_1=\log2.
\]
The schedule in Proposition~\ref{prop:thue-morse} gives
$L_1=\log2$. A control word of length $N$ is safe exactly
when its first components match the initial $Z$ word.
Thus $r_{\rm inv}(N,Q)=|\mathcal L_N(Z)|\le8N$ and $\hinv(Q)=0$.

For $\eta\in\Theta$, use the Borel selector
\[
 s_\eta(z,\theta,y,r)=(z_0,\one_{\{\theta=\eta\}}y_0).
\]
Its four fibers are nonempty, Borel, and assigned different
controls. Its actual name set is $Z\times Y_\eta$:
the slice $\theta=\eta$ realizes the full product, and the
other slices add only the zero second sequence. Hence
\[
 \Patt_{s_\eta}(F)=\Patt_Z(F)\Patt_{Y_\eta}(F).
\]
The binary-addition argument in Proposition~\ref{prop:thue-morse}
also gives all patterns on any finite subset of
$\{4^j:j\ge0\}$: extend the prescribed digits to the largest
exponent and apply the same parity correction. Thus both factors
realize $2^n$ patterns on any $n$ positions in $B_\eta$.
It follows that $p_{s_\eta}^*(n)=4^n$ and $h^*(s_\eta)=\log4$.
Their ordinary word entropy is zero by the word bounds for
$Z$ and $Y_\eta$, and their name languages vary Hausdorff
continuously with $\eta$.

Fix $A$. By Lemma~\ref{lem:sparse-bundle}, the normalized
logarithm of the second pattern factor tends to zero almost
surely. Therefore
\[
 h_A(s_\eta)=h_{\rm seq,coord}^A(Z)\le\log2<\log4
 \quad\text{for }\nu\text{-almost every }\eta.
\]
This disproves pointwise common attainment. In the designated
family one also has
$\max_A\inf_\eta h_A(s_\eta)=\log2<\min_\eta h^*(s_\eta)=\log4$.
For the full Borel class the lower-rate selector $(z_0,0)$
is available, and the exact equality $L_1=M_1$ remains valid.
\end{proof}

In this finite-control example the compactness assertion concerns
the name languages. It does not assert compactness of the Borel
selector class. In particular, the strict gap for the designated
family is not a counterexample to the full-class minimax identity.


\subsection{Finite switching and a change of sampling alignment}

\begin{lemma}\label{lem:finite-switch}
Let $s_1,\ldots,s_m$ be legal feedbacks at the same fixed period
$\tau$, let $\lambda:Q\to\{1,\ldots,m\}$ be Borel, and set
$s(x)=s_{\lambda(x)}(x)$. Suppose there is one integer $K$
such that, for every $x\in Q$, the index sequence
$\lambda(T_s^kx)$ changes at most $K$ times along the entire
infinite trajectory. Then
\[
 h^*(s)\le\max_j h^*(s_j).
\]
The hypothesis concerns the feedback index, not changes of
the control labels produced by a single feedback.
\end{lemma}
\begin{proof}
Put $H=\max_jh^*(s_j)$. For every $\eta>0$ there is
$C_\eta\ge0$ with
$\log p_j^*(k)\le k\tau(H+\eta)+C_\eta$ for all $j,k$;
take $p_j^*(0)=1$. This follows from the finitely many
maximal-pattern limits and a bound for their finitely many
remaining small lengths.

Fix $F=\{a_1<\cdots<a_n\}$. An actual trajectory has at
most $K+1$ constant-index stages. The stages meeting $F$
divide its ordered indices into $r\le\min(K+1,n)$ consecutive
groups. For a group $a_b,\ldots,a_c$ assigned index $j$,
the trajectory starting at $T_s^{a_b}x$ follows $s_j$ through
$a_c$. Its group pattern therefore has at most
$p_j^*(c-b+1)$ possibilities, using the translated window
$\{a_i-a_b:b\le i\le c\}$. This window depends on the
group, not on the unobserved time at which its stage began.
Unobserved stages may be omitted, while stages on either side
are kept as separate groups when needed. Taking a Cartesian
product gives an upper bound without presuming independent
states or feasible concatenation.

There are at most $\binom{n-1}{r-1}$ groupings and $m^r$
index choices. Hence, uniformly in $F$,
\[
 \Patt_s(F)\le e^{n\tau(H+\eta)}
 \sum_{r=1}^{\min(K+1,n)}
       \binom{n-1}{r-1}m^re^{rC_\eta}.
\]
The sum is bounded by a polynomial of degree at most $K$ in
$n$. Taking the maximum over windows, normalized logarithms,
and then $\eta\downarrow0$ proves the claim. Safety of $s$
is pointwise inherited from the selected legal blocks, so the
argument applies to its full infinite trajectories.
\end{proof}

When $K=0$, each whole name comes from one of the original
feedbacks, and the analogous inequality holds for every fixed
$h_A$. One switch suffices to invalidate that conclusion.

\begin{example}[One waiting step changes a fixed sampled rate]
\label{ex:one-wait}
There is a compact system with three distinct continuous controls,
two clopen safe feedbacks $s_{\mathrm{wait}},s_{\mathrm{read}}$, and a
clopen safe mixture $s_{\mathrm{mix}}$ with at most one feedback-index
switch on every trajectory, for which a single sampling sequence
$A$ satisfies
\[
 h_A(s_{\mathrm{wait}})=h_A(s_{\mathrm{read}})=0,
 \qquad h_A(s_{\mathrm{mix}})=\log2.
\]
Their maximal-pattern rates are respectively $0,\log2,\log2$.
\end{example}
\begin{proof}
Let $D=\{4^j:j\ge1\}$ and
\[
 Y=\bigcup_{t\ge0}
 \{y\in\{0,1\}^{\Nzero}:\supp y\subset(D-t)\cap\Nzero\}.
\]
This union is closed. A positive difference between two powers
of four determines those powers uniquely, by its $4$-adic
valuation. Thus two ones in a convergent sequence eventually
fix the same powers and the same nonnegative translate $t$;
every other one of the limit uses that translate. Zero and
single-one sequences also belong to the union. This proves
closedness, while shifting replaces $t$ by $t+1$, so
$\sigma Y\subset Y$. Surjectivity is not assumed.

At any finite set of positions in $D$ all binary patterns occur,
so $p_Y^*(n)=2^n$. On the other hand, take
$A=(0,4^1+1,4^2+1,\ldots)$ and its first $n$ positions $F_n$.
A pattern there has at most one one. A difference from zero
to $4^j+1$ is not divisible by four; a difference between two
positive sampled positions uniquely forces the translate
$t=-1$, which is excluded. Zero and all single-one patterns
occur, whence $\Patt_Y(F_n)=n+1$.

Use $Q=Y\times\{0,1\}$, $X=Q\sqcup\{\dagger\}$, and
controls $0,1,2$. At $(y,r)$ the matching control $u=y_0$
sends the state to $(\sigma y,0)$, the other binary control
sends it to $\dagger$, and control $2$ sends it to $(y,0)$.
The cemetery point is fixed. These three maps are continuous
and distinct. The flag $r$ is a coordinate of this system's
state, rather than feedback memory.

Set $s_{\mathrm{wait}}=2$, $s_{\mathrm{read}}=y_0$, and take
$s_{\mathrm{mix}}=2$ when $r=1$, $s_{\mathrm{mix}}=y_0$ when $r=0$.
They are clopen and globally safe. The mixture waits once
from flag one and then reads forever; from flag zero it reads
immediately. Their actual infinite name sets are exactly
$\{2^\infty\}$, $Y$, and $Y\cup2Y$. Since $0\in F_n$,
the latter two pieces are distinguished by their first sampled
label; the waiting piece reads all patterns on the first
$n-1$ powers of four. Consequently
\[
 \Patt_{\mathrm{wait}}(F_n)=1,\quad
 \Patt_{\mathrm{read}}(F_n)=n+1,\quad
 \Patt_{\mathrm{mix}}(F_n)=2^{n-1}+n+1.
\]
These equalities prove the sampled rates. On any $n$-position
window the mixture has at most $2^{n+1}$ patterns; its displayed
sampled lower bound gives maximal rate $\log2$, consistent
with Lemma~\ref{lem:finite-switch}.
\end{proof}

The constant waiting feedback makes the full-class values
$L_1=M_1=0$ in this example. It refutes a fixed-sampling
stability claim for finite switching, rather than the full-class
minimax identity. The change comes from restarting the reading
coordinates one step later on a one-sided language.


\section{Capacity peaks and probability-entropy bounds}
\label{app:capacity-boundaries}

The inverse-capacity theorem concerns support counts. Two tempting
substitutes are a strict growth discount for the actual feedback
transfer operator, and probability-entropy bounds for its names.
The following examples separate these objects. They give boundaries
of the estimates, not a strict inequality between $L_1$ and $M_1$.

\subsection{Invariant name measures need not certify the counting rate}

For $1<\beta<2$, put $Q=[0,M]$, $M=(\beta-1)^{-1}$, and use
$g_u(x)=\beta x-u$, $u=0,1$. The maps can be continuously clipped
on $[-1,\beta M]$, without changing their values on $Q$.
The safety domains are $[0,M/\beta]$ and $[1/\beta,M]$.
With normalized Lebesgue measure the scalar inverse capacities
are $1/\beta<1$, and the minimum safety-cover size is two.
Theorem~\ref{thm:capacity-main} gives $M_1=\log2$.
Every actual length-$n$ fiber lies in an interval of length
$M\beta^{-n}$, so consecutive observations give
$\log\beta\le L_1\le\log2$.

\begin{proposition}\label{prop:shannon-boundary}
For every integer $r\ge2$, take $\beta=2^{1/r}$ in this system.
The greedy feedback
\[
 s_g(x)=\begin{cases}0,&x<1/\beta,\\1,&x\ge1/\beta\end{cases}
\]
is Borel and globally safe. There is an arithmetic sampling
sequence $A_q=(0,q,2q,\ldots)$ with $h_{A_q}(s_g)=\log2$.
Nevertheless, for every sampling sequence $A$ and every
shift-invariant probability $\mu$ on its completed name system,
\[
 h_{A,\mu}(s_g;\mathcal P)\le
 H_2\left(\frac1{r+1}\right)<\log2,
\]
where $\mathcal P$ observes the single control coordinate,
$H_2(p)=-p\log p-(1-p)\log(1-p)$, and $h_{A,\mu}$ uses Shannon
entropy divided by the number of samples.
\end{proposition}
\begin{proof}
The interval $[0,1)$ is invariant under the greedy rule.
After a digit one there, the next state is below $\beta-1$.
Since $\beta^r=2$ and $2(\beta-1)<1$, the next $r$ digits
are zero. Let $Y_r$ be the binary subshift whose one-positions
have mutual distance at least $r+1$. Points below $M$ enter
$[0,1)$ after a finite initial run of ones: as long as the
state is at least one, the iterate is
$M-\beta^k(M-x)$. The endpoint $M$ is fixed. Thus the
completed name system is contained in the closed shift-invariant set
\[
 R_r=\{1^\infty\}\cup\bigcup_{k\ge0}1^kY_r.
\]
Closedness follows by separating bounded prefix lengths from
prefix lengths tending to infinity.

An invariant probability on $R_r$ has the form
$\mu=\lambda\delta_{1^\infty}+(1-\lambda)\nu$ with $\nu$
invariant and supported on $Y_r$: the increasing sets
$\sigma^{-n}Y_r$ all have measure $\mu(Y_r)$ and their union
is $R_r\setminus\{1^\infty\}$.
For $\nu$, invariance and $\sum_{j=0}^r\omega_j\le1$ give
$p=\nu(\omega_0=1)\le1/(r+1)$. For every $n$-element window $F$,
subadditivity and the component indicator give
\[
 H_\mu(\omega|_F)
 \le H_2(\lambda)+(1-\lambda)nH_2\left(\frac1{r+1}\right).
\]
Division by $n$ proves the claimed bound, including
nonergodic mixtures.

Choose $q$ so large that
$\beta^{q-1}(\beta-1)>1$, and put
$R=\beta^{q-1}/(\beta^q-1)<1$.
For any binary sequence $b$, let $d_{qj}=b_j$ and set the
other digits to zero. Define
$x_k=\sum_{i\ge k}d_i\beta^{-(i-k+1)}$.
Then $0\le x_k\le R<1$.
If $d_k=1$, then $x_k\ge1/\beta$; if $d_k=0$, then
$x_k\le R/\beta<1/\beta$.
The identity $x_{k+1}=\beta x_k-d_k$ therefore furnishes an
actual infinite greedy trajectory. All binary patterns occur
at the sampled positions, proving $h_{A_q}(s_g)=\log2$.
\end{proof}

The inequality is restricted to the fixed control-coordinate
partition. It makes no assertion about the supremum over all
measurable partitions, and it supplies no counting-entropy upper
bound for any other feedback.

\subsection{One feedback can saturate the transfer norm at every depth}

For this construction take $Q=[0,1]$ and
\[
 g_0(x)=\beta x,\qquad g_1(x)=\beta x-(\beta-1),\qquad
 \psi_0(y)=cy,\quad\psi_1(y)=cy+1-c,\quad c=\beta^{-1}.
\]
Here $K_0=[0,c]$, $K_1=[1-c,1]$, and again $M_1=\log2$. Continuous clipping on $[-2,2]$ gives a compact ambient realization without changing the affine formulas on $Q$.
For a legal selector $s$, with $S_u=s^{-1}(u)$, its actual
transfer operator on $L^\infty([0,1])$ is
\[
 (\mathcal P_sf)(y)=
 \sum_{u=0}^1c\,\one_{S_u}(\psi_u(y))f(\psi_u(y)).
\]
Affine change of variables gives
$\int h(T_sx)f(x)\,dx=\int h(y)\mathcal P_sf(y)\,dy$.
In particular this is the operator of the actual closed loop,
with all selector restrictions retained. Nondegenerate affine inverse maps preserve null sets, so it is well defined on $L^\infty$, and
$\opnorm{\mathcal P_s^m}\le(2c)^m$.

\begin{proposition}\label{prop:full-depth-saturation}
For every rational $1<\beta<2$, there is one globally safe
Borel selector $s_\beta$ such that
\[
 \opnorm{\mathcal P_{s_\beta}^m}=(2/\beta)^m
 \quad(m\ge1),\qquad
 \Patt_{s_\beta}(F)=2^{|F|}
 \quad(F\subset\Nzero\text{ finite}).
\]
At each depth the norm equality holds on a set of positive
Lebesgue measure. Thus no strict exponential discount of the
base $2/\beta$ holds even for this one feedback.
\end{proposition}
\begin{proof}
For a finite binary word $w$, let
$\psi_w=\psi_{w_0}\circ\cdots\circ\psi_{w_{|w|-1}}$,
including the empty word with $\psi_\varnothing=\mathrm{id}$.
If $\beta=p/q$ in lowest terms, $q\ge2$, then different words
of the same length give different affine maps. Otherwise their
translation difference would produce a nonzero polynomial
$P(\beta)=0$ with coefficients in $\{-1,0,1\}$ and leading
coefficient $\pm1$. Multiplication by $q^{\deg P}$ and reduction
modulo $q$ contradict coprimality. Maps of different lengths
have different rational slopes, so equality at a point can occur
only at a rational point. Consequently, at any irrational $y$,
all finite-word nodes $\psi_w(y)$ are distinct.

We use a compact nowhere dense set of positive measure. Starting
with $[0,1]$, remove at stage $k$ a central open interval of
length $4^{-k}$ from each remaining interval. The component
length after stage $k$ is
$2^{-k-1}+2^{-2k-1}$, by the splitting recurrence.
Thus the construction is defined at every stage, the removed
lengths sum to $1/2$, and the limiting compact set is nowhere
dense and has measure $1/2$. Its affine copy $F_J$ in a
nondegenerate closed interval $J$ has measure $|J|/2$.

Inductively choose intervals $J_j\subset(\beta-1,1)$ and set
$F_j=F_{J_j}$. Let
\[
 D_{j-1}=\bigcup_{i<j}\ \bigcup_{|w|\le i}\psi_w(F_i).
\]
This finite union is compact and nowhere dense.
The set
\[
 (\beta-1,1)\setminus
       \bigcup_{|w|\le j}\psi_w^{-1}(D_{j-1})
\]
is open dense. Choose an irrational $y_j$ there. The finitely
many nodes with $|w|\le j$ are distinct and avoid $D_{j-1}$.
Their pairwise distances, their distances to $D_{j-1}$ when
nonempty, and the distances of $y_j$ to the window endpoints
have a positive minimum $d_j$. Taking
$J_j=[y_j-d_j/4,y_j+d_j/4]$ makes all the intervals
$\psi_w(J_j)$, $|w|\le j$, disjoint and disjoint from
$D_{j-1}$, since their Lipschitz constants are at most one.
This separates nodes of different lengths as well as nodes
from different stages.

Define the disjoint $F_\sigma$ sets
\[
 C_u=\bigcup_{j\ge1}\ 
      \bigcup_{\substack{1\le|w|\le j\\w_0=u}}\psi_w(F_j),
 \qquad u=0,1.
\]
They lie in $\intr K_u$. Set
\[
 s_\beta(x)=
 \begin{cases}
  1,&x\in C_1\cup(Q\setminus K_0),\\
  0,&\text{otherwise}.
 \end{cases}
\]
This is a single Borel state rule. It chooses zero on $C_0$,
one on $C_1$, and is safe everywhere because $K_0\cup K_1=Q$.
It treats every endpoint by the same formula and gives infinite
safe trajectories for every initial state.

For $w$ of length $m$ and $y\in F_m$, the states
$x_i=\psi_{w_i\cdots w_{m-1}}(y)$, $0\le i<m$, with $x_m=y$,
satisfy $s_\beta(x_i)=w_i$ and $g_{w_i}(x_i)=x_{i+1}$.
Thus every word is a real prefix of this same feedback.
Expansion of the actual transfer operator gives
\[
 (\mathcal P_{s_\beta}^m1)(y)
 =\sum_{|w|=m}c^m\one_{E_w(s_\beta)}(\psi_w(y))
 =(2c)^m\qquad(y\in F_m).
\]
Since $|F_m|>0$, this proves the essential-norm lower bound;
the general upper bound proves equality. The sets $F_m$ need
not have a common positive measure bound.
Every pattern on a finite $F$ can be completed to a word of
length $1+\max F$, which has an actual witness by the
construction. The binary alphabet gives the matching count.
\end{proof}

The selector has counting rate $\log2$ along every sampling
sequence. It therefore refutes a transfer-norm estimate, not
the full-class equality $L_1=M_1$. Changing the feedback can
destroy the constructed inverse trees.

\subsection{A conditional collision bound can miss full support}

For an affine feedback on $Q=[0,1]$, fix sampled times
$a_1<\cdots<a_n$ and a terminal time $T=a_n+1$.
Let $E_\omega$ be the actual initial-state fiber of a sampled
name $\omega$. Its joint density with the terminal state is
defined by
\[
 \int_{E_\omega}h(T_s^Tx)\,dx
 =\int_Q h(y)f_\omega(y)\,dy.
\]
Piecewise affine change of variables supplies these densities.
Put $\rho=\sum_\omega f_\omega$, so $\int\rho=1$, and set
\[
 R_n=\int_{\rho>0}
       \frac{\sum_\omega f_\omega(y)^2}{\rho(y)}\,dy.
\]
If $N_n$ counts all nonempty actual sampled fibers, including
zero-measure fibers, Cauchy's inequality gives
$1/N_n\le R_n\le1$. Hence $-\log R_n$ is a lower bound for
$\log N_n$, not an upper bound.

\begin{example}\label{ex:collision-boundary}
Use three controls on $Q=[0,1]$,
\[
 g_a(x)=\tfrac43x,\qquad g_b(x)=4x-3,\qquad g_d(x)=16x-4,
\]
with safety domains $[0,3/4]$, $[3/4,1]$, and $[1/4,5/16]$.
Continuous clipping on $[-4,12]$ preserves these formulas on
$Q$. The minimum safety-cover size is two and the largest
single inverse capacity is $3/4<1$, so $M_1=\log2$.
The selector choosing $a$ on $[0,3/4)$ and $b$ on $[3/4,1]$
is globally safe and Borel; the third control remains available
in the full feedback class.

For this selector, every $a,b$ word has a real inverse path
ending at each $y\in(0,1)$, through the inverse maps
$\psi_a(y)=3y/4$ and $\psi_b(y)=3/4+y/4$.
All intermediate inverse states lie in the appropriate open
selector domains. The endpoints are handled by the stated rule.
Summing the unsampled letters, with $p_a=3/4$, $p_b=1/4$,
gives for every sampling sequence
\[
 f_\omega(y)=\prod_{j=1}^np_{\omega_j}\quad\text{a.e.},\qquad
 \rho(y)=1,\qquad N_n=2^n,\qquad R_n=(5/8)^n.
\]
Thus $h_A(s)=\log2$ for every $A$, whereas
$\lim_n-n^{-1}\log R_n=\log(8/5)<\log2$.
The terminal density is already uniform; the loss comes from
the unequal probabilities of the labels. The calculation is
the familiar nonuniform full-branch coding phenomenon, and
does not establish a full-class counting-entropy gap.
\end{example}

For a finite atom refinement of the same physical feedback,
coarse densities are sums of fine densities. Therefore the
fine support count is no smaller, while its collision quantity
is no larger. The equalities for the unrefined selector are
not asserted for arbitrary refinements.


\section*{Use of AI tools}
OpenAI ChatGPT/Codex tools assisted the argument exploration,
candidate proofs, comparison of primary literature, manuscript
integration, and automated source, compilation, and PDF checks.
These activities included mathematical work and were not limited
to language editing. Model versions for the complete earlier
research history cannot be reliably reconstructed from the record.

\begin{thebibliography}{99}

\bibitem{Debauche2022}
Debauche, V., Della Rossa, M., \& Jungers, R.~M. (2022).
Comparison of path-complete Lyapunov functions via template-dependent lifts.
\emph{Nonlinear Analysis: Hybrid Systems, 46}, 101237.
\doi{10.1016/j.nahs.2022.101237}.
Author version: \href{https://arxiv.org/abs/2110.13474v1}{arXiv:2110.13474v1}.

\bibitem{Goodman1974}
Goodman, T.~N.~T. (1974). Topological sequence entropy.
\emph{Proceedings of the London Mathematical Society, 29}(2), 331--350.
\doi{10.1112/plms/s3-29.2.331}.

\bibitem{HuangMaassYe2004}
Huang, W., Maass, A., \& Ye, X. (2004).
Sequence entropy pairs and complexity pairs for a measure.
\emph{Annales de l'Institut Fourier, 54}(4), 1005--1028.
\doi{10.5802/aif.2041}.

\bibitem{HuangYe2009}
Huang, W., \& Ye, X. (2009).
Combinatorial lemmas and applications to dynamics.
\emph{Advances in Mathematics, 220}(6), 1689--1716.
\doi{10.1016/j.aim.2008.11.009}.

\bibitem{Hulse1982}
Hulse, P. (1982). Sequence entropy and subsequence generators.
\emph{Journal of the London Mathematical Society, 26}(3), 441--450.
\doi{10.1112/jlms/s2-26.3.441}.

\bibitem{KamaeZamboniSystems2002}
Kamae, T., \& Zamboni, L. (2002).
Maximal pattern complexity for discrete systems.
\emph{Ergodic Theory and Dynamical Systems, 22}(4), 1201--1214.
\doi{10.1017/S0143385702000585}.

\bibitem{Kawan2013}
Kawan, C. (2013).
\emph{Invariance entropy for deterministic control systems:
An introduction}. Lecture Notes in Mathematics, Vol.~2089.
Springer. \doi{10.1007/978-3-319-01288-9}.

\bibitem{KL2007}
Kerr, D., \& Li, H. (2007).
Independence in topological and C$^*$-dynamics.
\emph{Mathematische Annalen, 338}, 869--926.
\doi{10.1007/s00208-007-0097-z}.

\bibitem{KL2009}
Kerr, D., \& Li, H. (2009).
Combinatorial independence in measurable dynamics.
\emph{Journal of Functional Analysis, 256}(5), 1341--1386.
\doi{10.1016/j.jfa.2008.12.014}.

\bibitem{Romagnoli2007}
Romagnoli, P.-P. (2007).
A remark on the topological entropies of covers and partitions.
\emph{Studia Mathematica, 182}(3), 273--281.
\doi{10.4064/sm182-3-6}.

\bibitem{Shi2021}
Shi, L., Zheng, W.~X., Shao, J., \& Cheng, Y. (2021).
Sub-Super-stochastic matrix with applications to bipartite tracking
control over signed networks.
\emph{SIAM Journal on Control and Optimization, 59}(6), 4563--4589.
\doi{10.1137/19M1239982}.
Author version: \href{https://arxiv.org/abs/2004.01867v2}{arXiv:2004.01867v2}.

\bibitem{TomarKawanZamani2022}
Tomar, M.~S., Kawan, C., \& Zamani, M. (2022).
Numerical over-approximation of invariance entropy via finite abstractions.
\emph{Systems \& Control Letters, 170}, 105395.
\doi{10.1016/j.sysconle.2022.105395}.

\end{thebibliography}

\end{document}
